% ============================================================
% Geometric Dosimetry — monograph_v2.tex
% Author : R. Wydaeghe
% Affiliation: Ghent University
% ============================================================

\documentclass[11pt,a4paper]{article}

% ---------------------------------------------------------------
% GEOMETRY — margins are intentionally wide; do not change
% ---------------------------------------------------------------
\usepackage[margin=2.4cm]{geometry}

% ---------------------------------------------------------------
% ENCODING AND FONTS
% ---------------------------------------------------------------
\usepackage[utf8]{inputenc}
\usepackage[T1]{fontenc}
\usepackage{lmodern}           % CM-compatible outlined fonts
% Alternative: Times-like math (closer to AMS / IEEE style)
%   \usepackage{newtxtext,newtxmath}  % requires newtx package
\usepackage{microtype}         % character-level justification; always last font pkg
\linespread{1.04}              % slight openness; 1.0 is pure single-spacing

% ---------------------------------------------------------------
% MATHEMATICS — core
% ---------------------------------------------------------------
\usepackage{amsmath,amssymb,amsthm,mathtools}
\usepackage{bm}                % bold math (vectors, matrices)
\usepackage{physics}           % \abs, \norm, \dd, \grad, …
\usepackage{mathrsfs}          % \mathscr{} script letters (e.g. for operators)
\usepackage{bbm}               % \mathbbm{1} blackboard-bold indicator
\usepackage{cancel}            % \cancel, \bcancel for struck-through terms
\usepackage{xfrac}             % \sfrac{}{} for inline slanted fractions

% Allow long equation arrays to break across pages
\allowdisplaybreaks[2]

% ---------------------------------------------------------------
% MATHEMATICS — highlighted / boxed equations
% ---------------------------------------------------------------
\usepackage{empheq}            % \begin{empheq}[box=\fbox]{align}
% Define a shaded box option for main results:
\usepackage[most]{tcolorbox}
\tcbset{highlight math style={%
    enhanced, colframe=black!70, colback=gray!6,
    arc=2pt, boxrule=0.5pt, left=3pt, right=3pt}}

% ---------------------------------------------------------------
% SI UNITS
% ---------------------------------------------------------------
\usepackage{siunitx}           % \SI{28}{\giga\hertz}, \si{\watt\per\kilo\gram}
\sisetup{
  detect-all,                  % inherit surrounding font
  range-phrase = {--},
  range-units  = single,
  per-mode     = symbol,
  inter-unit-product = {\,}
}

% ---------------------------------------------------------------
% FIGURES AND TABLES
% ---------------------------------------------------------------
\usepackage{graphicx}
\graphicspath{{figures/}}
\usepackage{booktabs}          % \toprule, \midrule, \bottomrule
\usepackage{array}             % extended column specifiers
\usepackage{tabularx}          % full-width tables
\usepackage{multirow}
\usepackage{makecell}          % multi-line table cells
\usepackage{caption}
\captionsetup{font=small, labelfont=bf, labelsep=period,
              skip=4pt}        % compact captions, bold label, period separator
\usepackage{subcaption}
\usepackage{threeparttable}    % table notes for lit-review tables
\usepackage{float}
\usepackage[section]{placeins} % keep floats within their section

% Relax float placement rules
\renewcommand{\topfraction}{0.9}
\renewcommand{\bottomfraction}{0.9}
\renewcommand{\textfraction}{0.1}
\renewcommand{\floatpagefraction}{0.7}
\setcounter{topnumber}{4}
\setcounter{bottomnumber}{4}
\setcounter{totalnumber}{8}

% ---------------------------------------------------------------
% SECTION HEADINGS (titlesec)
% ---------------------------------------------------------------
\usepackage{titlesec}
\titleformat{\section}{\large\bfseries}{\thesection}{1em}{}
\titleformat{\subsection}{\normalsize\bfseries}{\thesubsection}{1em}{}
\titleformat{\subsubsection}{\normalsize\bfseries}{\thesubsubsection}{1em}{}
\titlespacing*{\section}{0pt}{1.6ex plus .4ex minus .2ex}{0.8ex plus .2ex}
\titlespacing*{\subsection}{0pt}{1.2ex plus .3ex minus .2ex}{0.5ex plus .1ex}

% ---------------------------------------------------------------
% MISCELLANEOUS
% ---------------------------------------------------------------
\usepackage{enumitem}          % fine-grained list control
\usepackage{xcolor}
\usepackage{soul}              % \hl{} for inline erratum todos
\usepackage{xspace}            % smart spacing after macro names
\usepackage{todonotes}[disable]% \todo{} comments — set [disable] for final version
\usepackage{fancyhdr}          % running headers

% Running header: empty header, page number centred in footer
\pagestyle{fancy}
\fancyhf{}
\cfoot{\thepage}
\renewcommand{\headrulewidth}{0pt}   % no rule when header is empty

% ---------------------------------------------------------------
% BIBLIOGRAPHY
% ---------------------------------------------------------------
\usepackage[numbers,sort&compress]{natbib}

% ---------------------------------------------------------------
% HYPERLINKS — load near last, before cleveref
% ---------------------------------------------------------------
\usepackage[
  colorlinks = true,
  linkcolor  = {blue!55!black},
  citecolor  = {green!45!black},
  urlcolor   = {blue!65!black},
  pdfstartview = FitH,
  bookmarksnumbered = true,
  bookmarksopen     = true
]{hyperref}
\usepackage{bookmark}          % fixes bookmark levels in PDF outline
\usepackage{doi}               % live DOI hyperlinks (must load AFTER hyperref)

% cleveref must come AFTER hyperref
\usepackage[nameinlink,capitalize]{cleveref}

% ---------------------------------------------------------------
% THEOREM-LIKE ENVIRONMENTS
% (all share the same counter as theorem, numbered per section)
% ---------------------------------------------------------------
\usepackage{aliascnt}

\theoremstyle{definition}      % upright body text for all env.

\newtheorem{theorem}{Theorem}[section]

\newaliascnt{proposition}{theorem}
\newtheorem{proposition}[proposition]{Proposition}
\aliascntresetthe{proposition}

\newaliascnt{lemma}{theorem}
\newtheorem{lemma}[lemma]{Lemma}
\aliascntresetthe{lemma}

\newaliascnt{corollary}{theorem}
\newtheorem{corollary}[corollary]{Corollary}
\aliascntresetthe{corollary}

\newaliascnt{definition}{theorem}
\newtheorem{definition}[definition]{Definition}
\aliascntresetthe{definition}

\newaliascnt{remark}{theorem}
\newtheorem{remark}[remark]{Remark}
\aliascntresetthe{remark}

% cleveref names for theorem-like environments
\crefname{proposition}{Proposition}{Propositions}
\Crefname{proposition}{Proposition}{Propositions}
\crefname{lemma}{Lemma}{Lemmas}
\Crefname{lemma}{Lemma}{Lemmas}
\crefname{corollary}{Corollary}{Corollaries}
\Crefname{corollary}{Corollary}{Corollaries}
\crefname{definition}{Definition}{Definitions}
\Crefname{definition}{Definition}{Definitions}
\crefname{remark}{Remark}{Remarks}
\Crefname{remark}{Remark}{Remarks}

% ---------------------------------------------------------------
% CUSTOM MATH COMMANDS
% ---------------------------------------------------------------
% --- vectors and fields ---
\newcommand{\khat}{\hat{\bm{k}}}
\newcommand{\nhat}{\hat{\bm{n}}}
\newcommand{\rr}{\mathbf{r}}
\newcommand{\EE}{\mathbf{E}}
\newcommand{\HH}{\mathbf{H}}
% --- dosimetric quantities ---
\newcommand{\Sinc}{S_{\mathrm{inc}}}
\newcommand{\Sab}{S_{\mathrm{ab}}}
\newcommand{\Teff}{T_{\mathrm{eff}}}
\newcommand{\Tavg}{T_{\mathrm{avg}}}
\newcommand{\Aab}{A_{\mathrm{ab}}}
\newcommand{\Aperp}{A_{\perp}}
\newcommand{\SARwb}{\mathrm{SAR}_{\mathrm{wb}}}
% --- material parameters ---
\newcommand{\ntilde}{\tilde{n}}
\newcommand{\epstilde}{\tilde{\varepsilon}_r}
% --- operators ---
\DeclareMathOperator{\ReLU}{ReLU}
% (\erf already provided by the physics package)
\DeclareMathOperator{\RE}{Re}
\DeclareMathOperator{\IM}{Im}
% --- other shorthands ---
\newcommand{\Tbar}{\bar{T}}
\newcommand{\diff}{\mathrm{d}}        % upright differential
\newcommand{\ident}{\mathbbm{1}}       % indicator / identity  (requires bbm)
\newcommand{\Reals}{\mathbb{R}}
\newcommand{\Complex}{\mathbb{C}}
\newcommand{\half}{\tfrac{1}{2}}

% ---------------------------------------------------------------
% PART HEADINGS (article class has no \part)
% ---------------------------------------------------------------
\newcommand{\partheading}[2]{%
  \clearpage
  \phantomsection
  \addcontentsline{toc}{part}{#1}%
  \begin{center}
    \vspace*{1.5cm}
    {\LARGE\bfseries #1}\\[0.6cm]
    {\large #2}\\[1.5cm]
  \end{center}
}

% ---------------------------------------------------------------
% KEYWORDS command — plain comma list, wraps naturally
% ---------------------------------------------------------------
\newcommand{\keywords}[1]{%
  \par\smallskip\noindent\textbf{Keywords:}\enspace\parbox[t]{\dimexpr\linewidth-6.5em}{#1}\par\smallskip}

% ---------------------------------------------------------------
% TITLE / AUTHOR — PRIMEarxiv style with \hrule flanking
% (redefine \maketitle to add the rules)
% ---------------------------------------------------------------
\title{\textbf{Geometric Dosimetry:\\[0.3em]
Closed-Form Absorption Laws \\
from 100\,MHz to 100\,GHz}}
\author{
  Robin Wydaeghe \\
  Ghent University \& imec \\
  Ghent, Belgium \\
  \normalfont\texttt{robin.wydaeghe@ugent.be}
}
\date{\today}

% PRIMEarxiv-style title block: rule -- title/author -- rule
\makeatletter
\renewcommand{\maketitle}{%
  \thispagestyle{plain}%        % no header on title page
  \begin{center}
    \rule{\linewidth}{1pt}\\[0.4em]
    {\LARGE\bfseries \@title}\\[0.5em]
    {\large \@author}\\[0.3em]
    {\small \@date}
    \\[0.4em]\rule{\linewidth}{1pt}
  \end{center}
  \vspace{0.5em}
}
\makeatother

% ============================================================
\begin{document}
% PDF metadata (set here to avoid xkeyval comma-parsing issues)
\hypersetup{
  pdftitle   = {Geometric Dosimetry: Closed-Form Absorption Laws from 100 MHz to 100 GHz},
  pdfauthor  = {Robin Wydaeghe},
  pdfsubject = {Electromagnetic Dosimetry, mmWave, MIMO Beamforming},
  pdfkeywords= {Absorbed power density, Fresnel, pseudo-Brewster, ambient occlusion,
                ReLU network, exposure operator, dosimetry, 5G, 6G}
}
\maketitle

% ============================================================
\begin{abstract}
Brute-force electromagnetic simulation of absorbed power density on
human bodies at millimetre-wave frequencies ($>6$\,GHz) is intractable:
the skin depth (${\sim}\,0.5$\,mm) is five orders of magnitude smaller
than the body (${\sim}\,1$\,m), requiring ${\sim}\,10^{12}$ mesh cells.
We show that for biological tissue the absorbed power density admits a
closed-form expression that separates material response from geometry.
A pseudo-Brewster compensation between TE and TM polarisations
renders the unpolarised Fresnel transmission nearly angle-independent,
collapsing the material factor to a single scalar~$T_0$.  The
remaining geometric factor is the positive-part cosine of the surface
normal, governed by projected area, Cauchy's formula, and an
exposure fraction identical to ambient occlusion in computer graphics.
For multiple sources, the result is a single-hidden-layer ReLU neural
network, making the dosimetry chain differentiable and compatible with
ray-tracing engines.
%
The framework extends in two directions.  An absorption Stokes
vector computes total power for arbitrary polarisation; the
worst-case polarisation correction on a human body is 16\,\% and
drops below 2.5\,\% with multipath diversity ($N\ge 20$ paths).
An exact flux-weighted transmission~$\bar{T}$ makes total-power and
compliance results exact down to 100\,MHz.  Validation against Mie
theory and anatomical phantoms confirms errors of 5\,\% locally and
10\,\% for total absorbed power, within the 20\,\% uncertainty of
tissue-property data.
%
Part~III extends the framework to coherent beamforming, where $M$
antenna elements radiate with controlled complex weights.  Two
well-bounded approximations collapse the coherent double sum to a
squared norm $S_{\mathrm{ab}}(\mathbf{r}) =
\|\tilde{\mathbf{G}}(\mathbf{r})\mathbf{x}\|^2$, from which an
\emph{exposure operator}~$\mathbf{Q}$ encodes the complete
body-absorption response to any precoder and enables closed-form
exposure-constrained beamforming.
\end{abstract}

\keywords{absorbed power density, Fresnel transmission, pseudo-Brewster compensation, ambient occlusion, exposure operator, mmWave, dosimetry, 5G/6G, compliance}

\tableofcontents
\newpage

% ============================================================
\section*{Preface}
% ============================================================

I set out to interpret some FDTD results I had for millimetre-wave absorption on the human body, a
problem that by brute force requires trillions of mesh cells.  To my surprise, the TE and
TM Fresnel transmissions through tissue compensate each other, causing the
material factor to collapse to a single constant.  This pseudo-Brewster
compensation is a consequence of the high refractive index of living tissue.

Once the material factor becomes a constant, dosimetry becomes geometry.
I did not expect the geometry to have the structure it does.  The
direction-averaged absorption satisfies Cauchy's projected-area formula,
first proved in 1841.  The self-shadowing correction is ambient
occlusion in computer graphics.  The multi-source formula is a ReLU
neural network whose weights come from a ray tracer.  At every turn, the
answer was a known result in another field.

This work is long because the framework touches many areas, and I tried to be comprehensive.  
The derivations start from Maxwell's equations.  The five irreducible limitations are stated in
\cref{sec:limitations-v2}.  Where the framework is approximate, the error
is bounded and smaller than the uncertainty in the tissue data.

I believe the elegance of the mathematics repays the effort of reading through these pages.

\newpage

% ============================================================
\section{Introduction}\label{sec:introduction}
% ============================================================

\subsection{Problem definition}

Millimeter-wave (mmWave) wireless systems operating above 6\,GHz are
central to 5G and future 6G networks.  Regulatory compliance requires
evaluation of the absorbed power density~$\Sab$ on the human
body~\cite{ICNIRP2020}, defined as the electromagnetic power deposited
per unit surface area in the superficial tissue layer.  At mmWave
frequencies the skin depth is $\delta\sim 0.3$ to $0.5$\,mm, so all
penetrating power is absorbed within less than a millimetre of the
surface.

Full-wave methods such as the finite-difference time-domain (FDTD)
technique require a spatial mesh fine enough to resolve this thin
absorption layer~\cite{Dimbylow2002,Hirata2021}.  Because the human
body spans roughly 1\,m, the resulting mesh contains of the order of
$10^{12}$~cells.  This is far beyond present computational resources.
An analytical or semi-analytical framework that replaces the volumetric
simulation with a surface computation is therefore needed.

\subsection{Summary of the approach}

We show that the absorbed power density at a surface point~$\rr$
illuminated by a plane wave of power density~$\Sinc$ and propagation
direction~$\khat$ is given exactly by what we call the
\emph{exact absorption law} (derived in \cref{sec:local-law}):
\begin{equation}\label{eq:exact-law}
  \boxed{\Sab(\rr)
    \;=\; \Sinc\;\cdot\; \Teff(\rr)\;\cdot\;
    \ReLU\!\bigl[\nhat(\rr)\cdot(-\khat)\bigr],}
\end{equation}
where $\Teff$ is the Fresnel power-absorption coefficient, which
depends on the local incidence angle and polarisation state, and
$\ReLU(x)=\max(0,x)$ gates illumination (front-facing versus
back-facing).  This equation holds under the locally flat-surface
assumption, for any polarisation and any frequency at which the body
is electromagnetically opaque.  The word ``exact'' distinguishes it
from the simplified form that follows; the only approximation is the
flat-interface model of the tissue boundary.

For biological tissue at millimetre-wave frequencies, two independent
simplifications collapse the exact absorption law to a formula that
depends on geometry alone.  First, a pseudo-Brewster
compensation between TE and TM polarisations renders the unpolarised
Fresnel transmission nearly independent of incidence angle, so that
the full angle-dependent transmission~$\Tavg(\theta)$ may be replaced
by the constant~$T_0$ (the normal-incidence value).  Second,
for circularly polarised or unpolarised illumination the
polarisation-dependent correction vanishes identically.  Under both
conditions the exact absorption law reduces to the
\emph{geometric absorption law}:
\begin{equation}\label{eq:geometric-law}
  \boxed{\Sab(\rr)
    \;\approx\; \Sinc\;\cdot\; T_0\;\cdot\;
    \ReLU\!\bigl[\nhat(\rr)\cdot(-\khat)\bigr].}
\end{equation}
The tissue physics enters only through the scalar~$T_0$; all spatial
variation is determined by the body's shape.

The formula connects to three fields.
First, in integral geometry, total absorbed power equals
$\Sinc\,T_0$ times the projected (shadow) area.
Direction-averaged absorption satisfies a generalisation of
Cauchy's formula~\cite{Cauchy1841,Hadwiger1957}, exact for non-convex
bodies when weighted by an exposure fraction~$\eta(\rr)$ that encodes
self-shadowing.  This gives a closed-form dosimetry formula for
reverberation chambers and other isotropic environments.
Second, in computer graphics, the exposure fraction is
identical to ambient occlusion~\cite{Zhukov1998}, one of the
most optimised computations in real-time rendering.
Finally, in machine learning and differentiable ray tracing,
for multiple incident waves, $\Sab$ is a single-hidden-layer
ReLU network~\cite{NairHinton2010} whose weights are the path
directions and powers from a ray
tracer~\cite{SionnaRT}.  The entire dosimetry chain is differentiable.
This supports gradient-based optimisation of antenna configurations
subject to exposure constraints.

The exact absorption law (\cref{sec:local-law}) holds for arbitrary
polarisation at any frequency where the body is electromagnetically
opaque.  The simplifications that follow
(\cref{sec:pseudo-brewster,sec:geometry,sec:compliance,sec:computation,sec:applications})
specialise to the millimetre-wave band, where the pseudo-Brewster
compensation collapses the material factor to a single constant.
Part~II develops the full generality:
\cref{sec:polarisation} introduces the polarisation-aware absorption
Stokes vector, \cref{sec:sub6} extends total-power results to
sub-6\,GHz, and \cref{sec:limitations-v2} enumerates the few
irreducible limitations.
Part~III extends the framework to coherent MIMO beamforming, where $M$
antenna elements radiate with controlled complex weights and fields add
as amplitudes rather than powers.  Two well-bounded approximations
collapse the coherent depth integral to a squared norm
$\Sab(\rr)=\|\tilde{\mathbf{G}}(\rr)\mathbf{x}\|^2$; integrating over
the body yields the \emph{exposure operator}~$\mathbf{Q}$, whose
eigendecomposition encodes the body's complete absorption response to any
precoder and enables a closed-form solution to exposure-constrained
beamforming (\cref{sec:ecbf}).

\subsection{Scope and accuracy}

The core results apply from 100\,MHz to 100\,GHz.  The geometric
form is optimal in the 10--60\,GHz band (Part~I), where
the Fresnel approximation error is 2 to 6\,\%.  Total-power and
compliance results extend to sub-6\,GHz via the exact averaged
transmission~$\Tbar$ (Part~II, \cref{sec:sub6}).  Comparison with
exact Mie theory for lossy spheres (a worst-case geometry) shows
errors of roughly 10\,\% for total absorbed power at body-relevant
sizes, dominated by diffraction into the geometric shadow.  The
uncertainty in tissue dielectric properties is itself of
order 20\,\%, so the framework error is comparable to the parametric
uncertainty of any simulation.

\subsection{Structure of this work}

\Cref{sec:literature} surveys the relationship between this work
and prior results in electromagnetic dosimetry, diffuse-field
dosimetry, optics, integral geometry, and computer graphics.
\Cref{sec:summary} summarises the main findings across all three
parts.
\Cref{sec:limitations-v2} enumerates the five irreducible limitations
of the complete framework.
\Cref{sec:conclusion} gives the conclusions.

Part~I establishes the millimetre-wave dosimetry framework.
\Cref{sec:local-law} derives the exact, polarisation-aware local
absorption law from Poynting's theorem and Fresnel's equations, then
introduces the TE/TM decomposition that separates polarisation from
geometry.  \Cref{sec:pseudo-brewster} establishes the pseudo-Brewster
compensation and quantifies its accuracy across tissues and
frequencies.  \Cref{sec:geometry} develops the geometric framework:
projected area, self-shadowing, Cauchy's formula, and the absorption
directivity~$D(\khat)$.  \Cref{sec:compliance} applies the results to
whole-body SAR compliance, deriving closed-form criteria that depend
only on body mass, surface area, and a single material constant.
\Cref{sec:computation} presents the computational framework, including
the neural-network interpretation and the matrix formulation.
\Cref{sec:corrections} discusses curvature corrections and diffraction
smoothing as systematic refinements.
\Cref{sec:validation} collects validation results against Mie theory
and anatomical phantoms.
\Cref{sec:applications} develops three applications: exposure-aware wireless
network design with differentiable ray tracing, real-time exposure monitoring,
and reverberation chamber dosimetry.

Part~II extends the framework beyond the mmWave, unpolarised
specialisations of Part~I.
\Cref{sec:polarisation} develops polarisation-aware dosimetry via
the absorption Stokes vector, derives the cylinder bound, and shows
that multipath diversity suppresses polarisation corrections to below
2.5\,\%.
\Cref{sec:sub6} extends total-power and compliance results down to
100\,MHz by replacing $T_0$ with the exact averaged
transmission~$\Tbar$, and identifies where the local surface map
breaks down.

Part~III extends the framework to coherent MIMO illumination.
\Cref{sec:coherent-setup} states the inherited assumptions.
\Cref{sec:field-channel} defines propagation paths, the field channel
matrix, the UE channel, and the signal and exposure branches.
\Cref{sec:fresnel-operator} derives the per-path transmitted field
under Approximation~1 (TM polarisation direction, $\lesssim 4\,\%$
error on cross-terms) and defines the Fresnel transmission operator.
\Cref{sec:coherent-absorption-law} derives the coherent absorption law
via the universal depth-coupling Approximation~2
($\leq 0.5\,\%$) and states the main theorem.
\Cref{sec:exposure-operator} introduces the exposure
operator~$\mathbf{Q}$ and its eigendecomposition into exposure modes.
\Cref{sec:MRT-body-section} analyses how MRT precoding distributes
absorbed power and introduces the exposure-signal alignment~$\rho$.
\Cref{sec:coherent-hotspot} analyses coherent hotspot formation and
the Fresnel filter's role in shaping the hotspot.
\Cref{sec:ecbf} solves the exposure-constrained beamforming problem.
\Cref{sec:mu-mimo} extends the framework to multi-user MIMO.
\Cref{app:coherent-cost} quantifies the computational cost of the
coherent pipeline.
Appendices~\ref{app:refractive-index} to~\ref{app:coherent-cost}
collect supplementary data, derivations (including the Fresnel equations
from Maxwell's boundary conditions in
\cref{app:fresnel-derivation}), proofs, approximation accuracy, and
the computational cost of the coherent pipeline.\relax

% ============================================================
\section{Literature review}\label{sec:literature}
% ============================================================

The framework draws on results from electromagnetic dosimetry, optics,
diffuse-field theory, integral geometry, and wireless communications.
The contribution of this work lies, in part, in connecting them, and in part, 
in expanding on existing results.
\Cref{tab:lit:comparison} summarises which elements of the framework
are addressed by prior work;
\cref{tab:lit:predictions} compares the framework's quantitative
predictions with independent observations in the literature.

% ------------------------------------------------------------------
\subsection{Absorbed power density and the local surface law}
% ------------------------------------------------------------------

Whole-body average SAR under plane-wave exposure has been evaluated by
FDTD simulation of anatomical voxel
models~\cite{Dimbylow2002,Christ2006,Hirata2007corr} for around three
decades; recent reviews survey computational methods above
6\,GHz~\cite{Hirata2021,Zhadobov2011,Colombi2018}.
A consistent empirical finding across these studies is that the
absorption cross section correlates linearly with body surface
area~\cite{Hirata2007corr,Dimbylow2002}.
Above 10\,GHz, a series of numerical studies confirmed that absorbed
power is proportional to the body's projected area times a nearly
constant transmission
coefficient~$T_{\mathrm{tr}}$~\cite{Kodera2024,Diao2024,Colombi2018}.
Kodera~\textit{et al.}~\cite{Kodera2024} showed that
$\SARwb = T_{\mathrm{tr}} \times \Aperp \times \Sinc / W$,
with~$T_{\mathrm{tr}}$ extracted from a one-dimensional
multilayer slab model, reproduces full three-dimensional results to
within 5\,\% from 10 to 100\,GHz.
Diao~\textit{et al.}~\cite{Diao2024} obtained
$T \approx 0.52$ at 28\,GHz from full anatomical FDTD; the Fresnel
prediction is $T_0 = 0.536$, a 3\,\% discrepancy.
In all these studies, $T_{\mathrm{tr}}$ is computed from tissue
dielectric properties by one-dimensional or full-wave simulation,
not from Fresnel theory, and no physical mechanism is identified for
its near-constancy across incidence angle and polarisation.

At the local level, several groups identified the transmitted Poynting
flux through the skin surface as the appropriate dosimetric quantity
at mmWave frequencies.
Funahashi~\textit{et al.}~\cite{Funahashi2018} proposed the
transmitted power density $\Sinc \times T_{\mathrm{tr}}$ as a
dosimetric metric and showed it correlates with temperature rise to
within 15\,\% from 30 to 300\,GHz.
Li~\textit{et al.}~\cite{Li2019} showed numerically that the
transmitted flux is nearly insensitive to incidence angle and
polarisation from 6\,GHz to 1\,THz on a planar skin model.
The unpolarised average $\Tavg(\theta)$ and its near-constancy were
not examined.
Samaras and Kuster~\cite{Samaras2019} computed TE and TM
transmittances separately versus incidence angle above 6\,GHz using
multilayer planar skin models.
For TM incidence, the worst-case transmittance over all
layer-thickness combinations reaches 0.92 at the pseudo-Brewster
angle (thick stratum corneum, 10\,GHz); this exceeds the value
assumed in the ICNIRP draft guidelines by up to 70\,\%.
The unpolarised average was not evaluated.
Diao~\textit{et al.}~\cite{Diao2021} confirmed that the heating
factor is approximately flat versus incidence angle at mmWave
frequencies.
Sasaki~\textit{et al.}~\cite{Sasaki2017} computed $T_0(f)$ by Monte
Carlo simulation of skin-layer thickness variability from 10\,GHz to
1\,THz, providing a reference dataset at normal incidence.
Christ~\textit{et al.}~\cite{Christ2020a} showed that a thick stratum
corneum acts as a quarter-wave matching layer above 15\,GHz.
On palms and soles this invalidates the single-$T_0$ approximation;
the present framework assumes a homogeneous skin half-space
throughout.

The physical mechanism behind the angular near-constancy of the
transmission is classical in optics.
The pseudo-Brewster angle for lossy dielectrics was characterised by
Potter~\cite{Potter1970} and named by
\"{O}hman~\cite{Ohman1977}.
Azzam~\cite{Azzam2015} showed that for lossless dielectric substrates
with refractive index $n \ge 2+\sqrt{3} \approx 3.73$ the unpolarised
reflectance deviates from its normal-incidence value by less than 1\,%
over $0$--$60^\circ$~\cite{BornWolf1999}.
Biological tissue at mmWave frequencies has $|\ntilde| \approx 3$--6
with a large imaginary part; it falls in the same refractive-index
magnitude range, yet this result has not been applied in the dosimetry
literature.

% ------------------------------------------------------------------
\subsection{Diffuse-field dosimetry and geometric structure}
% ------------------------------------------------------------------

Reverberation chamber (RC) measurements provide the direction-averaged
absorption cross section (ACS, $\sigma_a \equiv P_{\mathrm{abs}}/\Sinc$)
of the human body under isotropic illumination.
Andersen~\textit{et al.}~\cite{Andersen2007} introduced room
electromagnetics and the formula
$P_{\mathrm{abs}} = S_{\mathrm{total}} \,\langle\sigma_a\rangle$
for diffuse-field exposure.
Bamba~\textit{et al.}~\cite{Bamba2012} measured whole-body SAR of
volunteers in an RC at 1.8\,GHz; a companion study~\cite{Bamba2014}
derived a frequency-dependent absorption efficiency $\eta \approx 0.5$
from FDTD simulations of ellipsoidal body models over 1.45--5.8\,GHz.
No physical origin for $\eta$ was identified.
Flintoft~\textit{et al.}~\cite{Flintoft2014} measured the
direction-averaged ACS of 60 volunteers from 1 to 12\,GHz.
From 7\,GHz the normalised ACS
$\langle Q^a\rangle \equiv \langle\sigma_a\rangle/\langle G\rangle$
holds at approximately~0.41, where
$\langle G\rangle = \eta\,\mathrm{BSA}/4$ is the
Cauchy projected area~\cite{Cauchy1841,Vouk1948} and
$\mathrm{BSA}$ denotes the total body surface area.
Below 3\,GHz a dip appears, attributed qualitatively to standing
waves in the subcutaneous fat layer; no quantitative model was given.
Zhang~\cite{Zhang2017thesis,ZhangRobinson2020} measured the ACS of
48~subjects from 1 to 18\,GHz in an RC.
Defining the absorption coefficient
$\xi \equiv \langle\sigma_a\rangle/(0.25\times\mathrm{BSA})$,
he found $\xi \approx 0.45$--0.65 above 6\,GHz, gently rising with
frequency, and showed that the ACS--BSA correlation exceeds
$R = 0.9$ throughout 6--18\,GHz.
The minimum of $\xi(f)$ near 4--6\,GHz shifts to lower frequency with
increasing fat thickness, consistent with a Fabry--P\'{e}rot resonance in
the subcutaneous fat layer.
In this thesis, $\xi$ is systematically characterised but its
physical content ($\xi = \Tbar \times \Aab/A$, a product of the
flux-averaged Fresnel transmission and the non-convexity correction)
is not thoroughly derived from first principles.
Hallbj\"{o}rner~\textit{et al.}~\cite{Hallbjorner2005} showed that
the direction-averaged ACS of electrically large lossy objects in an
RC depends only on permittivity, not shape, in the skin-depth limit;
the mechanism (pseudo-Brewster compensation) was not identified.
In all these studies the Cauchy mean projected area and a transmission
coefficient appear as separate empirical ingredients;
they are not connected to Fresnel theory.

Cauchy's projected-area formula~\cite{Cauchy1841} is a classical
result of integral geometry (see
also~\cite{Hadwiger1957,SchneiderWeil2008}): for a convex body, the
mean projected area averaged over all directions is one quarter of the
surface area.
Ambient occlusion~\cite{Zhukov1998,Landis2002} is among the most
highly optimised computations in real-time computer
graphics~\cite{Pharr2016,AkenineMoller2018}; it quantifies the
fraction of the hemisphere from which a surface point is unoccluded.
Neither concept has been applied to electromagnetic dosimetry.

% ------------------------------------------------------------------
\subsection{Coherent MIMO and exposure-constrained beamforming}
% ------------------------------------------------------------------

When multiple antenna elements radiate with controlled complex
weights, the absorbed power is a quadratic form in the transmitted
signal: $\mathrm{SAR}(\mathbf{x}) = \mathbf{x}^H \mathbf{S}\,
\mathbf{x}$, where $\mathbf{S}$ is a Hermitian positive
semi-definite SAR matrix whose entries are
$[\mathbf{S}]_{ij} = (\sigma/\rho)\int_V \mathbf{E}_i^*(\rr)\cdot
\mathbf{E}_j(\rr)\,dV$.
This structure was introduced by
Hochwald~\textit{et al.}~\cite{Hochwald2014}, who validated it via
FDTD at 1.8\,GHz and demonstrated that SAR-aware phase coding
achieves 2.5\,dB gain over the Alamouti code and more than 3.5\,dB
gain over SAR-unaware beamforming.
In that work, $\mathbf{S}$ is calibrated numerically for a specific
device and body configuration; its entries have no closed-form
expression.
Ying~\textit{et al.}~\cite{Ying2015} established the full
capacity-optimisation framework: the precoder maximising capacity
subject to $\mathrm{SAR} \le \mathrm{SAR}_0$ takes the closed form
$\mathbf{x}^* = (\lambda\mathbf{S} + \nu\mathbf{I})^{-1}
\mathbf{h}^*$, with strong duality proved by induction.
Ying~\textit{et al.}~\cite{Ying2017} extended this to the multi-user
MIMO uplink, where each user~$u$ has a per-user SAR matrix
$\mathbf{S}^{(u)}$ and the exposure constraint is
$\mathrm{tr}(\mathbf{S}^{(u)}\mathbf{Q}^{(u)}) \le \mathrm{SAR}_0$,
with~$\mathbf{Q}^{(u)}$ the transmit covariance of device~$u$.
Ebadi-Shahrivar~\textit{et al.}~\cite{Ebadi2019} showed that
worst-case SAR over all beamforming vectors equals
$\lambda_{\max}(\mathbf{S})$, the leading eigenvalue, reducing
regulatory compliance testing to a single eigendecomposition.
Castellanos~\textit{et al.}~\cite{Castellanos2020}
derived per-element scalar Fresnel transmission coefficients at
28\,GHz and validated the per-path decomposition against FDTD.
Their calibration scalar $\beta \approx 0.70$ absorbs both the
Fresnel transmission and the depth-decay physics into a single number;
its relation to $T_0$ was not examined.
Zhou~\textit{et al.}~\cite{Zhou2026} derived pointwise exposure
matrices from dipole radiation geometry at 30\,GHz and coupled them
to a thermal dynamics model; their transmission
coefficient~$T_{\mathrm{tr}} = 0.8$ is taken as a parameter.
In all these works the exposure matrix is FDTD-calibrated; its
entries have no closed-form physical content.
All prior work assumes the uplink (handset to base station);
the downlink base-station geometry is not addressed.
It is in principle possible to couple the antenna channels via
multi-run FDTD
superposition~\cite{Shikhantsov2020,Wydaeghe2026npj}, but this
becomes intractable as frequency rises and yields no general closed
form.

% ------------------------------------------------------------------
\subsection{Numerical approaches and their limitations}
% ------------------------------------------------------------------

Full three-dimensional FDTD remains the regulatory standard for
dosimetric assessment~\cite{ICNIRP2020,Hirata2021}.
A single plane-wave simulation on an anatomical phantom at 28\,GHz
requires of order $10^{12}$~mesh cells; a systematic campaign
covering 15~frequencies, four phantoms, and two exposure scenarios
amounted to 550~individual simulations and several weeks of
GPU-accelerated computation~\cite{WydaegheEtAl2026PMB}.
Several groups have pursued shortcuts that reduce this cost.
Kodera~\textit{et al.}~\cite{Kodera2024} showed that the
decomposition $\SARwb = T_{\mathrm{tr}} \times \Aperp
\times \Sinc / W$, with $T_{\mathrm{tr}}$ from an analytical
one-dimensional multilayer slab model, reproduces full
three-dimensional results to within 5\,\% from 10 to 100\,GHz.
Kushiyama and Nagaoka~\cite{Kushiyama2025} precomputed spherical-wave
lookup tables (approximately 2.2\,TB per antenna geometry) allowing
field superposition without re-running FDTD per source configuration.
Zhang~\cite{Zhang2017thesis} proposed a surface-integral-equation
method (ABCD impedance matching combined with the method of moments)
whose memory scales with surface area rather than volume; this
approach was validated in two dimensions at 1--6\,GHz.
All these methods retain the numerical tissue-parameter input and do
not provide a closed-form expression for the transmission
coefficient or the exposure operator.


% ------------------------------------------------------------------
% Tables
% ------------------------------------------------------------------

\begin{table}[!htbp]
\centering
\begin{threeparttable}
\caption{Prior art on the elements of the geometric dosimetry
  framework.
  $\bullet$ = derived analytically or proved;
  $\circ$ = observed empirically without derivation;
  $\triangle$ = partial (see notes);
  --- = not addressed.}\label{tab:lit:comparison}
\small
\begin{tabularx}{\textwidth}{X
  c@{\hskip 4pt}c@{\hskip 4pt}c@{\hskip 4pt}c@{\hskip 4pt}c@{\hskip 4pt}c}
\toprule
& \makecell{$\Tavg$\\$\approx T_0$}
& \makecell{Local\\$\Sab(\rr)$}
& \makecell{Cauchy\\$\Aab/4$}
& \makecell{$\eta$\\$=$ AO}
& \makecell{Compl.\\formula}
& \makecell{Expos.\\op.\ $\mathbf{Q}$} \\
\midrule
Kodera \textit{et al.}\ 2024~\cite{Kodera2024}
  & $\circ$ & ---   & $\circ$ & ---   & ---   & --- \\
Funahashi \textit{et al.}\ 2018~\cite{Funahashi2018}
  & ---     & $\triangle$\tnote{1} & ---   & ---   & ---   & --- \\
Li \textit{et al.}\ 2019~\cite{Li2019}
  & $\circ$ & $\triangle$\tnote{2} & ---   & ---   & ---   & --- \\
Samaras \& Kuster 2019~\cite{Samaras2019}
  & $\triangle$\tnote{3} & ---   & ---   & ---   & ---   & --- \\
Diao \textit{et al.}\ 2021~\cite{Diao2021}
  & $\circ$ & ---   & ---   & ---   & ---   & --- \\
Diao \textit{et al.}\ 2024~\cite{Diao2024}
  & $\circ$ & ---   & ---   & ---   & ---   & --- \\
Sasaki \textit{et al.}\ 2017~\cite{Sasaki2017}
  & $\triangle$\tnote{4} & ---   & ---   & ---   & ---   & --- \\
Bamba \textit{et al.}\ 2012, 2014~\cite{Bamba2012,Bamba2014}
  & $\circ$ & ---   & ---   & ---   & ---   & --- \\
Flintoft \textit{et al.}\ 2014~\cite{Flintoft2014}
  & ---     & ---   & $\circ$ & ---   & ---   & --- \\
Zhang \& Robinson 2020~\cite{ZhangRobinson2020}
  & $\circ$ & ---   & $\circ$ & ---   & ---   & --- \\
Azzam 2015~\cite{Azzam2015}
  & $\bullet$\tnote{5} & ---   & ---   & ---   & ---   & --- \\
Gosselin \textit{et al.}\ 2011~\cite{Gosselin2011}
  & ---     & ---   & ---   & ---   & $\circ$ & --- \\
Hochwald \textit{et al.}\ 2014~\cite{Hochwald2014}
  & ---     & ---   & ---   & ---   & ---   & $\circ$ \\
Ying \textit{et al.}\ 2015, 2017~\cite{Ying2015,Ying2017}
  & ---     & ---   & ---   & ---   & $\circ$ & $\bullet$\tnote{6} \\
Castellanos \textit{et al.}\ 2020~\cite{Castellanos2020}
  & ---     & $\triangle$\tnote{7} & ---   & ---   & ---   & $\triangle$\tnote{7} \\
\midrule
This work
  & $\bullet$ & $\bullet$ & $\bullet$ & $\bullet$ & $\bullet$ & $\bullet$ \\
\bottomrule
\end{tabularx}
\begin{tablenotes}\footnotesize
  \item[1] Normal incidence only; no angular factor.
  \item[2] Flat planar skin model, normal incidence;
           no three-dimensional body geometry.
  \item[3] TE and TM transmittances computed separately;
           unpolarised average not evaluated.
  \item[4] Normal incidence only; no angular analysis.
  \item[5] Derived for high-index optical substrates
           ($|\ntilde|>2.5$); not applied to biological tissue
           or dosimetry.
  \item[6] SU-MIMO uplink: QCQP proved, precoder derived
           (Ying~2015).  MU-MIMO: additive per-user structure
           (Ying~2017).  SAR matrix empirically calibrated in
           both; no physical derivation of entries.
  \item[7] Per-element scalar Fresnel transmission at 28\,GHz;
           no body-surface integral, no depth coupling, no
           closed-form $\mathbf{Q}$.
\end{tablenotes}
\end{threeparttable}
\end{table}

\begin{table}[!htbp]
\centering
\caption{Framework predictions compared with independent observations
  in the dosimetry literature.  Detailed validation is collected in
  \cref{app:lit:validation}.}\label{tab:lit:predictions}
\small
\begin{tabularx}{\textwidth}{XXl}
\toprule
Published observation & Framework prediction & Agreement \\
\midrule
Bamba~\textit{et al.}~\cite{Bamba2012,Bamba2014}:
  absorption efficiency $\eta \approx 0.5$ at 1.5--6\,GHz
  in a diffuse field
& Framework identifies $\eta$ with the Fresnel
  coefficient $T_0$; pseudo-Brewster compensation
  explains the weak frequency dependence
& $T_0 = 0.47$--0.48; discrepancy 2--4\,\% \\[4pt]
%
Kodera~\textit{et al.}~\cite{Kodera2024}:
  ACS proportional to projected area above 10\,GHz
& Pseudo-Brewster compensation gives
  $\Tavg(\theta) \approx T_0$, so that
  total absorbed power reduces to $T_0\,\Aperp\,\Sinc$
& Within 5\,\% above 10\,GHz \\[4pt]
%
Diao~\textit{et al.}~\cite{Diao2024}:
  transmission coefficient $T \approx 0.52$ at 28\,GHz
  from full anatomical FDTD
& Normal-incidence Fresnel value
  $T_0 = 0.536$ for skin at 28\,GHz
& 3\,\% \\[4pt]
%
Flintoft~\textit{et al.}~\cite{Flintoft2014}:
  normalised ACS $\langle Q^a\rangle \approx 0.41$ at
  7--11\,GHz, 60~volunteers
& $\langle Q^a\rangle = T_0 \times \eta \approx
  0.48 \times 0.85 = 0.41$
& 3--8\,\% \\[4pt]
%
Zhang~\cite{Zhang2017thesis,ZhangRobinson2020}:
  absorption coefficient $\xi \approx 0.45$--0.65 above
  6\,GHz, 48~subjects (1--18\,GHz)
& $\xi = \Tbar \times \Aab/A$;
  predicted $\xi \approx 0.43$--0.49
& Within measurement scatter \\[4pt]
%
Flintoft~\textit{et al.}~\cite{Flintoft2014}:
  ACS dip near 3\,GHz, attributed to subcutaneous
  fat layer
& Fabry--P\'{e}rot model:
  $f_{\min} = c/(4\,n_f\,d_f) \approx 3.9$\,GHz
& Qualitative \\[4pt]
%
Li~\textit{et al.}~\cite{Li2019}:
  transmitted flux insensitive to incidence angle
  and polarisation, 30\,GHz--1\,THz
& Pseudo-Brewster: for tissue ($|\ntilde| > 2.5$),
  the ratio $\Tavg(\theta)/T_0$ stays near unity
& Within 5.6\,\% for skin \\[4pt]
%
Diao~\textit{et al.}~\cite{Diao2021}:
  heating factor approximately constant versus
  incidence angle at mmWave
& TE transmittance decreases and TM transmittance
  increases with angle; their average $\Tavg$
  remains near $T_0$
& Within 6\,\% for skin \\
\bottomrule
\end{tabularx}
\end{table}

% ------------------------------------------------------------------
% Closing paragraph
% ------------------------------------------------------------------

The results reviewed above span electromagnetic dosimetry, optics,
diffuse-field theory, integral geometry, and wireless communications.
In each case partial results exist: empirical observations, numerical
confirmations, or algebraic structures.
No single physical mechanism connects them;
\cref{tab:lit:comparison,tab:lit:predictions} summarise the
comparison.


% ============================================================
\section{Summary}\label{sec:summary}
% ============================================================

The absorbed power density on biological tissue reduces to a hierarchy of
closed-form expressions, one per level of approximation.
Parts~I--III derive each result in full; this section states the principal
findings with their accuracy and cross-references to where each argument lives.

% ------------------------------------------------------------------
\subsection{Part~I: A Millimetre-Wave Framework}
\label{sum:part1}
% ------------------------------------------------------------------

\subsubsection*{Local law and pseudo-Brewster compensation}

At a surface point~$\rr$ with outward normal~$\nhat(\rr)$, an incident
plane wave with power density~$\Sinc$ and direction~$\khat$ induces an
incidence cosine $\mu(\rr)=\nhat(\rr)\cdot(-\khat)$.
The wave decomposes locally into TE and TM components; the Fresnel-weighted
average $\Teff(\rr)=|e_s|^2T_s(\theta)+|e_p|^2T_p(\theta)$ is the
\emph{effective transmission}.
Energy conservation at a locally flat surface then gives the
\emph{exact absorption law} (\cref{sec:local-law}):
\begin{equation}\label{sum:eq:Sab-exact}
  \boxed{\Sab(\rr) = \Sinc\;\Teff(\rr)\;\ReLU\!\bigl(\mu(\rr)\bigr).}
\end{equation}
The ReLU gate is exact: back-facing points absorb nothing.

The effective transmission decomposes as (\cref{sec:Teff-decompose})
\begin{equation}\label{sum:eq:Teff-decomp}
  \Teff(\rr) = \Tavg\!\bigl(\theta(\rr)\bigr)
    + \frac{q(\rr)}{2}\,\Delta T\!\bigl(\theta(\rr)\bigr),
\end{equation}
where $\Tavg=\tfrac{1}{2}(T_s+T_p)$ is the unpolarised baseline and
$q\in[-1,1]$ is the local TM excess.
Two independent conditions each reduce \cref{sum:eq:Sab-exact} to the
\emph{geometric absorption law}:
\begin{equation}\label{sum:eq:geometric-law}
  \boxed{\Sab(\rr)
    \;\approx\; \Sinc\;\cdot\; T_0\;\cdot\;
    \ReLU\!\bigl[\nhat(\rr)\cdot(-\khat)\bigr].}
\end{equation}
First, for circularly polarised or unpolarised illumination,
$\langle q\rangle=0$ and the correction vanishes identically.
Second, for biological tissue at mmWave frequencies, $T_s(\theta)$
decreases monotonically with angle while $T_p(\theta)$ first rises then
falls; their average $\Tavg(\theta)$ stays close to $T_0$ across the
relevant angular range—the \emph{pseudo-Brewster compensation}
(\cref{sec:pseudo-brewster}).
Azzam~\cite{Azzam2015} showed this near-constancy holds whenever
$|\ntilde|>2.5$; biological tissue has $|\ntilde|\approx3$--6.
\Cref{tab:sum:materials} quantifies the compensation: $\Tavg/T_0$ stays
within 5.6\,\% for skin and within 8.2\,\% for all tissues listed.
All spatial variation in \cref{sum:eq:geometric-law} is determined by
the body's shape alone.

\begin{table}[!htbp]
\centering
\caption{Pseudo-Brewster compensation for representative tissues at
  28\,GHz.  $|\ntilde|$ is the complex refractive index modulus;
  $T_0$ is the normal-incidence power-absorption coefficient.
  The last column gives the variation of $\Tavg(\theta)/T_0$ over
  $\theta\in[0^\circ,75^\circ]$.  All tissues satisfy $|\ntilde|>2.5$,
  placing them in the Azzam~\cite{Azzam2015} regime.}
  \label{tab:sum:materials}
\begin{tabular}{lccr}
\toprule
Tissue & $|\ntilde|$ & $T_0$ & $\Tavg/T_0$ variation \\
\midrule
Skin   & 4.84 & 0.54 & $\le 5.6\,\%$ \\
Muscle & 5.62 & 0.48 & $\le 4.8\,\%$ \\
Fat    & 2.05 & 0.77 & $\le 8.2\,\%$ \\
Water  & 6.62 & 0.45 & $\le 3.9\,\%$ \\
\bottomrule
\end{tabular}
\end{table}

The approximation is exact to four significant figures near 40\,GHz;
below 40\,GHz it introduces a conservative error (underestimate);
above 80\,GHz the overestimate reaches 3.5\,\%.
For predictions below 10\,GHz, replace $T_0$ by the exact
flux-averaged transmission~$\Tbar$ (\cref{sec:sub6}).

\subsubsection*{Geometric consequences}

Integrating \cref{sum:eq:geometric-law} over the illuminated surface
(\cref{sec:geometry}) gives total absorbed power as $\Sinc T_0$ times
the \emph{projected area}
$\Aperp(\khat)\equiv\int_{\Sigma_+}\cos\theta\,\diff A$: the shadow cast
by the body perpendicular to~$\khat$.

For a non-convex body, some surface points are self-shadowed.

\begin{definition}[Exposure fraction]\label{sum:def:eta}
The \emph{exposure fraction} $\eta(\rr)$ is the cosine-weighted fraction
of the upper hemisphere from which $\rr$ is unoccluded:
\begin{equation}\label{sum:eq:eta-def}
  \eta(\rr) \;\equiv\; \frac{1}{\pi}\int_{S^2}
  \ReLU\!\bigl[\nhat(\rr)\cdot(-\khat)\bigr]\;
  O(\rr,\khat)\;\diff\Omega(\khat).
\end{equation}
For convex bodies $\eta\equiv1$; for concave regions $\eta<1$.
\end{definition}

This quantity is mathematically identical to \emph{ambient occlusion} in
computer graphics~\cite{Zhukov1998}, giving access to GPU implementations
that compute it at interactive frame rates.
The \emph{absorption area} $\Aab\equiv\int_\Sigma\eta(\rr)\,\diff A$
equals~$A$ for a convex body; for the Thelonious phantom
$\Aab/A\approx0.87$.

The direction-averaged total absorbed power satisfies a generalisation
of Cauchy's classical projected-area formula~\cite{Cauchy1841}
(proved in \cref{sec:cauchy}):
\begin{equation}\label{sum:eq:gen-cauchy}
  \boxed{\langle P_{\mathrm{abs}}\rangle
    = \Sinc\,T_0\,\frac{\Aab}{4}.}
\end{equation}
The classical Cauchy formula ($\langle\Aperp\rangle=A/4$ for convex
bodies) is the special case $\eta\equiv1$.
The \emph{absorption directivity} $D(\khat)\equiv4\Aperp(\khat)/\Aab$
factors the result as $P_{\mathrm{abs}}(\khat)=\Sinc\,T_0\,(\Aab/4)\,D(\khat)$
(\cref{sec:geometry}): the electromagnetic environment enters only
through~$\Sinc$; the tissue through~$T_0$; and the body geometry
through $\Aab$ and $D(\khat)$, both precomputed from the mesh.
Since $D(\khat)$ is smooth on $S^2$, a spherical harmonic expansion
with $L\le6$ (49~coefficients) achieves an RMS error below 1\,\%.

\subsubsection*{Compliance and computation}

Using $\SARwb=P_{\mathrm{abs}}/m$ and the hemisphere bound $D\le2$,
the worst-case whole-body SAR compliance criterion (\cref{sec:compliance}) is
\begin{equation}\label{sum:eq:compliance-worst}
  \boxed{\Sinc < \frac{0.16\,m}{T_0\,\Aab}.}
\end{equation}
Via the Du~Bois formula~\cite{DuBois1916}, the ratio $m/\Aab$ scales as
$\mathrm{BMI}^{0.575}h^{0.425}$, so children are the binding case
for population-level limits.

For $N$ paths with directions $\mathbf{K}\in\mathbb{R}^{3\times N}$,
surface normals $\mathbf{N}$, and path power densities $\mathbf{s}$,
the spatial dosimetry map (\cref{sec:computation}) is
\begin{equation}\label{sum:eq:Sab-mat}
  \mathbf{S}_{\mathrm{ab}} = T_0\;\ReLU(\mathbf{N}\mathbf{K})\;\mathbf{s},
\end{equation}
the forward pass of a single-hidden-layer ReLU neural network.
The analogy is structural, not approximate, and makes the dosimetry chain
differentiable through any ray-tracing engine.

For a point source at distance $d>3\lambda$, the law holds pointwise
with spatially varying inputs; integrating over the illuminated surface
yields the \emph{solid-angle formula}
$P_{\mathrm{abs}}=P_t\,T_0\,\Omega_{\mathrm{body}}/(4\pi)$
(\cref{sec:near-field}).
This is the emissive view-factor formula of radiative heat transfer, with~$T_0$
replacing emissivity.

% ------------------------------------------------------------------
\subsection{Part~II: Generalisations}
\label{sum:part2}
% ------------------------------------------------------------------

\subsubsection*{Polarisation correction}

Part~I set $q=0$.  Retaining it exactly and integrating over the body
surface yields the \emph{absorption Stokes vector} representation
(\cref{sec:polarisation}):
\begin{equation}\label{sum:eq:Pabs-Stokes}
  \boxed{P_{\mathrm{abs}} = \mathbf{m}(\khat)\cdot\mathbf{s}_{\mathrm{inc}},}
\end{equation}
where $\mathbf{s}_{\mathrm{inc}}=(S_0,S_1,S_2,S_3)^\top$ is the incident
Stokes vector and $\mathbf{m}(\khat)=(\tilde{A}_\perp,
-\tfrac{1}{2}B_c,-\tfrac{1}{2}B_s,0)^\top$ is the body's response vector.
Both $\tilde{A}_\perp\equiv\int_{\Sigma_+}\Tavg(\theta)\cos\theta\,\diff A$
(the Fresnel-weighted projected area, exact) and
$\tilde{B}=B_c+iB_s\equiv\int_{\Sigma_+}\Delta T(\theta)\cos\theta\,
e^{2i\alpha(\rr)}\,\diff A$ are precomputed from the body mesh and Fresnel
coefficients, independent of the incident polarisation state.
The zero fourth component confirms that circular polarisation handedness
is irrelevant.

The \emph{polarisation directivity} $D_B(\khat)\equiv|\tilde{B}|/(2\tilde{A}_\perp)$
bounds the fractional correction: the geometric worst case (an infinite
cylinder at broadside) gives $D_B^{\mathrm{cyl}}=27.9\,\%$ for skin at
28\,GHz; the maximum over all directions on the Thelonious phantom is
16.0\,\%.
In an indoor environment with $N\ge20$ mutually independent paths, the
standard deviation of the fractional polarisation error is
$D_B/\sqrt{2N}\le2.5\,\%$, comparable to the pseudo-Brewster
approximation itself.
The unpolarised geometric law is conservative for vertically polarised
illumination: horizontal limbs are purely TE, and $T_s(\theta)\le T_0$
for all $\theta$, so the simplified formula overestimates limb absorption.

\subsubsection*{Sub-6\,GHz extension}

The approximation $\Tavg(\theta)\approx T_0$ weakens below 10\,GHz.
The direction-averaged total power formula is made exact at any
frequency above ${\sim}100$\,MHz by replacing $T_0$ with the
\emph{flux-averaged transmission} (\cref{sec:sub6})
\begin{equation}\label{sum:eq:Tbar}
  \Tbar(f) \equiv 2\int_0^1 \Tavg(\mu,f)\;\mu\;\diff\mu,
\end{equation}
computed once per frequency from the Fresnel coefficients:
\begin{equation}\label{sum:eq:Cauchy-exact}
  \boxed{\langle P_{\mathrm{abs}}\rangle = \Sinc\,\Tbar\,\frac{\Aab}{4}.}
\end{equation}
See \cref{tab:Tbar} for $T_0$, $\Tbar$, and their ratio across
0.9--100\,GHz.
Below 40\,GHz the $T_0$-based formula underestimates absorbed power
and is therefore conservative for compliance.
The local surface map $\Sab(\rr)\approx\Sinc T_0\ReLU(\mu)$ loses
physical meaning below 6\,GHz when multi-layer resonances
(Fabry--P\'{e}rot effects between skin, fat, and muscle) become
significant; total power results via $\Tbar$ remain valid throughout.

% ------------------------------------------------------------------
\subsection{Part~III: Coherent MIMO Illumination}
\label{sum:part3}
% ------------------------------------------------------------------

\subsubsection*{Coherent absorption law}

Parts~I--II assume incoherent illumination: paths arrive with random
relative phases, so cross-terms average to zero.
In coherent MIMO beamforming with $M$ antenna elements and precoding
vector $\mathbf{x}\in\mathbb{C}^M$ ($\|\mathbf{x}\|^2=P$), fields add
as amplitudes; the coherent peak can exceed the incoherent sum by a
factor of~$N$ and the Part~I--II framework cannot capture it.

Each path~$n$ carries a \emph{polarisation-amplitude vector}
$\boldsymbol{\psi}_n\in\mathbb{C}^3$ encoding amplitude, polarisation,
and all propagation interactions (\cref{sec:coherent-setup}).
In the coherent regime, exact complex Fresnel amplitude coefficients
$t_s, t_p$ must be retained; replacing them by the real constant
$\sqrt{T_0}$ destroys the angle-dependent phases in cross-terms.
The \emph{Fresnel transmission operator} for path~$n$ at~$\rr$
(\cref{sec:fresnel-operator}) is
\begin{equation}\label{sum:eq:Fn}
  \mathbf{F}_n(\rr) =
  \Bigl(
    t_{s,n}\;\hat{e}_{s,n}\hat{e}_{s,n}^\top
    + t_{p,n}\;\hat{e}_{p,n}\hat{e}_{p,n}^\top
  \Bigr)H\!\bigl(\mu_n(\rr)\bigr).
\end{equation}

The absorbed power density is the Ohmic loss integrated over depth
(\cref{sec:coherent-absorption-law}); expanding the squared modulus
produces cross-terms weighted by depth coupling factors~$\Lambda_{nn'}$.
Two well-bounded approximations collapse this double sum to a squared
norm: Approximation~1 (replacing the refracted TM direction by the
incident direction) introduces $O(1/|\ntilde|^2)\lesssim4\,\%$ error on
cross-terms, with self-terms exact; Approximation~2 (setting all
normalised depth couplings $\Gamma_{nn'}=1$) is exact to $0.44\,\%$
at 28\,GHz.  Both affect only cross-terms; the combined error is
$\lesssim5\,\%$.
\hl{todo: both numbers in this paragraph are wrong. See the todos at
prop approx1 (bound unsound, O(1) per pair, rescued only on the
integrated operator) and prop approx2 (alpha-beta swap in
validate\_approx2.py, true error 2.5 percent at 28 GHz, up to 4
percent in band).}

Under these approximations, with exposure channel matrix
$\tilde{\mathbf{G}}(\rr)\in\mathbb{C}^{3\times M}$ collecting the
Fresnel-filtered, depth-weighted field contributions of each transmit
element (defined in full in \cref{sec:coherent-absorption-law}), the
collapsed result is the \emph{coherent absorption law}:
\begin{equation}\label{sum:eq:coherent-law}
  \boxed{\Sab(\rr) = \|\tilde{\mathbf{G}}(\rr)\,\mathbf{x}\|^2.}
\end{equation}
When paths have mutually random phases, cross-terms average to zero
and \cref{sum:eq:coherent-law} reduces to the Part~I multi-source
formula (\cref{sec:coherent-absorption-law}).

\subsubsection*{Exposure operator and beamforming}

Integrating \cref{sum:eq:coherent-law} over the body surface yields
the \emph{exposure operator} (\cref{sec:exposure-operator}):
\begin{equation}\label{sum:eq:Q-def}
  P_{\mathrm{abs}} = \mathbf{x}^H\,\mathbf{Q}\,\mathbf{x},
  \qquad
  \boxed{\mathbf{Q} \equiv \int_\Sigma
    \tilde{\mathbf{G}}(\rr)^H\,\tilde{\mathbf{G}}(\rr)\,\diff A
    \;\in\;\mathbb{C}^{M\times M}.}
\end{equation}
$\mathbf{Q}$ is Hermitian positive-semidefinite; its eigendecomposition
$\mathbf{Q}=\sum_k\lambda_k\mathbf{v}_k\mathbf{v}_k^H$ reveals the
\emph{exposure modes}.
The dominant eigenvector~$\mathbf{v}_1$ maximises body absorption;
$\lambda_{\max}$ is the maximum absorbed power per unit transmit power;
$\mathrm{tr}(\mathbf{Q})$ is the expected absorption under isotropic
transmission.
This eigendecomposition plays the same role in
Part~III as the absorption directivity $D(\khat)$ plays in Part~I.

Under maximum-ratio transmission (MRT), the absorbed power is
$P^{\mathrm{MRT}}_{\mathrm{abs}}=P\,\rho\,\lambda_{\max}(\mathbf{Q})$,
where the \emph{exposure-signal alignment} (\cref{sec:MRT-body-section})
\begin{equation}\label{sum:eq:rho-def}
  \boxed{\rho \equiv \frac{\mathbf{h}^\top\mathbf{Q}\,\mathbf{h}^*}
    {\|\mathbf{h}\|^2\,\lambda_{\max}(\mathbf{Q})} \;\in\;[0,1]}
\end{equation}
is the single number quantifying the exposure cost of MRT.
A small $\rho$ means the dominant signal paths arrive at grazing
incidence or from self-shadowed directions: body geometry gives inherent
protection.
The signal and body fields diverge structurally: the signal sums
scalars $h_j^*h_j$, while the body field sums vectors
$h_j^*\tilde{\mathbf{g}}_j(\rr)$.
Three structural differences prevent the body hotspot from coinciding
with the signal focus (\cref{sec:MRT-body-section}).
First, transmitted polarisations can partially cancel at the body surface.
Second, grazing-incidence paths contribute strongly to the signal but
weakly to absorption.
Third, additional per-path phases accumulate between the UE position
and any body-surface point.

The exposure-constrained beamforming problem (maximise received signal
power subject to a body absorption limit) is a QCQP with two PSD
constraints; the S-procedure guarantees a global closed-form solution
(\cref{sec:ecbf}):
\begin{equation}\label{sum:eq:optimal-x}
  \boxed{\mathbf{x}^\star = \sqrt{P}\;
    \frac{(\lambda\mathbf{Q}+\nu\mathbf{I})^{-1}\mathbf{h}^*}
         {\|(\lambda\mathbf{Q}+\nu\mathbf{I})^{-1}\mathbf{h}^*\|}.}
\end{equation}
The matrix $(\lambda\mathbf{Q}+\nu\mathbf{I})^{-1}$ suppresses
high-absorption modes by $\sim1/(\lambda\lambda_k)$; null-space
directions of~$\mathbf{Q}$ are unconstrained.
For MU-MIMO with $K$ users, each person~$u$ contributes an exposure
operator~$\mathbf{Q}^{(u)}$ built from the paths arriving at their
body; per-person constraints fit directly into WMMSE solvers
(\cref{sec:mu-mimo}).


% ============================================================
\section{Limitations}\label{sec:limitations-v2}
% ============================================================

The geometric dosimetry framework, developed across Parts~I--III,
has five irreducible limitations.  Each is specific, quantified, and
distinct from the relaxable specialisations treated in
\cref{sec:polarisation,sec:sub6,sec:coherent-absorption-law}.

First, for reactive near-field sources ($d < \lambda/(2\pi)$, i.e.\ 1.7\,mm at 28\,GHz),
evanescent waves and antenna--body impedance coupling require
full-wave simulation.  For all other distances, the framework
applies with spatially varying inputs (\cref{sec:near-field}).

Second, coherent multi-bounce interference between body parts is treated incoherently
(\cref{sec:inter-body}).  The body's curved geometry randomises the
phase of reflected waves at mmWave frequencies, making the
incoherent treatment physically justified.  In tight concavities
at sub-6\,GHz, standing waves could form, but the body-averaged
effect is sub-percent (\cref{sec:inter-body}).

Third, the geometric-optics framework systematically underestimates
absorption near edges and small features (fingers, ears) due to diffraction at shadow boundaries.  The
error is conservative (it underestimates) and exceeds 10\,\% for
body parts smaller than approximately $10\lambda$.
\Cref{sec:corrections} quantifies the diffraction smoothing.

Fourth, the surface absorption map $\Sab(\rr)$ loses physical meaning below 6\,GHz
when the skin depth exceeds the surface layer thickness.
Total absorbed power ($\SARwb$) remains valid
(\cref{sec:sub6}), but the depth-resolved SAR requires
volumetric computation.

Finally, there is an inherent tissue dielectric uncertainty.
All numerical results depend on the assumed dielectric model
(Gabriel 1996~\cite{Gabriel1996}, IT'IS v5.0~\cite{ITISv5}).  At
mmWave frequencies, measured skin permittivity varies by
approximately 20\,\% between
studies~\cite{AlekseevZiskin2007,Zhadobov2011}.
This propagates linearly into $T_0$ and hence into all power
predictions.


% ============================================================
\section{Conclusion}\label{sec:conclusion}
% ============================================================

The absorbed power density on biological tissue exposed to
electromagnetic radiation from 100\,MHz to 100\,GHz reduces to a
framework that separates material response from geometry.  Two
independent observations anchor the approach.

First, a pseudo-Brewster compensation between TE and TM
polarisations renders the unpolarised Fresnel transmission
approximately angle-independent for biological tissue, collapsing
the material factor to a single scalar~$T_0$.  At 40\,GHz the
compensation is exact; across the full 0.3--100\,GHz range the
RMS error remains below 5\,\%.  The exact averaged
transmission~$\Tbar$ eliminates this error entirely for
total-power quantities.

Second, for circularly polarised or unpolarised illumination, the
polarisation-dependent correction vanishes identically.  For
arbitrary polarisation, the total absorbed power is the inner product
of the incident Stokes vector with a precomputed absorption Stokes
vector~$\mathbf{m}(\khat)$.  The worst-case polarisation correction
on a human body is 16\,\%, bounded universally by 27.9\,\% (the
cylinder limit), and suppressed below 2.5\,\% in multipath
environments with 20 or more paths.

Under the mmWave specialisation, the local absorbed power density
becomes $\Sab(\rr)\approx\Sinc\,T_0\,\ReLU[\nhat(\rr)\cdot(-\khat)]$,
which depends on geometry alone.  Three consequences follow.
First, total absorbed power equals $\Sinc\,T_0$ times the projected
(shadow) area, connecting dosimetry to integral geometry.
Direction-averaged absorption satisfies a generalisation of Cauchy's
formula: $\langle P_{\mathrm{abs}}\rangle = \Sinc\,\Tbar\,\Aab/4$
(exact at any frequency), where the absorption area~$\Aab$ is the
surface area weighted by the exposure fraction~$\eta$, identical to
ambient occlusion in computer graphics.
Second, closed-form whole-body SAR compliance criteria depend only on
body mass, surface area, and $\Tbar$, with explicit anthropometric
scaling across the human population.
Finally, for multiple incident waves the formula is a
single-hidden-layer ReLU neural network, making the full dosimetry
chain differentiable and directly compatible with modern ray-tracing
engines for gradient-based optimisation of wireless systems subject
to exposure constraints.

Validation against exact Mie theory~\cite{BohrenHuffman1983} for
lossy spheres (a worst-case geometry) confirms framework errors of
order 10\,\% for body-scale objects at millimetre-wave frequencies.  The dominant uncertainty
remains the tissue dielectric properties themselves.  The framework
transforms a $10^{12}$-cell volumetric simulation into $O(M)$ surface
operations: a reduction of six orders of magnitude that preserves
the essential physics.

The five irreducible limitations (\cref{sec:limitations-v2}) are
reactive near-field coupling ($d < \lambda/(2\pi)$), coherent
multi-bounce interference, diffraction at shadow boundaries,
volumetric SAR below 6\,GHz, and tissue dielectric uncertainty.
All other restrictions of the geometric absorption law
(polarisation, frequency range, near-field sources) have been
relaxed in Parts~I--II without loss of the framework's computational
efficiency.

Part~III extends the framework to coherent MIMO beamforming,
where fields add as complex amplitudes rather than powers.  Two
well-bounded approximations collapse the coherent double sum to a
squared norm: $\Sab(\rr) =
\|\tilde{\mathbf{G}}(\rr)\mathbf{x}\|^2$.  Integrating over
the body yields the exposure operator~$\mathbf{Q}$, whose
eigendecomposition defines exposure modes in direct analogy to
Part~I's absorption directivity.  The exposure-signal
alignment~$\rho$ determines whether MRT is benign or hazardous for
body absorption; exposure-constrained beamforming reduces to a QCQP
with a closed-form solution, and MU-MIMO extends to $K$ simultaneous
users via per-person exposure operators.  The combined approximation
error is $\lesssim 5\,\%$, the same order as the Fresnel approximation
of Part~I.


% ============================================================
\partheading{Part~I: A Millimetre-Wave Framework}
{Sections~\ref{sec:local-law}--\ref{sec:applications}: a self-contained
dosimetry framework for the 10--60\,GHz band.}
% ============================================================

% ============================================================
\section{Local absorption law}\label{sec:local-law}
% ============================================================

The absorbed power density at a point on the body surface follows from
first principles.  The derivation below is exact; all approximations
are introduced in subsequent sections.

% -------------------------------------------------------
\subsection{Setup and notation}\label{sec:setup}

Consider a monochromatic plane wave propagating in direction $\khat$
with time-averaged Poynting vector
$\mathbf{S}_{\mathrm{inc}}=\Sinc\,\khat$ (units: W/m$^2$).  The wave
is incident on a body whose surface~$\Sigma$ is locally flat on the
scale of the wavelength. We assume no coherent reflections from internal tissue interfaces or other body parts; these are secondary effects at mmWave frequencies. This assumption is revisited in
\cref{sec:corrections}.

At a surface point~$\rr$ with outward unit normal~$\nhat(\rr)$, define
the incidence cosine
\begin{equation}\label{eq:mu-def}
  \mu(\rr) \;\equiv\; \nhat(\rr)\cdot(-\khat)
  \;=\; \cos\theta_i(\rr),
\end{equation}
where $\theta_i$ is the local incidence angle.  A front-facing point
has $\mu>0$; a point facing away from the source has $\mu\le 0$.

% -------------------------------------------------------
\subsection{Power flux and absorption}\label{sec:power-flux}

The inward power flux through a surface element $\diff A$ at $\rr$ is
\begin{equation}
  \diff P_{\mathrm{in}} = \Sinc\,\mu\,\diff A.
\end{equation}
This follows from the definition of the Poynting flux through an
oriented surface: the beam carries $\Sinc$\,W/m$^2$ normal to $\khat$,
but a tilted surface element intercepts only the
component~$\mu=\cos\theta_i$ per unit of its own area.

The reflected wave propagates in the specular direction, making the
same angle~$\theta_i$ with the normal.  Its power density is
$|r(\theta_i)|^2\,\Sinc$, where $r(\theta_i)$ is the Fresnel
amplitude reflection coefficient appropriate to the polarisation.  The
outward reflected flux is therefore
\begin{equation}
  \diff P_{\mathrm{ref}} = |r(\theta_i)|^2\,\Sinc\,\mu\,\diff A.
\end{equation}

For a lossy half-space in which the skin depth~$\delta$ is much
smaller than any body dimension (at mmWave frequencies,
$\delta\sim 0.3$ to $0.5$\,mm, far smaller than body curvature radii),
all transmitted power is absorbed within the surface layer.  Energy
conservation at the interface gives
\begin{equation}
  \diff P_{\mathrm{abs}} = \diff P_{\mathrm{in}} - \diff P_{\mathrm{ref}}
  = \Sinc\,\mu\,(1-|r|^2)\,\diff A.
\end{equation}

\begin{definition}[Power-absorption coefficient]
The power-absorption coefficient for a given polarisation is
\begin{equation}\label{eq:T-def}
  T(\theta_i) \;\equiv\; 1 - |r(\theta_i)|^2.
\end{equation}
\end{definition}

The absorbed power density is then
\begin{equation}\label{eq:Sab-exact-single-pol}
  \boxed{\Sab(\rr) = \Sinc \cdot T\!\bigl(\theta_i(\rr)\bigr)
    \cdot \mu(\rr).}
\end{equation}
This holds for each polarisation separately.  The quantity
$\Sab = \Sinc\cdot T(\theta)\cdot\cos\theta$ equals the transmitted
power density crossing the skin surface per unit skin area, often written
$S_{\mathrm{tr}}$ in the dosimetry literature~\cite{Li2019,Funahashi2018}.
Li~\textit{et al.}~\cite{Li2019} showed empirically that the ratio
$S_{\mathrm{tr}}/\Delta T_{\mathrm{skin}}$ (transmitted power density
per degree of surface temperature rise) is nearly angle- and
polarisation-independent from 6\,GHz to 1\,THz; the pseudo-Brewster
analysis of \cref{sec:pseudo-brewster} shows why.
We now derive the transmission coefficients.

% -------------------------------------------------------
\subsection{Fresnel coefficients for a lossy half-space}
\label{sec:fresnel}

The body surface is modelled as a planar interface~\cite{BornWolf1999}
between free space ($n_1=1$) and a lossy medium with complex refractive index
$\ntilde = \sqrt{\epstilde}$, where
\begin{equation}\label{eq:eps-complex}
  \epstilde = \varepsilon_r - i\frac{\sigma}{\omega\varepsilon_0}.
\end{equation}
Here $\varepsilon_r$ is the relative permittivity and $\sigma$ the
conductivity of the tissue~\cite{Gabriel1996,ITISv5}; $\omega=2\pi f$
is the angular frequency.

Snell's law generalises to complex angles
(\cref{app:snell}):
$\sin\theta_t = \sin\theta_i/\ntilde$.  Define the normal component
of the wave vector in the lossy medium,
\begin{equation}\label{eq:xi-def}
  \xi \;\equiv\; \ntilde\cos\theta_t
  \;=\; \sqrt{\ntilde^2 - \sin^2\theta_i}
  \;=\; \sqrt{\ntilde^2 - 1 + \mu^2},
  \qquad \RE(\xi)>0.
\end{equation}
For lossless media, $\xi$ is real; for lossy media, $\xi$ is complex
with positive real part, encoding exponential decay into the medium.

The Fresnel amplitude reflection coefficients follow from continuity
of the tangential electric and magnetic fields at the interface
(derived from Maxwell's equations in \cref{app:fresnel-derivation}).

\paragraph{TE (s) polarisation.}
The electric field is perpendicular to the plane of incidence.
Matching tangential $E$ and $H$ at the boundary yields
\begin{equation}\label{eq:rs}
  r_s = \frac{\mu - \xi}{\mu + \xi}.
\end{equation}
This has the form of an impedance mismatch: $\mu=\cos\theta_i$
characterises the free-space side, $\xi$ the medium side.

\paragraph{TM (p) polarisation.}
The magnetic field is perpendicular to the plane of incidence.
The constitutive relation $\mathbf{D}=\epstilde\,\varepsilon_0\,\EE$
introduces an additional factor of~$\ntilde^2$:
\begin{equation}\label{eq:rp}
  r_p = \frac{\ntilde^2\,\mu - \xi}{\ntilde^2\,\mu + \xi}.
\end{equation}

The corresponding power-absorption coefficients are
\begin{equation}\label{eq:Ts-Tp}
  T_s(\theta) = 1 - |r_s|^2
  = 1 - \left|\frac{\mu-\xi}{\mu+\xi}\right|^2, \qquad
  T_p(\theta) = 1 - \left|\frac{\ntilde^2\mu-\xi}{\ntilde^2\mu+\xi}\right|^2.
\end{equation}
Samaras and Kuster~\cite{Samaras2019} computed $T_s$ and $T_p$ separately
as functions of angle above 6\,GHz using multilayer planar skin models.
They noted that TM transmittance peaks near the pseudo-Brewster angle,
reaching as high as 0.92 for thick stratum-corneum skin at 10\,GHz; the
TE counterpart was not discussed as a compensating effect.

At normal incidence ($\mu=1$, $\xi=\ntilde$) both polarisations are
degenerate:
\begin{equation}\label{eq:T0}
  T_0 \;\equiv\; T_s(0) = T_p(0)
  = 1 - \left|\frac{1-\ntilde}{1+\ntilde}\right|^2
  = \frac{4\,\RE(\ntilde)}{|1+\ntilde|^2}
  = \frac{4n}{(1+n)^2+\kappa^2},
\end{equation}
where $\ntilde = n - i\kappa$ with $n,\kappa>0$.
This scalar represents the power-absorption coefficient at normal incidence ($\theta=0$).
In the dosimetry literature it often appears as $T_{\mathrm{tr}}$
(Kodera~\textit{et al.}~\cite{Kodera2024}) or
$T$ (Diao~\textit{et al.}~\cite{Diao2021,Diao2024}); those works extract it
numerically or analytically from one-dimensional slab models rather than from Fresnel theory.
For skin at 28\,GHz ($\varepsilon_r=16.55$, $\sigma=25.8$\,S/m): $\ntilde=4.47-1.85i$, giving
$T_0=0.536$.

% -------------------------------------------------------
\subsection{Exact polarisation-aware law}\label{sec:exact-pol}

A plane wave with propagation direction $\khat$ has its electric field
$\EE_0$ in the plane perpendicular to $\khat$; it is always fully
polarised.  At a surface point~$\rr$ with incidence
plane spanned by $\khat$ and $\nhat(\rr)$, define the local TE and TM
unit vectors
\begin{equation}\label{eq:es-ep}
  \hat{e}_s(\rr) = \frac{\khat\times\nhat}{|\khat\times\nhat|},
  \qquad
  \hat{e}_p(\rr) = \hat{e}_s\times\khat.
\end{equation}
Both lie in the plane perpendicular to $\khat$ and form an orthonormal
basis there.  The wave's electric field decomposes as
$\EE_0 = a_s\,\hat{e}_s + a_p\,\hat{e}_p$ (with $a_s, a_p \in \mathbb{C}$), and the local TE and TM
energy fractions are
\begin{equation}
  |e_s|^2 = \frac{|a_s|^2}{|\EE_0|^2}, \qquad
  |e_p|^2 = \frac{|a_p|^2}{|\EE_0|^2}, \qquad
  |e_s|^2 + |e_p|^2 = 1.
\end{equation}

The effective power-absorption coefficient at $\rr$ is
\begin{equation}\label{eq:Teff}
  \Teff(\rr) = |e_s(\rr)|^2\,T_s(\theta) + |e_p(\rr)|^2\,T_p(\theta).
\end{equation}
The exact absorbed power density is therefore
\begin{equation}\label{eq:Sab-exact}
  \boxed{\Sab(\rr) = \Sinc\;\Teff(\rr)\;\ReLU\!\bigl(\mu(\rr)\bigr),}
\end{equation}
where $\ReLU(\mu)=\max(0,\mu)$ gates illumination (front-facing versus
back-facing).  We call this the \emph{exact absorption law}.  It holds
for any polarisation, any frequency, and any body geometry where the
surface is locally flat on the scale of the wavelength.

% -------------------------------------------------------
\subsection{Decomposition into baseline and polarisation correction}
\label{sec:Teff-decompose}

Define the \emph{unpolarised (average) transmission} $\Tavg(\theta)$ and the
\emph{polarisation splitting} $\Delta T(\theta)$ as
\begin{equation}\label{eq:Tavg-DeltaT}
  \Tavg(\theta) \equiv \tfrac{1}{2}\bigl[T_s(\theta)+T_p(\theta)\bigr],
  \qquad
  \Delta T(\theta) \equiv T_p(\theta) - T_s(\theta).
\end{equation}
For lossy dielectrics with $|\ntilde|>1$ one has $\Delta T\ge 0$ at
all angles.  Define also the \emph{local TM excess}
\begin{equation}\label{eq:q-def}
  q(\rr) \equiv |e_p(\rr)|^2 - |e_s(\rr)|^2 \;\in [-1,1].
\end{equation}
Then \cref{eq:Teff} can be written as the exact identity
\begin{equation}\label{eq:Teff-decomp}
  \boxed{\Teff(\rr) = \Tavg(\theta)
    + \frac{q(\rr)}{2}\,\Delta T(\theta).}
\end{equation}
This separates the effective transmission into an
unpolarised baseline~$\Tavg(\theta)$, which depends only on
the incidence angle, and a polarisation correction proportional
to $q\,\Delta T$.  The splitting $\Delta T$ grows steeply with
angle (reaching 0.83 at $75^\circ$ for skin at 28\,GHz); Diao and
Li~\cite{Diao2021} observed this angular enhancement of TM
heating on skin but framed it as a local compliance problem rather
than a compensating mechanism.  The splitting is the foundation of
the absorption Stokes vector developed in
\cref{sec:stokes-vector}.

\begin{remark}[Normal incidence]
At normal incidence ($\theta=0$) the TE/TM distinction is degenerate:
$T_s(0)=T_p(0)=T_0$ and hence $\Delta T(0)=0$.  The polarisation
correction vanishes regardless of $q$, so $\Teff=T_0$ for any
polarisation at normal incidence.
\end{remark}

% -------------------------------------------------------
\subsection{TM excess in terms of the polarisation ellipse}
\label{sec:q-ellipse}

Let the incident wave have ellipticity angle~$\chi\in[-\pi/4,\pi/4]$
(with $\chi=0$ for linear and $|\chi|=\pi/4$ for circular
polarisation) and orientation angle~$\psi_0\in[0,\pi)$ of the
polarisation ellipse, both defined with respect to an arbitrary
reference basis $(\hat{e}_1,\hat{e}_2)$ perpendicular to $\khat$.
The local TE/TM basis at $\rr$ is rotated by an angle~$\alpha(\rr)$
relative to this reference basis.  Then (see \cref{app:q-proof}):
\begin{equation}\label{eq:q-exact}
  \boxed{q(\rr)
    = -\cos(2\chi)\,\cos\!\bigl(2\alpha(\rr)-2\psi_0\bigr).}
\end{equation}

This formula is exact. It splits into three independent factors
controlling the polarisation correction:

\begin{center}
\begin{tabular}{lll}
\toprule
Factor & Depends on & Range \\
\midrule
$\cos(2\chi)$ & Wave polarisation &
  $[0,1]$; vanishes for circular \\
$\cos(2\alpha-2\psi_0)$ & Surface geometry $+$ wave orientation &
  $[-1,1]$ \\
$\Delta T(\theta)$ & Tissue $+$ incidence angle &
  $[0,{\sim}\,0.77]$ at 28\,GHz \\
\bottomrule
\end{tabular}
\end{center}

The correction vanishes when \emph{any one} of three conditions holds.
First, for circular polarisation ($|\chi|=\pi/4$), $\cos(2\chi)=0$.
Second, for normal incidence ($\theta=0$), $\Delta T(0)=0$.
Third, under azimuthal averaging, if $\alpha$ varies uniformly across the
surface over a full $2\pi$ range, the cosine term averages to zero.

% -------------------------------------------------------
\subsection{Unpolarised and circularly polarised illumination}
\label{sec:unpol}

A single monochromatic plane wave is always fully polarised.  The term
``unpolarised'' refers to an ensemble or time average in which the
polarisation fluctuates randomly over the observation time; it is
described by a coherency matrix $\mathbf{J} = \frac{1}{2}S_{\mathrm{inc}}\mathbf{I}_2$. In
any local TE/TM frame this matrix is diagonal:
\begin{equation}
  \mathbf{J} = \frac{1}{2}S_{\mathrm{inc}} \begin{pmatrix} 1 & 0 \\ 0 & 1 \end{pmatrix}.
\end{equation}
Under this condition, the local intensities are equal,
$\langle|a_s|^2\rangle=\langle|a_p|^2\rangle=S_{\mathrm{inc}}/2$, hence
$\langle q\rangle=1/2-1/2=0$ and the effective transmission simplifies to
$\langle\Teff\rangle=\Tavg$ at every surface point.

For circular polarisation ($|\chi|=\pi/4$), \cref{eq:q-exact} gives
$q=0$ pointwise, so $\Teff=\Tavg$ everywhere. This holds pointwise, not
only on average. It holds for any surface geometry.

Under either condition the exact law \cref{eq:Sab-exact} becomes
\begin{equation}\label{eq:Sab-Tavg}
  \Sab(\rr) = \Sinc\;\Tavg\!\bigl(\theta(\rr)\bigr)\;
  \ReLU\!\bigl(\mu(\rr)\bigr).
\end{equation}
The material factor is now the angle-dependent but
polarisation-independent function~$\Tavg(\theta)$.  One further
simplification reduces it to a constant; this is the subject of the
next section.


% ============================================================
\section{Pseudo-Brewster compensation}\label{sec:pseudo-brewster}
% ============================================================

\subsection{Physical mechanism}\label{sec:pB-mechanism}

For a lossless dielectric interface, the TM reflection coefficient
vanishes at the Brewster angle
$\theta_B=\arctan(n_2/n_1)$~\cite{BornWolf1999}.  For a lossy
dielectric the reflection minimum is finite but can be very
small; the angle at which $|r_p|^2$ is minimised is called the
\emph{pseudo-Brewster angle}~$\theta_{pB}$~\cite{Potter1970,Ohman1977}.

\begin{proposition}[Pseudo-Brewster angle]\label{prop:pB}
For a lossy half-space with $|\ntilde|\gg 1$,
\begin{equation}\label{eq:theta-pB}
  \theta_{pB} \approx \arctan|\ntilde|.
\end{equation}
The approximation is accurate to within $1^\circ$ for $|\ntilde|>3$.
\end{proposition}

At this angle, $T_p$ reaches a maximum (often exceeding 95\%), while
$T_s$ is small (typically below 20%).  Their average
$\Tavg(\theta_{pB})\approx\frac{1}{2}(0+1)=0.5$ is close to the
normal-incidence value $T_0\approx 0.5$ to $0.6$ for biological tissue.

The underlying mechanism is a compensation: as the incidence angle
increases from zero, $T_s$ \emph{decreases} monotonically (more
reflection for TE) while $T_p$ \emph{increases} towards the
pseudo-Brewster peak (less reflection for TM).  Their average remains
approximately constant over a wide angular range.
Azzam~\cite{Azzam2015} showed that for lossless dielectric substrates
with $|\ntilde|>2.5$, the unpolarised reflectance is nearly
angle-independent.  Biological tissue at mmWave frequencies falls in
this regime ($|\ntilde|\approx 4$--6,
see~\cite{Gabriel1996,AlekseevZiskin2007}), and the resulting
near-constancy of $\Tavg(\theta)$ is the basis of the simplification
exploited here.  For these refractive indices, $\Tavg$ stays within a
few percent of~$T_0$ up to about $75^\circ$.

\subsection{Numerical verification}\label{sec:pB-numerics}

\Cref{tab:fresnel-skin} shows $T_s$, $T_p$, and $\Tavg$ for human
skin at 28\,GHz ($\varepsilon_r=17$, $\sigma=25$\,S/m, giving
$\ntilde=4.49-1.79i$ and $|\ntilde|=4.84$).

\begin{table}[!htbp]
\centering
\caption{Fresnel transmission coefficients for skin at 28\,GHz.
  The ratio $\Tavg/T_0$ stays within 5.6\% of unity up to
  $75^\circ$.}\label{tab:fresnel-skin}
\begin{tabular}{ccccc}
\toprule
$\theta$ & $T_s$ (TE) & $T_p$ (TM) & $\Tavg$ & $\Tavg/T_0$ \\
\midrule
$0^\circ$  & 0.539 & 0.539 & 0.539 & 1.000 \\
$30^\circ$ & 0.489 & 0.591 & 0.540 & 1.002 \\
$45^\circ$ & 0.422 & 0.666 & 0.544 & 1.010 \\
$60^\circ$ & 0.321 & 0.791 & 0.556 & 1.032 \\
$70^\circ$ & 0.233 & 0.902 & 0.568 & 1.054 \\
$75^\circ$ & 0.182 & 0.952 & 0.567 & 1.053 \\
\bottomrule
\end{tabular}
\end{table}

The maximum deviation of $\Tavg$ from $T_0$ over $[0^\circ,75^\circ]$
is 5.6\%.

\Cref{fig:apd-angle} illustrates the compensation graphically.  The
upper panel shows $T_s$, $T_p$, and $\Tavg$ as functions of incidence
angle for skin at 28\,GHz.  As $\theta$ increases, $T_s$ decreases
monotonically (more TE reflection) while $T_p$ increases towards the
pseudo-Brewster peak near $78^\circ$ (less TM reflection).  Their
average $\Tavg$ remains nearly flat out to about $75^\circ$.  The
lower panel shows the absorbed power density normalised to
$\Sinc$, i.e.\ $T(\theta)\cos\theta$, for each polarisation and
for the simplified formula $T_0\cos\theta$.  The unpolarised curve
(green) tracks the simplified formula (dashed grey) closely; the
small gap between them is the Fresnel approximation error.

\begin{figure}[!htbp]
  \centering
  \includegraphics{apd_angle_dependence.pdf}
  \caption{Pseudo-Brewster compensation for skin at 28\,GHz
    ($\ntilde=4.49-1.79i$, $T_0=0.539$).  The upper panel shows the
    Fresnel power-absorption coefficients $T_s$ (TE), $T_p$ (TM), and
    $\Tavg=\frac{1}{2}(T_s+T_p)$ versus incidence angle; the dashed line marks
    $T_0$.  The TE and TM curves diverge strongly, but their average stays
    within 5.6\% of $T_0$ up to $75^\circ$.  The lower panel shows the
    normalised absorbed power density $\Sab/\Sinc=T(\theta)\cos\theta$.  The
    simplified expression $T_0\cos\theta$ (dashed grey) closely matches the
    full unpolarised result (green), confirming the geometric absorption law.}
  \label{fig:apd-angle}
\end{figure}

\subsection{Sphere-averaged absorption ratio}\label{sec:sphere-ratio}

A sphere illuminated by a plane wave provides a convenient scalar
measure of the Fresnel approximation error because every incidence
angle from $0^\circ$ to $90^\circ$ is present on the illuminated
hemisphere.  Define the \emph{sphere ratio}
\begin{equation}\label{eq:R-sphere}
  R \;\equiv\;
  \frac{T_0 \cdot \pi a^2}{2\pi a^2\int_0^1 \Tavg(\mu)\,\mu\,\diff \mu}
  \;=\; \frac{T_0}{2\int_0^1\Tavg(\mu)\,\mu\,\diff \mu}.
\end{equation}
If $\Tavg$ were exactly equal to $T_0$, one would have $R=1$ (since
$\int_0^1\mu\,\diff \mu=\frac{1}{2}$).  For skin at 28\,GHz:

\begin{center}
\begin{tabular}{lc}
\toprule
Polarisation & $R$ \\
\midrule
TE only  & 1.37 \\
TM only  & 0.77 \\
Unpolarised & 0.99 \\
\bottomrule
\end{tabular}
\end{center}

The TE and TM ratios deviate substantially from unity but compensate
when averaged.  The unpolarised sphere ratio of 0.99 indicates that
the constant-$T_0$ approximation underestimates total absorbed power
on a sphere by about 1\%. This is a conservative error.

\subsection{Material universality}\label{sec:material-universality}

All biological tissues at mmWave frequencies cluster in the
$(n,\kappa)$ region where the pseudo-Brewster compensation is
effective~\cite{Gabriel1996,ITISv5,AlekseevZiskin2007}.
\Cref{tab:materials} lists the main parameters.

\begin{table}[!htbp]
\centering
\caption{Pseudo-Brewster compensation across tissue types at
  28\,GHz.  The $\Tavg$ variation is the maximum deviation of
  $\Tavg/T_0$ from unity over $[0^\circ,75^\circ]$.}
\label{tab:materials}
\begin{tabular}{lccccc}
\toprule
Tissue & $\varepsilon_r$ & $\sigma$ (S/m) & $|\ntilde|$ & $T_0$
  & $\Tavg$ variation \\
\midrule
Skin   & 17.0 & 25.0 & 4.84 & 0.54 & 5.6\% \\
Muscle & 25.0 & 30.0 & 5.62 & 0.48 & 4.8\% \\
Fat    &  4.0 &  2.0 & 2.05 & 0.77 & 8.2\% \\
Water  & 25.0 & 55.0 & 6.62 & 0.45 & 3.9\% \\
\bottomrule
\end{tabular}
\end{table}

Fat has the lowest $|\ntilde|$ and the largest variation (8.2\%), but
it is rarely the outermost tissue layer at mmWave exposure sites.

\subsection{Frequency dependence}\label{sec:freq-dep}

As the frequency changes, so do the tissue dielectric properties and
hence $|\ntilde|$ and the quality of the pseudo-Brewster compensation.
\Cref{tab:freq} shows the relevant quantities for skin.

\begin{table}[!htbp]
\centering
\caption{Frequency dependence of the pseudo-Brewster compensation for
  skin.  The framework is most accurate in the 20 to 60\,GHz band.}
\label{tab:freq}
\begin{tabular}{lccccc}
\toprule
Frequency & $|\ntilde|$ & $T_0$ & $R_{\mathrm{sphere}}$
  & $\Tavg$ variation \\
\midrule
6\,GHz   & 6.07 & 0.48 & 0.95 & 14.2\% \\
10\,GHz  & 5.87 & 0.49 & 0.96 & 12.9\% \\
28\,GHz  & 4.84 & 0.54 & 0.99 & 5.7\%  \\
40\,GHz  & 4.30 & 0.57 & 1.00 & 2.4\%  \\
60\,GHz  & 3.68 & 0.62 & 1.02 & 5.9\%  \\
77\,GHz  & 3.33 & 0.66 & 1.02 & 9.6\%  \\
100\,GHz & 3.00 & 0.70 & 1.03 & 13.4\% \\
\bottomrule
\end{tabular}
\end{table}

The compensation is optimal near 40\,GHz, where $R_{\mathrm{sphere}}=1$
and $\Tavg$ varies by only 2.4\%.  Below 10\,GHz and above 80\,GHz
the $\Tavg$ variation exceeds 10\%, and the constant-$T_0$
approximation should be used with care or replaced by the
angle-dependent~$\Tavg(\theta)$.

\subsection{Geometric absorption law}\label{sec:geometric-law}

Combining the two independent simplifications, $\Teff\to\Tavg$ for
unpolarised or circularly polarised illumination
(\cref{sec:unpol}), and $\Tavg(\theta)\approx T_0$ from the
pseudo-Brewster compensation (\cref{sec:pB-mechanism}), the
exact absorption law \cref{eq:Sab-exact} reduces to the
geometric absorption law already stated in
\cref{eq:geometric-law}:
\begin{equation*}
  \Sab(\rr) \;\approx\; \Sinc\;T_0\;\ReLU\!\bigl[\nhat(\rr)\cdot(-\khat)\bigr].
\end{equation*}

This is the starting point for the geometric and computational
developments that follow.  The tissue physics enters only through the
scalar~$T_0$; all spatial variation is determined by the body's shape.


% ============================================================
\section{Geometric framework}\label{sec:geometry}
% ============================================================

\subsection{Reduction to geometry}\label{sec:geom-reduction}

With $T_0$ fixed, the absorbed power density is entirely determined by
the function
\begin{equation}\label{eq:f-geom}
  f(\rr) = \ReLU\!\bigl[\nhat(\rr)\cdot(-\khat)\bigr],
\end{equation}
which depends only on the shape of the body (through the
surface-normal field~$\nhat(\rr)$) and the direction of the
wave ($\khat$).  Electromagnetic dosimetry at mmWave frequencies thus
reduces to a problem in integral geometry.

\subsection{Total absorbed power and the projected area}
\label{sec:total-power}

Integrating $\Sab$ over the illuminated surface~$\Sigma_+$ (where
$\mu>0$) and using $T(\theta)\approx T_0$:
\begin{equation}\label{eq:Pabs-Aperp}
  P_{\mathrm{abs}}(\khat)
  = \Sinc\,T_0 \int_{\Sigma_+}\!\!\cos\theta\,\diff A
  = \Sinc\,T_0\,\Aperp(\khat),
\end{equation}
where
\begin{equation}\label{eq:Aperp-def}
  \Aperp(\khat) \equiv \int_{\Sigma_+}\cos\theta\,\diff A
\end{equation}
is the projected area (shadow area) of the body perpendicular
to~$\khat$.  Total absorbed power equals incident power density times
$T_0$ times a shadow area.  Kodera~\textit{et al.}~\cite{Kodera2024}
write the same relationship as
$\SARwb = T_{\mathrm{tr}}\times\mathrm{PA}\times\mathrm{IPD}/W$,
where $\mathrm{PA}$ is the projected area (here $\Aperp$),
$\mathrm{IPD}$ is the incident power density ($\Sinc$), and $W$ is the
body mass.  Their coefficient $T_{\mathrm{tr}}$ is extracted numerically
from a one-dimensional FDTD model; the present result identifies it as the
Fresnel normal-incidence transmission~$T_0$ and explains its near-constancy
via pseudo-Brewster compensation.

For a convex body, $\Aperp(\khat)$ equals the area of the orthogonal
projection of the body onto the plane perpendicular to $\khat$.  For a
non-convex body, \cref{eq:Aperp-def} overcounts (it includes
self-occluded front-facing elements); the true shadow area is the area
of the projection, which accounts for overlap.  The relationship
$P_{\mathrm{abs}}=\Sinc\,T_0\,\Aperp$ remains valid provided
$\Aperp$ is understood as the true projected area. A flux argument
guarantees this for any body.

\subsection{Self-shadowing and occlusion}\label{sec:self-shadow}

A human body is not convex: concavities such as armpits, the gap
between the legs, and the neck region can cause one part of the body to
shadow another.

\begin{definition}[Occlusion factor]
For a surface point~$\rr$ and an incidence direction~$\khat$, the
occlusion factor is
\begin{equation}\label{eq:O-def}
  O(\rr,\khat) = \begin{cases}
    1 & \text{if $\rr$ is visible from direction $\khat$
          (not blocked by another body part),}\\
    0 & \text{otherwise.}
  \end{cases}
\end{equation}
\end{definition}

For a non-convex body the local $\Sab$ becomes
\begin{equation}\label{eq:Sab-occluded}
  \Sab(\rr) = \Sinc\,T_0\,\ReLU\!\bigl[\nhat(\rr)\cdot(-\khat)\bigr]
  \;\cdot\; O(\rr,\khat).
\end{equation}
Computing $O$ requires ray casting against the body mesh. This is the
same operation a rendering engine performs. For convex bodies,
$O\equiv 1$ and the formula reduces to the geometric absorption law.

\begin{definition}[Exposure fraction]\label{def:eta}
The exposure fraction at a surface point~$\rr$ is
\begin{equation}\label{eq:eta-def}
  \eta(\rr) = \frac{1}{\pi}\int_{S^2}
  \ReLU\!\bigl[\nhat(\rr)\cdot(-\khat)\bigr]\;O(\rr,\khat)\;\diff \Omega(\khat).
\end{equation}
\end{definition}

This is the fraction of the cosine-weighted hemisphere from which
$\rr$ is visible.  For convex bodies, $O\equiv 1$ and $\eta\equiv 1$.
For non-convex bodies, $\eta<1$ in concavities (armpits, between the
legs) and $\eta=1$ on exposed surfaces (top of the head, shoulders).
On an anatomical phantom (Thelonious, 23\,826 triangles), the
distribution of $\eta$ has mean 0.80, median 0.92, with 26\% of
triangles having $\eta=1$ and a long tail extending to $\eta=0$ in
deep concavities.  Area-weighted, the mean rises to 0.87 and 38\% of
the surface area is fully exposed, because the low-$\eta$ concavities
are finely meshed while exposed regions carry larger triangles.

The exposure fraction~$\eta(\rr)$ is mathematically identical to the
ambient occlusion factor in computer
graphics~\cite{Zhukov1998}.  Zhukov~\textit{et al.}~\cite{Zhukov1998}
introduced the concept; Landis~\cite{Landis2002} coined the term
``ambient occlusion'' and brought it into production rendering.
This identification imports four decades of GPU-optimised algorithms
(BVH construction, ray casting, importance
sampling~\cite{Pharr2016,AkenineMoller2018}) directly into dosimetric
computation.
In the dosimetry literature the factor $\Aab/A = \int_\Sigma\eta\,\diff A/A$
appears as a posture-dependent ``self-shadowing correction'' $\gamma_s$
in Flintoft~\textit{et al.}~\cite{Flintoft2014},
where $\langle\sigma_a\rangle = \langle Q^a\rangle\cdot\gamma_s\cdot\mathrm{BSA}/4$
is the measured direction-averaged absorption cross-section and
$\langle Q^a\rangle$ is the absorption efficiency.
The framework identifies $\langle Q^a\rangle = T_0$ and
$\gamma_s = \Aab/A$, giving both quantities a physical derivation rather than
empirical status, and identifying it as standard ambient occlusion in computer graphics.

\subsection{Cauchy's formula and its generalisation}\label{sec:cauchy}

Cauchy's formula~\cite{Cauchy1841} is a classical result of integral
geometry~\cite{Hadwiger1957,SchneiderWeil2008}: the direction-averaged
projected area of a convex body equals one quarter of its surface area.
We extend it to non-convex bodies via the exposure
fraction~$\eta(\rr)$.

\begin{remark}[Motivation for directional averaging]
Direction-averaged absorbed power is the natural metric whenever the
illumination is statistically isotropic.  The most important example
is the reverberation chamber, widely used in bioelectromagnetics
for whole-body exposure of small animals and phantoms.  Inside a
well-stirred chamber the field arrives from all directions with equal
probability and uniform polarisation.  The direction-averaged quantity
$\langle P_{\mathrm{abs}}\rangle$ is then the actual absorbed power,
not merely a statistical summary.  The same situation arises in richly
scattered indoor environments and in any scenario where the
illumination direction is unknown or variable (a pedestrian in a city,
a worker moving through an industrial hall).  The framework developed
below gives a closed-form answer for
$\langle P_{\mathrm{abs}}\rangle$ in terms of a single geometric
scalar, the absorption area~$\Aab$.
\end{remark}

\begin{theorem}[Cauchy's formula for convex bodies]
For a convex body with surface area~$A$, the direction-averaged
projected area is
\begin{equation}\label{eq:cauchy}
  \langle \Aperp\rangle
  \equiv \frac{1}{4\pi}\int_{S^2}\Aperp(\khat)\,\diff\Omega
  = \frac{A}{4}.
\end{equation}
\end{theorem}

\begin{proof}
Exchange the order of integration (Fubini, valid since the integrand
is non-negative):
\[
  \langle\Aperp\rangle
  = \frac{1}{4\pi}\int_{S^2}\int_\Sigma[\nhat\cdot(-\khat)]_+\,\diff A\,\diff\Omega
  = \frac{1}{4\pi}\int_\Sigma
    \biggl(\int_{S^2}[\nhat\cdot(-\khat)]_+\,\diff\Omega\biggr)\,\diff A.
\]
The inner integral is independent of $\nhat$ by rotational symmetry
and equals $\pi$ (integrate $\cos\theta\sin\theta$ over the
hemisphere: $2\pi\int_0^{\pi/2}\cos\theta\sin\theta\,\diff\theta=\pi$).
Hence $\langle\Aperp\rangle = (\pi/4\pi)\,A = A/4$.
\end{proof}

\begin{definition}[Absorption area]\label{def:Aab}
For a body with surface $\Sigma$ and exposure fraction
$\eta(\rr)$, the \emph{absorption area} is
\begin{equation}\label{eq:Aab-def}
  \Aab \equiv \int_\Sigma\eta(\rr)\,\diff A.
\end{equation}
\end{definition}

For convex bodies, $\eta\equiv 1$ and $\Aab=A$.  For non-convex
bodies, $\Aab < A$; the deficit measures how much surface area is
hidden from isotropic illumination.

\begin{theorem}[Generalised Cauchy formula]\label{thm:gen-cauchy}
For any body (convex or not), the direction-averaged absorbed power is
\begin{equation}\label{eq:gen-cauchy}
  \boxed{\langle P_{\mathrm{abs}}\rangle
    = \Sinc\,T_0\,\frac{\Aab}{4}.}
\end{equation}
\end{theorem}

This is exact and follows from the same Fubini argument as the
classical Cauchy formula, with the occlusion factor~$O$ restricting
the integration to visible surface elements.  The absorption
area~$\Aab$ absorbs all the geometric complexity of self-shadowing
into a single scalar.  Replacing $T_0$ with the exact averaged
transmission $\Tbar$ (\cref{sec:Tbar}) makes the formula exact at
any frequency, without the pseudo-Brewster approximation.

For non-convex bodies, $\Aab\le A$ (since $\eta\le 1$).
Let $A_{\mathrm{CH}}$ denote the surface area of the convex hull.
For bodies with thin protruding limbs separated by gaps (such as
a standing human), $A_{\mathrm{CH}}$ can exceed~$A$ because the
hull's bridging surfaces across the gaps are larger than the
concavity surfaces they replace ($A_{\mathrm{CH}}/A\approx 1.20$
for Thelonious).  Cauchy's formula applied to the convex hull gives
a conservative upper bound:
$\langle P_{\mathrm{abs}}\rangle \le \Sinc\,T_0\,A_{\mathrm{CH}}/4$.

\subsection{Shadow deficit}\label{sec:shadow-deficit}

The projected area of a body is bounded by that of its convex hull:
$\Aperp^{\mathrm{body}}(\khat)\le\Aperp^{\mathrm{CH}}(\khat)$.  The
difference
\begin{equation}
  \Delta\Aperp(\khat)
  = \Aperp^{\mathrm{CH}}(\khat) - \Aperp^{\mathrm{body}}(\khat)
\end{equation}
arises from lines parallel to $\khat$ that pass through the convex
hull without intersecting the body. These are windows through open
channels such as the gap between the legs or between an arm and the
torso.  For a standing human body, the direction-averaged shadow
deficit is of order 2 to 3\% of the convex-hull shadow area, because
each window is visible only from a limited range of directions.

\subsection{Inter-body reflections}\label{sec:inter-body}

The geometric absorption law captures first-bounce absorption only.  At each
illuminated point, a fraction $R_{\mathrm{avg}}\approx 1-T_0\approx
0.46$ of the incident power is reflected.  When one body part's
reflection re-illuminates another (for example across the gap between
the legs or between an arm and the torso), the framework systematically
underestimates the absorbed power.  This subsection quantifies the
effect and shows it is negligible.

\begin{remark}[Recapture fraction and radiosity]
Let $f(\rr)$ denote the fraction of power reflected from~$\rr$ that
strikes another part of the body.  For a convex body, every reflected
ray escapes to infinity, so $f\equiv 0$ and the first-bounce law is
exact.  For a non-convex body, the exposure fraction already provides
a bound: under diffuse (Lambertian) reflection, the recaptured
fraction equals the fraction of the cosine-weighted hemisphere blocked
by other body parts, giving
\begin{equation}\label{eq:f-bound}
  f(\rr)\;\le\;1-\eta(\rr).
\end{equation}
For specular reflection (the dominant regime at millimetre
wavelengths, where skin roughness satisfies the Rayleigh criterion), the
recapture fraction is lower still, because a single reflected ray
either hits the body or escapes.

Summing the geometric series of successive bounces, the total absorbed
power at a point with recapture fraction~$f$ is enhanced by
\begin{equation}\label{eq:radiosity-C}
  C(f) = \frac{1}{1-\bar R\,f},
\end{equation}
where $\bar R=1-\Tbar\approx 0.46$ is the flux-weighted reflectance.
For the Thelonious phantom at 28\,GHz, the area-weighted global
recapture fraction is $f_{\mathrm{global}}\approx 0.09$ (diffuse
bound), yielding $C\approx 1.04$.  The specular estimate
($f_{\mathrm{spec}}\approx 0.3\,f_{\mathrm{diffuse}}$) reduces this
to $C\approx 1.01$: a body-averaged enhancement below 2\,\%.
\end{remark}

\begin{lemma}[Convex-hull energy bound]
A parameter-free upper bound on the total absorbed power follows from
energy conservation.  Under isotropic illumination, the power entering
the convex hull is $\Sinc\,A_{\mathrm{CH}}/4$ (Cauchy's formula
applied to the hull).  Because all power must either be absorbed or
exit through the hull,
\begin{equation}\label{eq:CH-bound}
  \langle P_{\mathrm{abs}}\rangle \;\le\; \Sinc\,\frac{A_{\mathrm{CH}}}{4}.
\end{equation}
For Thelonious, $A_{\mathrm{CH}}/A\approx 1.20$; the excess arises
from filling the gaps between the legs and between the arms and torso.
Combined with the first-bounce Cauchy formula, this brackets the true
absorbed power to within a few percent.
\end{lemma}

\begin{proposition}[Self-compensation mechanism]
A self-compensating mechanism limits the practical impact of inter-body
reflections.  Without reflections, the absorbed power density at a
point scales with~$\eta$; including all bounces it scales with
$\eta\cdot C$, where $C=1/(1-\bar R\,f)$ from \cref{eq:radiosity-C}.
The bound $f\le 1-\eta$ (\cref{eq:f-bound}) means that $C$ is largest
precisely where $\eta$ is smallest: a deep concavity sees less direct
power but recaptures more reflected power, so the two effects partially
cancel.  Concretely, for a point with $\eta=0.3$ (a deep concavity
such as between the legs), the diffuse bound gives
$f\le 0.7$, hence $C\le 1/(1-0.46\times 0.7)=1.47$.  The effective
exposure fraction $\eta\cdot C\le 0.44$, only 14~percentage points
above the uncorrected~$\eta=0.3$, recovering less than 20\,\% of the
70\,\% shadow deficit.  More realistically, specular reflection
reduces $f$ by roughly a factor of three, giving $C\approx 1.12$ and
$\eta\cdot C\approx 0.34$: reflections fill in less than 5\,\% of
the deficit.  On the Thelonious phantom, this pattern holds for every
concavity region.
\end{proposition}

Because the body-averaged correction ($\lesssim 1$\,\%) falls well
below diffraction (${\sim}10$\,\%) and tissue-property uncertainties
(${\sim}10$\,\%), inter-body reflections are omitted from the
framework's equations.  Their quantitative insignificance is included
in the error budget (\cref{sec:error-budget}).

\subsection{Absorption directivity}\label{sec:directivity}

\begin{definition}[Absorption directivity]
The \emph{absorption directivity} is
\begin{equation}\label{eq:D-def}
  D(\khat) = \frac{\Aperp(\khat)}{\langle\Aperp\rangle}
  = \frac{4\,\Aperp(\khat)}{\Aab}.
\end{equation}
\end{definition}

By construction, $\langle D\rangle=1$.  The directivity is a
dimensionless, shape-dependent but size-independent function on the
sphere~$S^2$.  It separates the body's response into
two parts: overall size, captured by $\Aab$ (or $A$ for convex bodies), and
shape anisotropy, captured by $D(\khat)$.

In terms of $D$, the total absorbed power factorises as
\begin{equation}\label{eq:Pabs-D}
  \boxed{P_{\mathrm{abs}}(\khat) = \Sinc\,T_0\,
    \frac{\Aab}{4}\;D(\khat).}
\end{equation}

\begin{proposition}[Separation into independent factors]
\Cref{eq:Pabs-D} makes the key structural point explicit: whole-body
absorption factorises into \emph{independent} contributions,
\[
  P_{\mathrm{abs}}(\khat)
  = \underbrace{\Sinc}_{\substack{\text{environment}\\\text{(field strength)}}}
    \cdot \underbrace{T_0}_{\substack{\text{material}\\\text{(tissue)}}}
    \cdot \underbrace{\frac{\Aab}{4}}_{\substack{\text{body size}\\\text{(geometry only)}}}
    \cdot \underbrace{D(\khat)}_{\substack{\text{posture/shape}\\\text{$+$ direction}}}.
\]
The environment enters only through the scalar incident power density
$\Sinc$; the tissue enters only through the scalar transmission~$T_0$.
All geometric complexity is isolated in the two purely geometric
quantities $\Aab$ (overall size) and $D(\khat)$ (directional anisotropy),
both precomputable from the body mesh and pose.
\end{proposition}

\begin{proposition}[Bounds on absorption directivity]
For any convex body, $\Aperp(\khat)\le A/2$ (the shadow area is at
most half the surface area), hence $D(\khat)\le 2$.  For non-convex
bodies the bound becomes $D\le 2A_{\mathrm{CH}}/\Aab$.
For Thelonious, $A_{\mathrm{CH}}/\Aab\approx 1.38$, giving
$D\le 2.76$; the actual maximum is~1.2, so the bound is loose.  From below, $\langle D\rangle=1$ implies
$D_{\max}\ge 1$; for a sphere, $D=1$ uniformly.
\end{proposition}

\paragraph{Human body values.}
For a standing adult (Thelonious phantom):
$D_{\max}\approx 1.22$ (side-on), $D_{\min}\approx 0.45$ (overhead).
The side-to-top ratio is roughly $2.7{:}1$ ($\approx 4.3$\,dB),
quantifying the directional sensitivity of whole-body absorption.  This
ratio is preserved across body sizes if the proportions are similar.
\Cref{fig:directivity} shows $D(\khat)$ on a Mollweide projection.
The equatorial belt (horizontal illumination) is near or above unity;
the poles (overhead and underfoot) are low.  The human body absorbs
most when illuminated from the side, least from above.

\begin{figure}[!htbp]
  \centering
  \includegraphics{body_directivity_thelonious.pdf}
  \caption{Absorption directivity $D(\khat)$ for the Thelonious
    phantom, shown in Mollweide projection.  The equator corresponds
    to horizontal illumination; the poles to vertical.  Side-on
    illumination ($D\approx 1.2$) produces roughly 2.7 times more
    absorbed power than overhead ($D\approx 0.45$).  The pattern is
    smooth, with modest front-back asymmetry due to anatomical
    detail.}
  \label{fig:directivity}
\end{figure}

\Cref{fig:directivity-hist} shows the distribution of $D$ values over
all sampled directions.  The histogram is strongly right-skewed: most
directions yield $D > 1$ because horizontal illumination covers a
larger fraction of the sphere than vertical.  The mode lies near
$D \approx 1.2$ (side-on), and the long tail extends down to
$D \approx 0.45$ (overhead).

\begin{figure}[!htbp]
  \centering
  \includegraphics{body_directivity_hist_thelonious.pdf}
  \caption{Histogram of the absorption directivity $D(\khat)$ for the
    Thelonious phantom over 4096 uniformly sampled directions.  The
    distribution is right-skewed: horizontal illumination directions
    (large~$D$) dominate the solid angle.  The mean is 1 by
    construction.}
  \label{fig:directivity-hist}
\end{figure}

\paragraph{Spherical harmonic representation.}
Since $D(\khat)$ is a smooth function on~$S^2$, it is efficiently
represented in Spherical Harmonics (SH):
\begin{equation}\label{eq:D-SH}
  D(\khat) \approx \sum_{l=0}^{L}\sum_{m=-l}^{l}
  d_{lm}\,Y_{lm}(\khat).
\end{equation}
The projected area of a convex body satisfies
$\Aperp(\khat)=\Aperp(-\khat)$, so $D$ is even and only even-$l$
harmonics contribute.  For a human body, $L\le 6$ (49~coefficients)
captures $D$ to better than 1\% RMS.  Total absorbed power then
reduces to a dot product of spherical-harmonic coefficient vectors.
This costs $O(L^2)$ per evaluation.

\Cref{fig:sh-fit} confirms this convergence.  The RMS fit error drops
below 3\% at $L=2$ (9~coefficients), below 1\% at $L=4$
(25~coefficients), and below 0.5\% at $L=6$ (49~coefficients).
The rapid convergence reflects the smoothness of $D(\khat)$: a
human body, despite its geometric complexity, presents a smooth
directional response because the projected area integrates out local
surface detail.

\begin{figure}[!htbp]
  \centering
  \includegraphics{body_directivity_sh_fit_thelonious.pdf}
  \caption{Spherical-harmonic fit error for the absorption
    directivity $D(\khat)$ of the Thelonious phantom, as a function
    of the maximum harmonic degree~$L$.  RMS error (green), 99th
    percentile error (blue), and maximum error (red) all
    decrease rapidly.  Errors are expressed relative to the mean
    directivity (percent).  At $L=6$ (49~coefficients) the maximum
    error is roughly 2\%.}
  \label{fig:sh-fit}
\end{figure}

\begin{remark}[Posture sensitivity]
The directivity values above are for a single standing pose.  To test
whether $D(\khat)$ changes with posture, 51~frames were sampled from
a backflip animation exported from Blender.  The body model differs
from the Thelonious phantom used elsewhere in this work; it has a
narrower directivity range ($D_{\max}\approx 1.06$--$1.13$ standing,
compared with ${\approx}\,1.22$ for Thelonious).  Across all 51
frames, $D_{\max}$ varies from 1.057 to 1.126, a spread of 6.5\%.
The median stays within 1\% of unity, and the spherical-harmonic fit
at $L=6$ has RMS error below 0.13\% for every frame.
\Cref{fig:posture-sensitivity} summarises the result.  The directivity
is not sensitive to posture, even during an extreme whole-body motion.
A single precomputed directivity table therefore remains a reasonable
approximation for compliance assessment.
\end{remark}

\begin{figure}[!htbp]
  \centering
  \includegraphics{posture_sensitivity.pdf}
  \caption{Posture sensitivity of the absorption directivity.  Top:
    representative silhouettes from a 51-frame backflip animation.
    Bottom: per-frame summary statistics of $D(\khat)$.  The shaded
    band spans $[D_{\min},\,D_{\max}]$.  $D_{\max}$ varies by 6.5\%
    across frames; the median stays near unity.}
  \label{fig:posture-sensitivity}
\end{figure}

\subsection{Arbitrary angular illumination}\label{sec:angular-spectra}

The isotropic average rarely describes a real deployment, but it serves
as a normalisation.  For a general angular power
spectrum~$\rho(\khat)$ on the sphere, normalised by
$\int_{S^2}\rho\,\diff \Omega=4\pi$ (so that $\rho=1$ is isotropic), the
total absorbed power is
\begin{equation}\label{eq:Pabs-rho}
  \boxed{P_{\mathrm{abs}}
    = \Sinc\,T_0\,\frac{\Aab}{4}\;
    \frac{1}{4\pi}\int_{S^2}\rho(\khat)\,D(\khat)\,\diff \Omega.}
\end{equation}

The integral $\frac{1}{4\pi}\int\rho\,D\,\diff\Omega$ is a scalar
coupling factor~$F$ that quantifies how much the actual illumination
pattern differs from the isotropic average.  For a standing human
body: an elevation-concentrated spectrum (modelling a base station)
yields $F\approx 1.15$; an urban-canyon ring of horizontal paths
gives $F\approx 1.18$; an isotropic upper hemisphere gives
$F\approx 1.00$ (because $D(\khat)\approx D(-\khat)$ for a nearly
symmetric body).  All these corrections are cheap to
evaluate from the precomputed $D(\khat)$ table.


% ============================================================
\section{Whole-body SAR compliance}\label{sec:compliance}
% ============================================================

\subsection{Closed-form compliance criterion}\label{sec:SAR-criterion}

The ICNIRP guidelines~\cite{ICNIRP2020} specify a whole-body average
SAR limit of $0.08$\,W/kg for the general public.  At mmWave
frequencies all transmitted power is absorbed in the skin layer, so
$\SARwb = P_{\mathrm{abs}}/m$, where $m$ is the body mass.  Using
\cref{eq:Pabs-D} and the conservative bound $\Aab\le A$:
\begin{equation}\label{eq:SARwb}
  \SARwb(\khat) \;\le\;
  \frac{\Sinc\,T_0}{4}\;\frac{A}{m}\;D(\khat).
\end{equation}

This has the same factorised structure as \cref{eq:Pabs-D}, now with a
mass normalisation: (environment) $\Sinc$, (material) $T_0$, (body size)
$A/m$ (conservatively; exact would use $\Aab/m$), and (posture/direction)
$D(\khat)$.

The condition $\SARwb<0.08$\,W/kg for direction~$\khat$ is guaranteed
whenever
\begin{equation}\label{eq:compliance-D}
  \Sinc < \frac{0.32\,m}{T_0\,A\,D(\khat)}.
\end{equation}

For worst-case exposure, using $D_{\max}\le 2$:
\begin{equation}\label{eq:compliance-worst}
  \boxed{\Sinc < \frac{0.16\,m}{T_0\,A}.}
\end{equation}

This depends on only three quantities: body mass~$m$, body surface
area~$A$ (BSA), and the tissue transmission coefficient~$T_0$ (which is
frequency-dependent).  The body surface area is defined here as the total
surface area of the skin, estimated from height and weight via the
Du~Bois formula (see \cref{sec:anthropometric}); it ranges from
roughly $0.4$\,m$^2$ for an infant to $2.0$\,m$^2$ for a large adult.
For linearly polarised illumination, the
compliance threshold shifts by at most 12\,\%
(\cref{sec:exact-bounds}).  At sub-6\,GHz, replacing $T_0$ with the
exact averaged transmission~$\Tbar$ gives exact results
(\cref{sec:Tbar}).

Gosselin~\textit{et al.}~\cite{Gosselin2011} derived a structurally equivalent
formula for whole-body SAR from base-station antennas at 300\,MHz--5\,GHz.
Their formula uses the normal-incidence transmission coefficient~$|t|^2$
computed from the tissue dielectric properties and
the cross-section-to-mass ratio of the body, validated against 432 FDTD
configurations across six frequencies, three human models, and twelve antennas.
A geometric factor $H_{\mathrm{eff}}/H_{\mathrm{body}}$ accounts for
partial illumination of the body by the antenna beam, where $H_{\mathrm{eff}}$
is the lesser of the body height, the antenna height, and the beam spread
at distance~$d$; this plays an analogous role to the directivity
factor~$D(\khat)$ in \cref{eq:compliance-D}, though $D(\khat)$ captures
the angular dependence of the body's projected area rather than partial
illumination.
\Cref{eq:compliance-worst} recovers the same transmission-times-area-over-mass
structure and extends its valid range to 100\,GHz.

\begin{remark}[Isotropic compliance]
In a reverberation chamber the illumination is isotropic, so
$D(\khat)$ averages to unity and the worst-case bound is not needed.
The exact whole-body SAR is $\SARwb = S_{\mathrm{total}}\,T_0\,\Aab/(4m)$.
Compliance is guaranteed whenever
$S_{\mathrm{total}} < 0.32\,m/(T_0\,\Aab)$.  This formula requires
only the precomputed absorption area~$\Aab$ and the body mass.
\end{remark}

\subsection{Body-size dependence}\label{sec:body-size}

\Cref{tab:compliance} evaluates \cref{eq:compliance-worst} for
representative body types at 28\,GHz ($T_0=0.536$).

\begin{table}[!htbp]
\centering
\caption{Worst-case whole-body SAR compliance threshold at 28\,GHz ($T_0=0.536$).
  The ICNIRP reference level for incident power density is 10\,W/m$^2$.
  Columns are shown for both the theoretical worst-case ($D=2.0$) and a
  more pragmatic bound ($D=1.5$).}\label{tab:compliance}
\begin{tabular}{lcccccc}
\toprule
Body type & $m$\,(kg) & $A$\,(m$^2$) & $m/A$ & $\Sinc^{\max, 2.0}$ & $\Sinc^{\max, 1.5}$ & vs.\ ICNIRP \\
\midrule
Adult     & 70  & 1.8 & 38.9 & 11.5 & 15.3 & compliant \\
Child     & 20  & 0.8 & 25.0 & 7.4  & \phantom{0}9.9  & non-compliant \\
Infant    &  8  & 0.4 & 20.0 & 5.9  & \phantom{0}7.9  & non-compliant \\
\bottomrule
\end{tabular}
\end{table}

The use of $D=1.5$ as a pragmatic bound increases the compliance threshold, but children and infants remain non-compliant at the 10\,W/m$^2$ level.
From ICNIRP's perspective, this indicates that the reference levels are conservative for adults but may be less so for smaller individuals under worst-case directional exposure.

Adults automatically comply with the whole-body SAR limit at the
ICNIRP reference level of 10\,W/m$^2$. This can be verified
analytically. Children and infants are more vulnerable because $m/A$
scales roughly as $m^{1/3}$ (mass grows faster with body size than
surface area).

\subsection{Anthropometric scaling}\label{sec:anthropometric}

Using the Du~Bois formula~\cite{DuBois1916} for body surface area,
\begin{equation}\label{eq:DuBois}
  A \approx 0.007184\times m^{0.425}\times h^{0.725}
  \quad\text{($m$ in kg, $h$ in cm)},
\end{equation}
the critical ratio $m/A$ scales as
$\mathrm{BMI}^{0.575}\cdot h^{0.425}$.  The SAR compliance threshold
increases with both body-mass index and height. Taller, heavier
individuals can tolerate more incident power before the SAR limit is
reached.

\begin{table}[!htbp]
\centering
\caption{Compliance threshold across the human population at 28\,GHz.
  The threshold varies by a factor of roughly~2 between an infant
  and a large adult.}\label{tab:anthropometric}
\begin{tabular}{lcccccc}
\toprule
Body type & $m$\,(kg) & $h$\,(cm) & BMI & $m/A$\,(kg/m$^2$)
  & $\Sinc^{\max}$\,(W/m$^2$) \\
\midrule
Infant (1\,yr) &  8 &  70 & 16.3 & 21.1 & 6.2  \\
Child (6\,yr)  & 20 & 110 & 16.5 & 25.8 & 7.6  \\
Adolescent     & 50 & 160 & 19.5 & 33.5 & 9.9  \\
Adult (ref.)   & 70 & 170 & 24.2 & 38.7 & 11.4 \\
Large adult    &100 & 180 & 30.9 & 45.6 & 13.5 \\
\bottomrule
\end{tabular}
\end{table}

Infants and small children are about 45\% more vulnerable than
reference adults.

\begin{remark}[Spherical phantoms do not bound human exposure]
For a given volume, the sphere minimises surface area and has
$D(\khat)=1$ uniformly.  A spherical phantom therefore absorbs the
least average power of any convex body of the same volume.  The human
body's elongated shape implies $D_{\max}>1$, so spherical phantoms
cannot provide an upper bound on worst-case human exposure; the actual
$D(\khat)$ is required.
\end{remark}


% ============================================================
\section{Computational framework}\label{sec:computation}
% ============================================================

\subsection{Multi-source environments embedded in neural networks}\label{sec:neural-net}

In a realistic wireless environment, the body is illuminated by $N$
plane waves (from scattering, reflections, multiple antennas), each
with power density~$S_i$ and direction~$\khat_i$.  For a convex body
(or using occlusion-corrected normals), the total absorbed power
density at a surface point~$\rr$ is
\begin{equation}\label{eq:Sab-multi}
  \Sab(\rr) = T_0\sum_{i=1}^{N}S_i\;\ReLU\!\bigl(\nhat(\rr)\cdot(-\khat_i)\bigr).
\end{equation}

This is a single-hidden-layer ReLU neural network~\cite{NairHinton2010,Goodfellow2016} with $N$ hidden
units, evaluated at the surface normal:
\begin{equation}\label{eq:nn}
  \boxed{\Sab(\rr) = \mathbf{a}^\top\ReLU(\mathbf{W}\,\nhat(\rr)),}
\end{equation}
where:
\begin{itemize}[nosep]
\item $\mathbf{W}=[-\khat_1;\ldots;-\khat_N]\in\mathbb{R}^{N\times 3}$
  is the weight matrix (path directions),
\item $\mathbf{a}=T_0\,[S_1,\ldots,S_N]^\top\in\mathbb{R}^N$ is the
  output-weight vector (path powers $\times$ tissue transmission),
\item the bias is zero,
\item the activation is $\ReLU$.
\end{itemize}

The electromagnetic environment, as computed by a ray tracer, defines
the ``weights'' of this network; the body ``evaluates'' it at each
surface point.

For non-convex bodies, each hidden unit is additionally gated by the
binary occlusion factor~$O(\rr,\khat_i)$:
\begin{equation}\label{eq:gated-relu}
  \Sab(\rr) = T_0\sum_{i=1}^{N}S_i\;
  \ReLU\!\bigl(\nhat(\rr)\cdot(-\khat_i)\bigr)\;O(\rr,\khat_i).
\end{equation}
This is a gated ReLU network in which the gates encode body
self-shadowing.

\subsection{Matrix formulation}\label{sec:matrix}

In discrete form, let the body mesh have $M$ triangles with
normal matrix $\mathbf{N}\in\mathbb{R}^{M\times 3}$, area
vector~$\mathbf{a}\in\mathbb{R}^M$, and let the ray tracer provide
$N$ paths with direction matrix
$\mathbf{K}=[-\khat_1;\ldots;-\khat_N]\in\mathbb{R}^{N\times 3}$ and
power vector~$\mathbf{s}\in\mathbb{R}^N$.  Define the positive-part
cosine matrix
$\bm{\mu}_+=\ReLU(\mathbf{N}\mathbf{K}^\top)\in\mathbb{R}^{M\times N}$
and the binary visibility matrix
$\mathbf{O}\in\{0,1\}^{M\times N}$.

The per-triangle $\Sab$ vector is
\begin{equation}\label{eq:Sab-mat}
  \mathbf{S}_{\mathrm{ab}} = T_0\;(\bm{\mu}_+\odot\mathbf{O})\;\mathbf{s}
  \;\in\mathbb{R}^M,
\end{equation}
where $\odot$ denotes the Hadamard (element-wise) product.  For convex
bodies, $\mathbf{O}=\mathbf{1}$ and the Hadamard product drops out.

To compute the spatially averaged $\Sab$ (averaging over a
$4\,\text{cm}^2$ area as specified by ICNIRP), let $\mathbf{G}\in\mathbb{R}^{M\times M}$ be a sparse,
row-stochastic averaging matrix encoding the neighbourhood structure.
Then $\bar{\mathbf{S}}_{\mathrm{ab}}=\mathbf{G}\,\mathbf{S}_{\mathrm{ab}}$.

The total absorbed power is $P_{\mathrm{abs}}=\mathbf{a}^\top\mathbf{S}_{\mathrm{ab}}$.

The corresponding whole-body SAR is $\SARwb = P_{\mathrm{abs}}/m$.

\subsection{Hierarchical computation}\label{sec:hierarchy}

The full per-triangle computation is $O(MN)$. It is necessary when the
spatial distribution of $\Sab$ is needed, but it is wasteful when only
aggregate quantities are required.  The projected-area identity
\cref{eq:Pabs-Aperp} provides a much faster path: total absorbed
power factors through the shadow area,
\begin{equation}\label{eq:Pabs-lookup}
  P_{\mathrm{abs}} = T_0\sum_{i=1}^{N}S_i\;\Aperp(\khat_i),
\end{equation}
where $\Aperp(\khat_i)$ is looked up from a precomputed table on
$S^2$.  In terms of the directivity:
\begin{equation}\label{eq:Pabs-D-multi}
  \boxed{P_{\mathrm{abs}} = T_0\,\frac{\Aab}{4}
    \sum_{i=1}^{N}S_i\;D(\khat_i).}
\end{equation}

This is $O(N)$. It is a weighted sum over arriving paths, with weights from
a precomputed table.  The body mesh ($M$ triangles) does not appear in
the inner loop.

These observations yield a natural three-level hierarchy:

\begin{table}[!htbp]
\centering
\caption{Computation hierarchy.  All levels are fed by the same
  ray-tracing output.}\label{tab:hierarchy}
\begin{tabular}{llll}
\toprule
Level & Output & Formula & Cost \\
\midrule
Bound     & Worst-case $P_{\mathrm{abs}}$
  & $\le T_0\,\frac{\Aab\,D_{\max}}{4}\sum_i S_i$
  & $O(1)$ \\[4pt]
Aggregate & Exact $P_{\mathrm{abs}}$, $\SARwb$
  & $T_0\frac{\Aab}{4}\sum_i S_i\,D(\khat_i)$
  & $O(N)$ \\[4pt]
Spatial   & Full $\Sab(\rr)$ map
  & $T_0\,(\bm{\mu}_+\odot\mathbf{O})\,\mathbf{s}$
  & $O(MN)$ \\
\bottomrule
\end{tabular}
\end{table}

For gradient-based optimisation, the aggregate level suffices in the
inner loop whenever the binding constraint is whole-body SAR.  With
$M\sim 10^4$ and $N\sim 10^2$, it is roughly $10^4$ times faster than
the spatial computation.  The spatial level is invoked only for the
final compliance check or when the peak spatially-averaged $\Sab$
limit is binding.

\subsection{Precomputed body response}\label{sec:precomputed}

The function $\Aperp(\khat)$ (equivalently $D(\khat)$) is the
body's angular transfer function, which maps the angular power spectrum
of the incident field to total absorbed power.  For a given
body mesh and pose, it is computed once during the ambient-occlusion
pass and stored as a look-up table on $S^2$ (for instance, on a
Fibonacci lattice of ${\sim}\,10^3$ directions).  At run time, the
value at an arbitrary path direction is obtained by interpolation.

Together with the spherical-harmonic representation
(\cref{sec:directivity}), this reduces the aggregate-level computation
to a dot product of coefficient vectors. It is essentially free compared to
the ray tracing itself.

\begin{table}[!htbp]
\centering
\caption{Precomputed quantities from the body mesh.  All are
  independent of the electromagnetic environment.}
\label{tab:precomputed}
\begin{tabular}{llll}
\toprule
Quantity & Symbol & Size & Description \\
\midrule
Surface normals & $\mathbf{N}$ & $M\times 3$ & Triangle normals \\
Triangle areas & $\mathbf{a}$ & $M$ & Triangle areas \\
Averaging matrix & $\mathbf{G}$ & $M\times M$ (sparse)
  & Spatial averaging for ICNIRP \\
Exposure fraction & $\bm{\eta}$ & $M$
  & Ambient occlusion per triangle \\
Absorption area & $\Aab$ & scalar & $=\mathbf{a}^\top\bm{\eta}$ \\
Projected-area table & $\Aperp(\khat)$ & $N_{\mathrm{dir}}$
  & Shadow area vs.\ direction \\
Absorption directivity & $D(\khat)$ & $N_{\mathrm{dir}}$
  & $=\Aperp/(\Aab/4)$, shape anisotropy \\
\bottomrule
\end{tabular}
\end{table}

\subsection{Diffuse-environment limit}\label{sec:diffuse-limit}

In richly scattered indoor environments where multipath creates
approximately isotropic illumination~\cite{Andersen2007}, the
direction-averaged results of \cref{sec:cauchy} apply directly:
\begin{equation}\label{eq:diffuse}
  \Sab(\rr) \approx \frac{S_{\mathrm{total}}\,T_0}{4}\;\eta(\rr),
  \qquad
  P_{\mathrm{abs}} \approx
  \frac{S_{\mathrm{total}}\,T_0\,\Aab}{4},
\end{equation}
where $S_{\mathrm{total}}=\sum_i S_i$.  The precomputed exposure
fraction~$\eta(\rr)$ \emph{is} the dosimetry map for isotropic
illumination.  No per-path computation is required.

In the isotropic limit, the generalised Cauchy formula reduces to
$\langle P_{\mathrm{abs}}\rangle = S_{\mathrm{total}}\,T_0\,\Aab/4$.
This result matches the BSA-based formula
$P_{\mathrm{abs}} = S_{\mathrm{total}}\,\eta\,\mathrm{BSA}/4$
used in the room-electromagnetics literature.
Andersen~\textit{et al.}~\cite{Andersen2007} introduced the diffuse-field
framework; Bamba~\textit{et al.}~\cite{Bamba2014} extracted
an absorption efficiency $\eta \approx 0.5$ from FDTD simulations at
1.5--6\,GHz, where $\mathrm{BSA} = A$ is the body surface area.
The framework identifies $\eta$ with $T_0$ (\cref{sec:pseudo-brewster}).
The factor $\Aab/A$ accounts for body non-convexity ($\Aab = \gamma_s A$
in the notation of Flintoft~\textit{et al.}~\cite{Flintoft2014}).
The statistical multipath exposure tool of
Vermeeren~\textit{et al.}~\cite{Vermeeren2008} provides a complementary
approach based on Monte Carlo sampling of the angular spectrum.

Applications of this limit are developed in \cref{sec:applications}.

\begin{remark}[Polarisation-aware extension]
For linearly polarised sources, the neural-network architecture of
\cref{eq:nn} extends to a Stokes inner product: the scalar look-up
$T_0\Aperp(\khat)$ is replaced by a four-component absorption Stokes
vector~$\mathbf{m}(\khat)$ (\cref{sec:polarisation}).  The cost is
one additional dot product per incidence direction.
\end{remark}


% -------------------------------------------------------
\subsection{Near-field and point-source extension}\label{sec:near-field}

The preceding sections assume far-field illumination: constant
$\Sinc$ and~$\khat$ across the body.  This subsection extends the
framework to point sources at arbitrary distance.  Three distance
regimes govern the validity of the exact absorption law.

\begin{definition}[Reactive near-field regime]
($d < \lambda/(2\pi)$).  In this regime the field is dominated by
non-propagating (evanescent) components.  The exact absorption law
\cref{eq:Sab-exact} is not applicable.  Dosimetry requires the full
near-field coupling logic of \cref{sec:setup}.

At 28\,GHz this boundary lies at 1.7\,mm from the
source; at 60\,GHz it shrinks to 0.8\,mm.  Only sources in direct
contact with the skin fall in this regime.
\end{definition}

\begin{definition}[Locally plane-wave regime]
($d > 3\lambda$).  The reactive field has decayed to negligible
levels, and the wavefront curvature is small enough that the
illumination is locally plane but non-uniform in power.  The
geometric absorption law \cref{eq:geometric-law} holds pointwise, provided $\Sinc$
is varied with distance and direction.  The total absorbed power
generalises to a solid-angle integral (\cref{eq:Pabs-solidangle}).

For a point source at distance $d > 3\lambda$ (3.2\,cm at 28\,GHz), the
exact absorption law applies at each surface point with spatially
varying inputs:
\begin{equation}\label{eq:Sab-point}
  \Sab(\rr)
  = \frac{P_t\,G\!\bigl(\hat{k}(\rr)\bigr)}{4\pi\,d(\rr)^2}
    \;\cdot\; T_0\;\cdot\;
    \ReLU\!\bigl[\nhat(\rr)\cdot\bigl(-\hat{k}(\rr)\bigr)\bigr],
\end{equation}
where $P_t$ is the total radiated power, $G(\hat{u})$ is the antenna
gain pattern, $d(\rr)=|\rr-\rr_s|$ is the source distance, and
$\hat{k}(\rr)=(\rr-\rr_s)/d(\rr)$.  The pseudo-Brewster
compensation and the ReLU cosine gate carry over without
modification.  The computational cost is $O(M)$ per source, the
same as the spatial computation for plane waves.
\end{definition}

\begin{definition}[Far-field regime]
($d > 5L$, where $L$ is the largest body dimension).  The wavefront
curvature is negligible over the entire body.  Distance variation of
$\Sinc$ is negligible.  Whole-body absorption is described by a single
incidence direction~$\khat$ and a constant power density~$\Sinc$.

When the source is many body lengths away, $d(\rr)\approx d_0$ and
$\hat{k}(\rr)\approx\hat{k}_0$ for all surface points.  The
plane-wave geometric absorption law applies with
$\Sinc = P_t G(\hat{k}_0)/(4\pi d_0^2)$.  The approximation error
scales as $(L/(2d))^2$, where $L$ is the transverse body extent in
the approach direction.  For a standing human ($L\approx 1.2$\,m)
illuminated from the front, the error is below 1\% at $d>4$\,m and
below 3\% at $d\approx 2$\,m.
\end{definition}

The scaling is direction-dependent: front and side illumination
produce a conservative (underestimating) error because the body's
elongated shape places most surface elements farther than the reference
distance $d_0$.  Overhead illumination can produce a non-conservative
error because the body centre (near the waist) is far below the
top of the head, causing the reference distance to overestimate the
actual source-to-surface distance.

\paragraph{Point-source total power and the solid-angle formula.}
Integrating \cref{eq:Sab-point} over the illuminated surface, the
factor $\mu(\rr)\,\diff A/d(\rr)^2$ is the differential solid angle
$d\Omega_s$ subtended by the surface element as seen from the source.
For an isotropic source ($G\equiv 1$):
\begin{equation}\label{eq:Pabs-solidangle}
  \boxed{P_{\mathrm{abs}}
    = \frac{P_t\,T_0}{4\pi}\;\Omega_{\mathrm{body}}(\rr_s),}
\end{equation}
where $\Omega_{\mathrm{body}}(\rr_s)
= \int_{\Sigma_+}\mu\,\diff A/d^2$ is the solid angle subtended by the
visible, front-facing body surface as seen from the source.
This formula connects the dosimetry framework to the view-factor
formalism of radiative heat transfer: absorbed power equals $T_0$
times the view factor times the radiated power.

In the far field, $\Omega_{\mathrm{body}}\to\Aperp(\hat{k}_0)/d_0^2$
and \cref{eq:Pabs-solidangle} recovers
$P_{\mathrm{abs}}=\Sinc\,T_0\,\Aperp$.  On the body surface
($\rr_s\to\Sigma$), $\Omega_{\mathrm{body}}\to 2\pi$ and
$P_{\mathrm{abs}}\to P_t\,T_0/2$.

For a directive antenna, the solid-angle integral acquires the gain
pattern as a weight:
\begin{equation}\label{eq:Pabs-directive}
  P_{\mathrm{abs}}
  = \frac{P_t\,T_0}{4\pi}\int_{\Omega_{\mathrm{body}}}
    G(\hat{u})\,d\Omega_s.
\end{equation}

\begin{remark}[Antenna--body decoupling via spherical harmonics]
\label{rem:SH-NF}
Expanding the antenna gain in spherical harmonics,
$G(\hat{u})=\sum_{lm}g_{lm}\,Y_{lm}(\hat{u})$, and exchanging the
sum with the surface integral in \cref{eq:Pabs-directive} gives
\begin{equation}\label{eq:Gamma-lm}
  P_{\mathrm{abs}}
  = \frac{T_0}{4\pi}\sum_{lm}g_{lm}\,\Gamma_{lm}(\rr_s),
  \qquad
  \Gamma_{lm}(\rr_s)
  \equiv \int_{\Sigma_+}
    \frac{\mu(\rr)\,Y_{lm}\!\bigl(\hat{k}(\rr)\bigr)}{d(\rr)^2}
    \;O(\rr,\hat{k})\;dA.
\end{equation}
The coefficients~$\Gamma_{lm}(\rr_s)$ encode the body's near-field
response to each spherical-harmonic antenna mode.  They depend on the
source position but not on the antenna pattern.  For a given
position~$\rr_s$, any antenna is evaluated via a dot product in
$O(L^2)$ operations.

For phone-scale dosimetry ($L\le 6$, 1\,cm resolution, 50\,cm extent),
the look-up table requires approximately 50\,MB of storage.  In the
far field, $\Gamma_{lm}\to\Aperp(\hat{k}_0)\,Y_{lm}(\hat{k}_0)/d_0^2$,
recovering the directivity--based framework of
\cref{sec:precomputed}.  The decoupling is a computational
acceleration for close-range dosimetry (e.g.\ testing many antenna
orientations at a fixed phone position); it does not introduce new
physics.
\end{remark}


% ============================================================
\section{Higher-order corrections}\label{sec:corrections}
% ============================================================

The geometric absorption law treats the surface as locally flat.  Two
systematic corrections refine the formula near curved regions and at
the shadow boundary.

\subsection{Curvature correction}\label{sec:curvature}

For a surface with principal radii of curvature $R_1$, $R_2$, the
first-order Physical Optics correction is
\begin{equation}\label{eq:curvature}
  \Sab(\rr) = \Sinc\,T_0\,\ReLU(\mu)
  \left[1 + \frac{\mu}{kR_1}+\frac{\mu}{kR_2}\right],
\end{equation}
where $k=2\pi/\lambda$ is the free-space wavenumber.  Since
$\ReLU(\mu)\cdot\mu=\ReLU(\mu)^2$, this separates cleanly into a
standard ReLU term plus a quadratic correction:
\begin{equation}\label{eq:curvature-nn}
  \boxed{S_{\mathrm{ab},j} = T_0\sum_i S_i\!\left[
    \ReLU(\mu_{ji}) + \frac{H_j}{k}\,\ReLU(\mu_{ji})^2\right],}
\end{equation}
where $H_j=1/R_{1,j}+1/R_{2,j}$ is twice the mean curvature at
triangle~$j$.

\begin{table}[!htbp]
\centering
\caption{Curvature correction magnitude at 28\,GHz
  ($k\approx 587$\,m$^{-1}$).  The correction is below the Fresnel
  approximation error for most body regions.}\label{tab:curvature}
\begin{tabular}{lccc}
\toprule
Body region & Radius $R$ & $1/(kR)$ & Correction \\
\midrule
Torso / head  & $>100$\,mm & ${<}\,0.2\%$ & negligible \\
Arm           & ${\sim}\,40$\,mm & 0.4\% & ${\sim}\,0.4\%$ \\
Finger        & ${\sim}\,8$\,mm & 2.1\% & ${\sim}\,2\%$ \\
Ear edge      & ${\sim}\,2$\,mm & 8.5\% & ${\sim}\,8\%$ \\
\bottomrule
\end{tabular}
\end{table}

\subsection{Diffraction smoothing}\label{sec:diffraction}

At the shadow boundary ($\mu=0$), diffraction smooths the sharp ReLU
discontinuity over a Fresnel-zone width~$\sqrt{\lambda R}$.  The
physical activation function becomes
\begin{equation}\label{eq:gelu}
  \ReLU_{\mathrm{phys}}(\mu) \approx \mu\;\tfrac{1}{2}\!\left[
    1 + \erf\!\left(\frac{\mu}{\sigma_j}\right)\right],
  \qquad \sigma_j = \sqrt{\frac{\lambda}{2\pi R_j}},
\end{equation}
where $R_j$ is the local radius of curvature at the shadow boundary.
This is the GELU (Gaussian Error Linear Unit) from the machine-learning
literature, but with a smoothing scale set by physics rather than
learned.

\begin{table}[!htbp]
\centering
\caption{Correspondence between asymptotic optics regimes and
  neural-network activations.}\label{tab:activations}
\begin{tabular}{lll}
\toprule
Activation & Formula & Physical origin \\
\midrule
ReLU  & $\max(0,\mu)$ & Geometric optics ($\lambda\to 0$) \\
GELU  & $\mu\,\Phi(\mu/\sigma)$ &
  Diffraction at a curved shadow boundary \\
\bottomrule
\end{tabular}
\end{table}

For compliance purposes the sharp ReLU is sufficient: ICNIRP's
4\,cm$^2$ spatial averaging already regularises the shadow boundary.
The GELU version is relevant for high-resolution local dosimetry or
comparison with point measurements near the terminator line.

\subsection{Hierarchy of approximations}\label{sec:hierarchy-approx}

\begin{table}[!htbp]
\centering
\caption{Three levels of fidelity, all differentiable.  The basic
  level suffices for most compliance work.}\label{tab:fidelity}
\begin{tabular}{llll}
\toprule
Level & Formula & Extra data & Accuracy \\
\midrule
Basic       & $T_0\,\ReLU(\mathbf{M})\,\mathbf{s}$
  & Normals only & ${\sim}\,5\%$ \\
$+$ Curvature & $+\,(H_j/k)\,\ReLU(\mathbf{M})^{\circ 2}\,\mathbf{s}$
  & $+$ mean curvature & ${\sim}\,1\%$ (curved) \\
$+$ Diffraction & ReLU $\to$ GELU$_{\sigma_j}$
  & $+$ local radius & Smooth boundary \\
\bottomrule
\end{tabular}
\end{table}

The hierarchy from geometric optics to physical optics and diffraction
maps onto activation functions: ReLU to ReLU + ReLU$^2$ and finally GELU.
Higher-order corrections would add further polynomial powers of $\ReLU(\mu)$,
forming a Luneburg--Kline asymptotic expansion.


% ============================================================
\section{Validation}\label{sec:validation}
% ============================================================

\subsection{Anatomical phantom}\label{sec:phantom}

The geometric absorption law was evaluated on the Thelonious phantom
(23\,826~triangles, approximately 7\,870\,cm$^2$ surface area)
illuminated by a plane wave from above, and compared against the full
angle-dependent Fresnel integration. To clarify, the comparison is between:
\begin{align}
  \Sab^{\mathrm{simplified}} &= \Sinc\;T_0\;\ReLU(\mu) \label{eq:val-simp} \\
  \Sab^{\mathrm{full}} &= \Sinc\;\Teff(\theta, \text{pol})\;\ReLU(\mu) \label{eq:val-full}
\end{align}
where $\Teff$ in \cref{eq:val-full} includes the full angular dependence $\Tavg(\theta)$ and the polarization-splitting correction $(q/2)\,\Delta T(\theta)$ from \cref{eq:Teff-decomp}.
The results are summarized in \cref{tab:phantom}.

\begin{table}[!htbp]
\centering
\caption{Framework vs.\ full Fresnel integration on Thelonious
  (skin at 28\,GHz, plane wave from above).}\label{tab:phantom}
\begin{tabular}{lccc}
\toprule
Metric & Simplified & Full Fresnel & Error \\
\midrule
Total absorbed power  & 48.12\,mW & 47.95\,mW & 0.35\% \\
Mean $\Sab$ (illuminated) & 0.185\,W/m$^2$ & 0.186\,W/m$^2$ & 0.5\% \\
Peak $\Sab$           & 0.539\,W/m$^2$ & 0.539\,W/m$^2$ & 0.0\% \\
\bottomrule
\end{tabular}
\end{table}

For the 4\,907 triangles with $\theta<75^\circ$, the local error
statistics are: mean $-2.6\%$, RMS $3.2\%$, range
$[-5.3\%,\;0.0\%]$.  The simplified formula slightly underestimates
$\Sab$ because $\Tavg(\theta)>T_0$ for $\theta>0$ in this frequency
range. This is a conservative error.

\Cref{fig:apd-direction} extends this analysis across all illumination
directions and all three polarisation states.  The theoretical
prediction $P_{\mathrm{abs}} = T_0\,\Aperp$ (dashed line at
105.9\,mW) falls within the interquartile range of the unpolarised
distribution.  TM polarisation produces systematically higher total
absorbed power than TE, consistent with the pseudo-Brewster peak
in~$T_p$.  The directional coefficient of variation is 19\% for
unpolarised illumination, 14\% for TM, and 28\% for TE.  The lower
variability of TM reflects its flatter angular transmission profile
over the relevant incidence angles on the body.

\begin{figure}[!htbp]
  \centering
  \includegraphics{apd_direction_analysis.pdf}
  \caption{Directional and polarisation analysis on the Thelonious
    phantom (skin at 28\,GHz, $\Sinc = 1$\,W/m$^2$).
    The upper-left panel shows a histogram of total absorbed power across 128
    illumination directions for TE (red), TM (blue), and unpolarised (green).
    The upper-right panel shows the probability density of local $\Sab$ on
    illuminated triangles.  The lower-left panel shows box plots by
    polarisation; the dashed line marks the theoretical prediction
    $T_0\,\Aperp = 105.9$\,mW.  The lower-right panel shows the coefficient of
    variation of total absorbed power and peak $\Sab$ across directions.}
  \label{fig:apd-direction}
\end{figure}

\subsection{Mie theory: exact wave-theoretic comparison}
\label{sec:mie}

The phantom validation tests the Fresnel approximation
($\Tavg\approx T_0$) but not the geometric-optics assumption itself.
To test the full framework including diffraction effects, we compare
against exact Mie theory for lossy dielectric spheres.

\subsubsection{Method}

For a lossy sphere of radius~$a$ and complex refractive index~$\ntilde$,
Mie theory~\cite{BohrenHuffman1983} provides an exact series solution
for the absorption efficiency~$Q_{\mathrm{abs}}$, which is the ratio of the
absorption cross section to the geometric cross section $\pi a^2$:
$P_{\mathrm{abs}}=\Sinc\,Q_{\mathrm{abs}}\,\pi a^2$.
Our framework predicts $P_{\mathrm{abs}}=\Sinc\,T_0\,\pi a^2$, so the
framework error is $(T_0/Q_{\mathrm{abs}}-1)\times 100\%$.

\begin{remark}[Absorption cross section]
The quantity $\sigma_a \equiv Q_{\mathrm{abs}}\,\pi a^2$ is the
\emph{absorption cross section} (ACS) of the sphere.
For a human body, the direction-averaged ACS~$\langle\sigma_a\rangle$
is the standard observable in reverberation chamber dosimetry
(see \cref{sec:app-reverb} and \cref{sec:Tbar}).
The generalised Cauchy formula (\cref{sec:cauchy}) gives
$\langle\sigma_a\rangle = T_0\,\Aab/4$, where
$\Aab = \int_\Sigma\eta(\rr)\,\diff A$ is the effective absorption area.
\end{remark}

Skin dielectric properties at each frequency are obtained from the
IT'IS~v5.0 database using the 4-Cole-Cole model; this guarantees
physical consistency of $\ntilde(f)$ and $T_0(f)$.

\subsubsection{Error decomposition}

The total framework error consists of two independent components.

\paragraph{Fresnel approximation error.}
This is quantified by the sphere ratio~$R$ defined in
\cref{eq:R-sphere}.  It is shape-dependent, frequency-dependent, and
independent of size.  The sphere gives the worst-case Fresnel error
among smooth convex shapes because it maximises the angular spread of
incidence angles.

\paragraph{Diffraction error.}
Diffraction bends waves around the sphere into the geometric shadow,
causing additional absorption that the framework misses.  This
contribution decreases with increasing size parameter
$x=\pi d/\lambda$ (Fock's theory predicts scaling as~$x^{-2/3}$).

\subsubsection{Frequency-dependent asymptotic limit}

The asymptotic error as $x\to\infty$ is not a universal constant but
varies with frequency because the skin's refractive index is
dispersive.  At ${\sim}\,39$\,GHz,
$R_{\mathrm{sphere}}=1$ exactly. This is a coincidence of the dispersion
curve.  Below this frequency the framework underestimates
($R<1$); above it, the framework overestimates ($R>1$).

\begin{table}[!htbp]
\centering
\caption{Asymptotic ($x\to\infty$) framework error for skin at
  various frequencies.  The Fresnel error changes sign near
  40\,GHz.}\label{tab:mie-asymptotic}
\begin{tabular}{lcccc}
\toprule
Frequency & $|\ntilde|$ & $T_0$ & $R_{\mathrm{sphere}}$
  & Asymptotic error \\
\midrule
6\,GHz   & 6.1 & 0.481 & 0.968 & $-3.2\%$  \\
10\,GHz  & 5.9 & 0.489 & 0.970 & $-3.0\%$  \\
28\,GHz  & 4.8 & 0.536 & 0.988 & $-1.2\%$  \\
40\,GHz  & 4.3 & 0.570 & 1.000 & $\phantom{+}0.0\%$ \\
60\,GHz  & 3.7 & 0.622 & 1.015 & $+1.5\%$  \\
100\,GHz & 3.0 & 0.701 & 1.035 & $+3.5\%$  \\
\bottomrule
\end{tabular}
\end{table}

\Cref{fig:mie} collects the Mie validation results.  The upper
panel shows the framework error as a function of size parameter
$x=\pi d/\lambda$ at 28\,GHz: the error converges from below towards
the Fresnel limit ($R_{\mathrm{sphere}}-1\approx -1.2\%$) as
$x\to\infty$.  Body-relevant sizes ($x\gtrsim 10$) fall in the
$-20\%$ to $-5\%$ range.  The middle panel shows the same error as a
function of frequency for four representative body-part diameters,
with the mmWave band (20 to 60\,GHz) shaded.  The lower panel
shows the Fresnel limit $R_{\mathrm{sphere}}(f)$, which crosses unity
near 39\,GHz.

\begin{figure}[!htbp]
  \centering
  \includegraphics{mie_validation_corrected.pdf}
  \caption{Mie-theory validation of the framework against lossy
    dielectric spheres with frequency-dependent IT'IS skin properties.
    The upper panel shows framework error versus size parameter at 28\,GHz;
    vertical dashed lines mark body-part sizes, and the error converges to the
    Fresnel limit (orange) as $x$ increases.  The middle panel shows framework
    error versus frequency for finger (17\,mm), arm (80\,mm), head (180\,mm),
    and torso (300\,mm) diameters; the mmWave band is shaded green.  The lower
    panel shows the frequency-dependent Fresnel limit
    $R_{\mathrm{sphere}}=T_0/\langle\Tavg\rangle$, crossing $R=1$ at
    ${\sim}\,39$\,GHz.}
  \label{fig:mie}
\end{figure}

\subsubsection{Error versus body-part size}

\begin{table}[!htbp]
\centering
\caption{Framework error versus frequency for representative body
  parts, computed from Mie theory.  Positive errors indicate
  overestimation.}\label{tab:mie-bodyparts}
\begin{tabular}{lcccc}
\toprule
Frequency & Finger (17\,mm) & Arm (80\,mm)
  & Head (180\,mm) & Torso (300\,mm) \\
\midrule
6\,GHz   & $-60.7\%$ & $-36.8\%$ & $-26.5\%$ & $-21.4\%$ \\
28\,GHz  & $-38.2\%$ & $-19.8\%$ & $-13.5\%$ & $-10.4\%$ \\
60\,GHz  & $-27.4\%$ & $-11.3\%$ & $-6.1\%$  & $-3.8\%$ \\
100\,GHz & $-19.1\%$ & $-5.1\%$  & $-1.2\%$  & $+0.4\%$ \\
\bottomrule
\end{tabular}
\end{table}

For body-scale objects (torso, head) in the core mmWave band
(28 to 60\,GHz), the total framework error is 3 to 14\%, dominated by
diffraction into the geometric shadow.  Fingers at low frequencies are
the worst case: their small size parameter admits significant
wave-resonance effects.

\subsubsection{Effect on the geometric bounds}

The geometric identities of \cref{sec:geometry} (Cauchy's formula,
hemisphere limit, $D\le 2$) are exact. They constrain the projected
area, which is purely geometric.  The Mie analysis shows that the
physics approximation $P_{\mathrm{abs}}=\Sinc\,T_0\,\Aperp$
introduces a systematic error, but this does not invalidate the
bounds.
First, the hemisphere bound
$P_{\mathrm{abs}}\le\Sinc\,T_0\,A/2$ holds for the
true Mie power at all body-relevant sizes: numerically,
$Q_{\mathrm{abs}}<2T_0$ for every combination of body-part
diameter and frequency tested ($x>5$).
Second, for the peak $\Sab$, the maximum occurs at normal incidence, where
the Fresnel approximation is exact ($\Teff(0)=T_0$ for any
polarisation), so the framework reproduces the peak faithfully.
Near the shadow boundary, diffraction smooths the sharp ReLU
transition (\cref{sec:diffraction}), causing the framework to
slightly overestimate the lit side and underestimate the shadow
side; the net effect is conservative for local exposure limits.
Finally, for total power or $\SARwb$, the framework underestimates by
${\sim}\,10\%$ at mmWave, because it misses shadow-region
absorption.  This is non-conservative by a small margin, negligible
relative to ICNIRP's 50$\times$ safety factor.

The sphere provides a conservative bound on framework accuracy: its
uniform curvature maximises both the Fresnel and diffraction errors.
A human body, with large flat regions on the torso, has smaller net
errors.

\subsection{Error budget}\label{sec:error-budget}

\begin{table}[!htbp]
\centering
\caption{Error budget for the framework at mmWave
  frequencies.}\label{tab:error-budget}
\begin{tabular}{llp{4.5cm}}
\toprule
Source & Magnitude & Notes \\
\midrule
Fresnel ($\Tavg\approx T_0$) &
  ${\sim}\,5\%$ local, ${\sim}\,1\%$ integrated &
  Worst over $\theta<75^\circ$ \\
Diffraction (shadow absorption) &
  ${\sim}\,10\%$ for torso at 28\,GHz &
  Sphere worst case \\
Curvature corrections &
  ${\sim}\,0.5\%$ & Most body regions \\
Near-field effects &
  ${\sim}\,1\%$ & Sources $>1\lambda$ away \\
Tissue-parameter uncertainty &
  ${\sim}\,20\%$ & IT'IS database variability \\
\midrule
Total (framework vs.\ exact EM) &
  $\sim$10 to 15\% & Sphere worst case \\
Total (incl.\ tissue uncertainty) &
  $\sim$15 to 20\% & Dominated by tissue $+$ diffraction \\
\bottomrule
\end{tabular}
\end{table}

\Cref{fig:error-budget} shows the error budget for two
scenarios: a best case (28\,GHz, torso, multipath) and a worst case
(6\,GHz, finger, single-path line-of-sight).  Each bar is classified
as conservative (underestimate, blue), non-conservative
(overestimate, red), or symmetric uncertainty (grey).

\begin{figure}[!htbp]
  \centering
  \includegraphics[width=\textwidth]{error_budget_comprehensive.pdf}
  \caption{Error budget for two representative scenarios.
    Blue bars are conservative errors (framework underestimates
    absorption); red bars are non-conservative; grey bars are
    symmetric uncertainties.  The dominant source is tissue
    dielectric uncertainty in both cases.}
  \label{fig:error-budget}
\end{figure}


% ============================================================
\section{Applications}\label{sec:applications}
% ============================================================

The framework reduces dosimetry to precomputed geometric quantities
and a single material constant.  Three application domains benefit
directly.


\subsection{Exposure-aware wireless network design}
\label{sec:app-network}

Modern GPU-accelerated ray-tracing engines such as
Sionna~RT~\cite{SionnaRT} provide differentiable propagation
simulation (see also~\cite{Shikhantsov2020} for dosimetric
coupling).  A typical pipeline proceeds in four steps.
First, the scene geometry is loaded and material properties are assigned.
Second, rays are traced from transmitters through reflections, diffractions,
and scattering events.
Third, at each receiver location, the $N$ arriving paths are collected with
directions~$\khat_i$, powers~$S_i$, and delays~$\tau_i$.
Finally, these are combined with antenna patterns to form the channel response.

Step~3 already provides all the plane waves arriving at the body.
The dosimetry formula attaches as a single differentiable
post-processing layer:
\begin{equation}\label{eq:diff-rt}
  \mathbf{S}_{\mathrm{ab}}
  = T_0\;\ReLU(\mathbf{N}\mathbf{K}^\top)\;\mathbf{s}.
\end{equation}
Every operation in the chain is differentiable: ray tracing, matrix
multiplication, the ReLU, and the weighted sum.  Gradients of
regulatory quantities ($\SARwb$, peak~$\Sab$) with respect to design
parameters propagate by standard backpropagation.  This enables
several optimisation problems to be solved end-to-end.

First, for exposure-aware beamforming, one can optimise antenna weights
to maximise signal quality subject to $\Sab\le 10$\,W/m$^2$.  The
constraint is a linear function of the path powers through
\cref{eq:diff-rt}.
Second, for compliance-by-design placement, one can optimise base-station
locations to minimise worst-case exposure across a population of
user positions.
Finally, for RIS phase optimisation, a reconfigurable intelligent
surface (RIS) modifies the propagation paths.  The dosimetry layer
keeps the chain scene $\to$ RIS phase $\to$ paths $\to$ $\Sab$
fully differentiable.

At the aggregate level (\cref{sec:hierarchy}), the computation
reduces to a weighted sum over the precomputed directivity table.
This costs $O(N)$ per evaluation and runs in the inner loop of a
gradient-based optimiser without touching the body mesh.

When polarisation-aware design is required, the absorption Stokes
vector (\cref{sec:stokes-vector}) replaces the scalar $T_0\,\Aperp$
lookup with a four-component inner product per path.  The gradient of
$P_{\mathrm{abs}}$ with respect to each Stokes parameter is a
component of~$\mathbf{m}(\khat)$, enabling polarisation-aware
beamforming at no additional algorithmic complexity.

\subsection{Real-time dosimetry}\label{sec:app-realtime}

The computational cost of the full spatial formula
(\cref{eq:Sab-mat}) is dominated by the $M\times N$ matrix-vector
product.  For typical values ($M\sim 10^4$ triangles, $N\sim 10^2$
paths) this amounts to ${\sim}\,10^6$ multiply-adds per frame.  A
modern GPU executes this in microseconds.

The exposure fraction~$\eta(\rr)$ is identical to ambient occlusion,
one of the most heavily optimised computations in real-time rendering.
Efficient GPU implementations (screen-space ambient occlusion,
hardware-accelerated ray-traced ambient occlusion) compute it at
interactive frame rates on consumer hardware.  The precomputed body
response (\cref{sec:precomputed}) adds negligible overhead.

The entire dosimetry computation, from ray-tracer
output to a full-body $\Sab$~map, fits within the budget of a single
rendered frame.  This opens two practical scenarios.
First, for digital-twin monitoring, in a digital-twin
deployment, a continuously updated ray-tracing scene provides the
path data.  The dosimetry layer produces a live exposure map of
every person in the scene.
Second, for deployment planning tools, an engineer places
antennas in a 3D model of a building and observes, in real time,
the effect on the $\Sab$ and SAR distribution as users move through the
space.

The low computational cost also makes the framework suitable for
Monte Carlo studies that require millions of random exposure
scenarios (varying body pose, position, and illumination) to
characterise population-level exposure statistics.

\subsection{Reverberation chamber dosimetry}\label{sec:app-reverb}

A reverberation chamber creates a statistically isotropic,
uniformly polarised electromagnetic field inside a shielded enclosure.
These chambers are the standard tool in bioelectromagnetics for
whole-body exposure of small animals, cell cultures, and phantoms.
The dosimetric question is: given a measured or calibrated scalar
power density~$S_{\mathrm{total}}$ inside the chamber, what is the
absorbed power in the specimen?

The framework gives an exact answer.  Because the field is isotropic,
the direction average applies and the generalised Cauchy formula
(\cref{eq:gen-cauchy}) yields
\begin{equation}\label{eq:reverb}
  P_{\mathrm{abs}} = \frac{S_{\mathrm{total}}\,T_0\,\Aab}{4}.
\end{equation}
The local distribution follows from \cref{eq:diffuse}:
\begin{equation}
  \Sab(\rr) = \frac{S_{\mathrm{total}}\,T_0}{4}\;\eta(\rr).
\end{equation}
The exposure fraction~$\eta(\rr)$ is the dosimetry map.  Points in
concavities ($\eta<1$) receive less power; exposed surfaces
($\eta=1$) receive the maximum.
Note that for small animals such as mice and rats (popular subjects in such chambers),
the exposure fraction~$\eta$ is close to unity due to their convex morphology.

No ray tracing, no direction sampling, and no angular integration
are needed.  The entire dose is determined by three precomputed
quantities: $T_0$ (from tissue data and frequency), $\Aab$ (from the
body mesh), and the spatial map~$\eta(\rr)$.  Replacing $T_0$
with~$\Tbar$ (\cref{sec:Tbar}) makes the result exact at any
frequency; the polarisation correction vanishes in an isotropic
field with random polarisation.  For whole-body SAR in
a reverberation chamber the compliance criterion from
\cref{sec:SAR-criterion} simplifies to
\begin{equation}\label{eq:reverb-SAR}
  \SARwb = \frac{S_{\mathrm{total}}\,T_0\,\Aab}{4\,m},
\end{equation}
which depends only on the body's absorption area and mass.

This closed-form result replaces the numerical simulations currently
used to characterise reverberation chamber exposures.  It applies
equally to animal phantoms and to human-body phantoms placed in a
mode-stirred environment.


% ============================================================
\partheading{Part~II: Generalisations and Fundamental Limits}
{Sections~\ref{sec:polarisation}--\ref{sec:sub6}: extensions beyond
Part~I towards the sub-6\,GHz band and arbitrary polarisation.}
% ============================================================


% ============================================================
\section{Polarisation-aware dosimetry}\label{sec:polarisation}
% ============================================================

Part~I set the polarisation correction to zero:
the pseudo-Brewster compensation reduces the Fresnel transmission
to a scalar~$T_0$, and the local TM excess $q\,\Delta T/2$ from
\cref{eq:Teff-decomp} is discarded.  The full polarisation-dependent
framework requires three ingredients.  The main results are
an absorption Stokes vector that computes total absorbed power
from the Stokes parameters of the incident wave, a universal bound
on the fractional correction (at most 27.9\,\% for any convex body,
15.8\,\% on a realistic human phantom), and proof that multipath
diversity suppresses the correction below 2.5\,\% once $N\ge 20$
paths are present.

% -------------------------------------------------------
\subsection{Absorption Stokes vector}\label{sec:stokes-vector}

Integrating the exact local law \cref{eq:Sab-exact} over the
illuminated surface, with the decomposition of
\cref{eq:Teff-decomp} and the TM excess~$q$ from
\cref{eq:q-exact}, gives the total absorbed power
\begin{equation}\label{eq:Pabs-pol-full}
  P_{\mathrm{abs}}
  = \Sinc\,\tilde{A}_\perp(\khat)
  - \frac{\Sinc\cos(2\chi)}{2}\,
    \RE\!\bigl[\tilde{B}(\khat)\,e^{-2i\psi_0}\bigr],
\end{equation}
where the \emph{Fresnel-weighted projected area} is
\begin{equation}\label{eq:Atilde-def}
  \tilde{A}_\perp(\khat)
  \equiv \int_{\Sigma_+}
    \Tavg\!\bigl(\theta(\rr)\bigr)\,\cos\theta(\rr)\;dA,
\end{equation}
and the complex \emph{polarisation response function} is
\begin{equation}\label{eq:Btilde-def}
  \tilde{B}(\khat)
  = B_c(\khat) + i\,B_s(\khat)
  \equiv \int_{\Sigma_+}
    \Delta T\!\bigl(\theta(\rr)\bigr)\,(\cos\theta(\rr)\;
    e^{2i\alpha(\rr)}\;dA.
\end{equation}
Here $\Sigma_+$ denotes the illuminated (front-facing) hemisphere,
$\alpha(\rr)$ is the angle between the local TE direction
$\hat{e}_s(\rr)$ and the fixed reference polarisation basis
$(\hat{e}_1,\hat{e}_2)$ perpendicular to~$\khat$,
and $\chi$ and $\psi_0$ are the ellipticity and orientation angles
of the incident wave's polarisation ellipse.

Both $\tilde{A}_\perp$ and $\tilde{B}$ are precomputable from the
body mesh and the Fresnel coefficients; they do not depend on the
incident wave's polarisation state.

The Stokes parameters of the incident wave in the reference frame
$(\hat{e}_1,\hat{e}_2)$ are
$S_0 = \Sinc$,
$S_1 = \Sinc\cos(2\chi)\cos(2\psi_0)$,
$S_2 = \Sinc\cos(2\chi)\sin(2\psi_0)$, and
$S_3 = \Sinc\sin(2\chi)$.
Substituting into \cref{eq:Pabs-pol-full} yields
\begin{equation}\label{eq:Pabs-Stokes}
  \boxed{P_{\mathrm{abs}}
    = \mathbf{m}(\khat)\cdot\mathbf{s}_{\mathrm{inc}},}
\end{equation}
where $\mathbf{s}_{\mathrm{inc}} = (S_0,S_1,S_2,S_3)^\top$ is the
Stokes vector and the \emph{absorption Stokes vector} of the body is
\begin{equation}\label{eq:m-def}
  \mathbf{m}(\khat)
  = \bigl(\tilde{A}_\perp,\;
    -\tfrac{1}{2}B_c,\;
    -\tfrac{1}{2}B_s,\;
    0\bigr)^\top.
\end{equation}

Three properties follow immediately.
First, the fourth component vanishes: circular polarisation
handedness ($S_3$) does not affect total absorbed power.
Only the degree of linear polarisation matters.
Second, for unpolarised light,
$\mathbf{s} = S_0(1,0,0,0)^\top$, and
$P_{\mathrm{abs}} = \Sinc\,\tilde{A}_\perp$.  Under the
pseudo-Brewster approximation ($\Tavg\approx T_0$),
$\tilde{A}_\perp \approx T_0\,\Aperp$,
recovering the geometric absorption law of Part~I.
Third, for multiple sources,
$P_{\mathrm{abs}} = \sum_i\mathbf{m}(\khat_i)\cdot\mathbf{s}_i$;
the computation is $O(N)$, differentiable, and compatible with the
neural-network formulation of \cref{sec:neural-net}.

\Cref{eq:Pabs-Stokes} connects the dosimetry framework
to standard Mueller calculus: the body surface acts as a
polarisation-dependent absorber whose response is a four-vector
on the Poincar\'e sphere.  The gradient of $P_{\mathrm{abs}}$
with respect to any Stokes parameter is a component
of~$\mathbf{m}$, which makes polarisation-aware network optimisation
straightforward.

\paragraph{Fractional polarisation sensitivity.}
Define the \emph{polarisation directivity}
\begin{equation}\label{eq:DB-def}
  D_B(\khat)
  \equiv \frac{|\tilde{B}(\khat)|}{2\,\tilde{A}_\perp(\khat)}.
\end{equation}
For a linearly polarised wave ($\chi=0$) at the most sensitive
orientation ($2\psi_0$ aligned with $\arg\tilde{B}$), the total
power deviates from the unpolarised prediction by $\pm D_B$:
\begin{equation}\label{eq:Pabs-DB}
  P_{\mathrm{abs}}
  = \Sinc\,\tilde{A}_\perp\bigl(1 \pm D_B\bigr).
\end{equation}

% -------------------------------------------------------
\subsection{Canonical shapes and the cylinder bound}
\label{sec:canonical-pol}

\paragraph{Sphere.}
A sphere's illuminated hemisphere has $\alpha(\rr)$ covering the full
$[0,2\pi)$ range uniformly, so the $e^{2i\alpha}$ integral in
\cref{eq:Btilde-def} vanishes by azimuthal symmetry:
$\tilde{B}=0$ for any direction~$\khat$ (proved in
\cref{app:sphere-proof}).  A sphere absorbs the same total power for
any polarisation state.

This result is not because TE and TM transmissions are equal
at each surface point; they differ substantially at oblique incidence.
A linearly polarised wave hitting a sphere produces a spatially
varying TE/TM decomposition.  In the E-plane (containing $\khat$ and
$\EE_0$), the wave is locally pure TM and transmits more power.  In
the H-plane, it is pure TE and transmits less.  The integral over
the illuminated hemisphere averages these contributions to zero.

\paragraph{Cylinder at broadside.}
An infinite vertical cylinder illuminated horizontally has all
surface normals in a single plane.  The angle $\alpha$ is constant
across the illuminated strip: a vertically polarised wave is pure TE
everywhere, and a horizontally polarised wave is pure TM everywhere.
No azimuthal averaging occurs.  This geometry maximises $|\tilde{B}|$
relative to $\tilde{A}_\perp$, giving the upper bound
\begin{equation}\label{eq:DB-cyl}
  D_B^{\mathrm{cyl}}
  = \frac{\displaystyle\int_0^1\Delta T(\mu)\,\mu\,\diff \mu}
         {\displaystyle 2\int_0^1\Tavg(\mu)\,\mu\,\diff \mu}.
\end{equation}
For skin at 28\,GHz, $D_B^{\mathrm{cyl}} = 27.9\,\%$.  This is the
theoretical worst case over all convex body shapes and all incidence
directions.  When the incidence direction makes an angle~$\gamma$
with the cylinder axis, the constant-$\alpha$ property breaks down
and $D_B$ decreases; the broadside value ($\gamma=90^\circ$) is the
strict maximum.

% -------------------------------------------------------
\subsection{Polarisation response of a human body}
\label{sec:pol-human}

\Cref{fig:pol-mollweide} shows the polarisation response of the
Thelonious phantom at 28\,GHz, computed from exact Fresnel-weighted
surface integrals over 2000 incidence directions.

\begin{figure}[!htbp]
  \centering
  \includegraphics[width=\textwidth]{polarization_mollweide_thelonious.pdf}
  \caption{Polarisation response of the Thelonious phantom at
    28\,GHz in Mollweide projection.
    (a)~Fresnel-weighted projected area
    $\tilde{A}_\perp(\khat)$;
    (b)~polarisation correction magnitude
    $|\tilde{B}(\khat)|$;
    (c)~fractional sensitivity $D_B(\khat)$;
    (d)~histogram of~$D_B$.
    The dashed line in~(d) marks the cylinder bound
    (27.9\,\%).  The maximum $D_B=16.0\,\%$ occurs for
    near-lateral illumination where the cylindrical limbs
    present coherent TE/TM alignment.}
  \label{fig:pol-mollweide}
\end{figure}

The worst-case fractional sensitivity is $D_B = 16.0\,\%$, which is
57\,\% of the cylinder bound.  The most sensitive directions lie about
$30^\circ$ from lateral, where the cylindrical limbs dominate the
integral.  Overhead illumination gives $D_B < 1\,\%$ because the
near-circular cross-section of the head produces the same azimuthal
averaging as a sphere.

\begin{remark}[Bilateral symmetry]
The human body is approximately bilaterally symmetric.  For a standing
human with limbs roughly parallel to the sagittal plane, the
$\tilde{b}_y$ term (which couples to $45^\circ$-polarisation) and
the circular terms $\tilde{a}_\perp$ averaged across the midline
tend to cancel.  This further suppresses the global polarisation
correction relative to a single isolated cylindrical limb. For frontal illumination, $|B_s/B_c| < 0.02$.  The
polarisation correction then reduces to
\begin{equation}\label{eq:DeltaP-frontal}
  \Delta P_{\mathrm{pol}}
  = -\frac{\Sinc\cos(2\chi)}{2}\;B_c(\khat)\;\cos(2\psi_0),
\end{equation}
and only the distinction between H-pol and V-pol matters.
\end{remark}

% -------------------------------------------------------
\subsection{Multipath convergence}\label{sec:multipath-pol}

In a realistic wireless environment, $N$ paths arrive with
independent random polarisation orientations
$\psi_{0,i}\sim\mathrm{Uniform}[0,2\pi)$.  For equal-power paths,
each carrying power $\Sinc/N$, the standard deviation of the
polarisation correction scales as
\begin{equation}\label{eq:pol-convergence}
  \frac{\sigma[\Delta P_{\mathrm{pol}}]}{P_{\mathrm{unpol}}}
  = \frac{D_B(\khat)}{\sqrt{2N}}.
\end{equation}

\Cref{tab:multipath-pol} lists the resulting fractional error for
the worst-case direction on Thelonious ($D_B = 16\,\%$).

\begin{table}[!htbp]
\centering
\caption{Multipath polarisation convergence on Thelonious at 28\,GHz. $\sigma/P$ is the standard deviation of the fractional polarisation correction over 2000 Monte Carlo trials. Theory: $D_B/\!\sqrt{2N}$ with $D_B^{\max} = 15.9\,\%$.}\label{tab:multipath-pol}
\begin{tabular}{r c c c}
\toprule
$N$ & Linear $\sigma/P$ (\%) & Elliptical $\sigma/P$ (\%) & Theory $D_B^{\max}/\!\sqrt{2N}$ (\%) \\
\midrule
  1  & 5.09 & 3.19 & 11.21 \\
  5  & 3.12 & 2.15 & 5.01 \\
 10  & 2.29 & 1.58 & 3.55 \\
 20  & 1.70 & 1.19 & 2.51 \\
 50  & 1.10 & 0.77 & 1.59 \\
100  & 0.80 & 0.54 & 1.12 \\
\bottomrule
\end{tabular}
\end{table}

Monte Carlo simulation (2000 realisations per~$N$) confirms the
$1/\sqrt{N}$ scaling.  Elliptical polarisation gives smaller errors
because $\cos(2\chi)<1$ reduces the correction amplitude.

For indoor 5G environments with $N\sim 20$--$100$ paths, the
polarisation error is 1--2.5\,\%: comparable to or smaller than the
pseudo-Brewster approximation error.  In a reverberation chamber
($N\to\infty$, random polarisation), the correction vanishes exactly.

\begin{figure}[!htbp]
  \centering
  \includegraphics[width=\textwidth]{multipath_convergence.pdf}
  \caption{Monte Carlo validation of the $\sigma/P \sim D_B/\sqrt{2N}$ scaling law for multipath polarisation convergence (Thelonious, 28\,GHz). Markers show Monte Carlo results (2000 trials per $N$); solid lines show the theoretical bounds.}
  \label{fig:multipath-convergence}
\end{figure}

\begin{remark}[Vertical polarisation from base stations]
A standing person illuminated horizontally by a vertically polarised
base station presents cylindrical limbs whose TE direction aligns
with the vertical.  The limbs see pure TE polarisation, for which
$T_s(\theta)\le T_0$ at all angles (\cref{app:TE-bound}).  The
simplified formula therefore overestimates limb absorption.  On the
torso, the incidence angle is small and $\Delta T\approx 0$, so
the polarisation correction is negligible.  For the most common 5G
exposure scenario, the unpolarised geometric absorption law is conservative.
\end{remark}

% -------------------------------------------------------
\subsection{Spherical harmonic representation and compression}
\label{sec:pol-computation}

The polarisation-aware framework adds two precomputed quantities
to the body's response table: $\tilde{A}_\perp(\khat)$ and
$\tilde{B}(\khat)$, alongside the geometric directivity
$D(\khat)$.  For a single incidence direction, the absorbed power
follows from the Stokes inner product \cref{eq:Pabs-Stokes}:
two scalar multiplications and one addition replace the single
$T_0\,\Aperp$ lookup of the unpolarised framework.  The multi-source
formula remains $O(N)$.

\begin{remark}[Spherical harmonic bandwidth]
The directivity $D(\khat)$ and the Fresnel-weighted area
$\tilde{A}_\perp(\khat)$ compress well: $L=4$
(25~coefficients) gives sub-percent error.  The near-identical
compression curves confirm that the pseudo-Brewster compensation
preserves angular structure:
$\tilde{A}_\perp\approx T_0\,\Aperp$.

\begin{figure}[!htbp]
  \centering
  \includegraphics[width=\textwidth]{sh_compression_polarization.pdf}
  \caption{Spherical harmonic compression error (RMS, \%) for the
    three body-response functions of Thelonious at
    28\,GHz.  The geometric directivity~$D(\khat)$ and the
    Fresnel-weighted area $\tilde{A}_\perp(\khat)$ reach sub-percent
    error at $L=4$ (25~coefficients).  The polarisation magnitude
    $|\tilde{B}(\khat)|$ requires $L=9$ (100~coefficients) for
    5\,\% error.}
  \label{fig:sh-compression-pol}
\end{figure}

The polarisation response $|\tilde{B}(\khat)|$ requires
$L=9$ (100~coefficients) for 5\,\% RMS error.  This higher bandwidth
arises from the cancellation of TM excess across the limb circumferences,
which makes the response more sensitive to local surface orientation
than the total-power terms.
\end{remark}
The compression curves show an even-$L$ staircase, a signature of
bilateral symmetry.

The total storage for the full polarisation-aware framework is
approximately 125 spherical harmonic coefficients per body per
frequency (25 for $\tilde{A}_\perp$, 100 for $\tilde{B}$).
The run-time overhead is one additional dot product per incidence
direction.

% -------------------------------------------------------
\subsection{Exact bounds}\label{sec:exact-bounds}

\begin{proposition}[Peak absorbed power density]
The peak $\Sab$ on any body, for any polarisation, occurs at normal
incidence ($\theta=0$), where $\Delta T=0$ and $\Teff=T_0$:
\begin{equation}\label{eq:peak-exact}
  \max_{\rr,\,\mathrm{pol}}\;\Sab(\rr)
  = \Sinc\,T_0.
\end{equation}
This is exact, not approximate.  At $\theta>0$, $T_p(\theta)$ can
increase towards the pseudo-Brewster peak, but the flux factor
$\cos\theta$ decreases faster: the product
$T_p(\theta)\cos\theta$ is monotonically decreasing in~$\theta$ for
all lossy dielectrics tested (skin, fat, muscle, bone, cornea across
0.9 to 100\,GHz).  For TE, $T_s(\theta)$ is itself monotonically
decreasing, so the product $T_s\cos\theta$ is trivially decreasing.

This strengthens the bound in \cref{app:peak-proof} from
$\max\Sab\le 1.017\,\Sinc\,T_0$ to equality at $1.000$.  The
peak-$\Sab$ compliance check ($\Sab < 10$\,W/m$^2$) requires only
$\Sinc < 10/T_0$, for any polarisation, on any body.
\end{proposition}

\paragraph{Polarisation-aware compliance.}
For linear polarisation at the most sensitive direction and
orientation, the worst-case total absorbed power is
\begin{equation}\label{eq:Pabs-worst-pol}
  P_{\mathrm{abs}}^{\mathrm{worst}}
  = \Sinc\bigl[\tilde{A}_\perp(\khat)
    + \tfrac{1}{2}|\tilde{B}(\khat)|\bigr].
\end{equation}
For Thelonious at 28\,GHz, this increases the effective projected
area from $\tilde{A}_\perp = 0.130$\,m$^2$ to 0.148\,m$^2$: a
12\,\% shift in the compliance threshold.  This margin is well within
the 50$\times$ safety factor applied by ICNIRP to derive reference
levels for the general public.


% ============================================================
\section{Extension above and below millimetre waves}\label{sec:sub6}
% ============================================================

Part~I specialises to the 10--60\,GHz band, where the
pseudo-Brewster compensation renders $\Tavg(\theta)\approx T_0$.
Below this band, the approximation error grows: the ratio
$R(f)\equiv T_0/\Tbar$ departs from unity and the local surface
map loses physical meaning as the skin depth exceeds the tissue
layer thickness.  Total-power quantities, however, remain valid at
any frequency where the body is electromagnetically opaque.
The frequency-dependent accuracy is quantified below, together with
the single replacement that makes the integrated framework exact.

% -------------------------------------------------------
\subsection{Exact averaged transmission}\label{sec:Tbar}

Define the \emph{flux-weighted averaged transmission}
\begin{equation}\label{eq:Tbar-def}
  \Tbar(f)
  \equiv 2\int_0^1\Tavg(\mu,f)\;\mu\;\diff \mu.
\end{equation}
This is the fraction of incident power absorbed by a convex body
in an isotropic, unpolarised field.  It is computed once per
frequency from the Fresnel coefficients (a single numerical
quadrature) and depends only on the complex refractive index.
Zhang~\cite{Zhang2017thesis,ZhangRobinson2020} defines an absorption coefficient
$\xi\equiv\langle\sigma_a\rangle/(0.25\times\mathrm{BSA})$, where
$\langle\sigma_a\rangle$ is the measured direction-averaged absorption cross section
and BSA is the body surface area.
The framework identifies $\xi = \Tbar\cdot\Aab/A$: the
Fresnel-weighted transmission $\Tbar$ multiplied by the non-convexity
correction $\Aab/A = \gamma_s\le 1$.
Bamba~\textit{et al.}~\cite{Bamba2014} obtained an absorption efficiency
$\eta\approx 0.5$ from FDTD simulations in diffuse fields at 1.5--6\,GHz;
the framework identifies this with $T_0$ at mmWave, and with $\Tbar$ below
6\,GHz.

Replacing $T_0$ with $\Tbar$ in the Cauchy formula makes the
direction-averaged absorbed power exact:
\begin{equation}\label{eq:Cauchy-exact}
  \boxed{\langle P_{\mathrm{abs}}\rangle
    = \Sinc\,\Tbar\,\frac{\Aab}{4}.}
\end{equation}
The compliance criterion becomes exact as well:
\begin{equation}\label{eq:compliance-exact}
  \Sinc < \frac{0.32\,m}{\Tbar\,\Aab\,D(\khat)}.
\end{equation}
Both identities hold at any frequency, for any body, without the
pseudo-Brewster approximation.

\Cref{tab:Tbar} lists $\Tbar$ and the ratio $R=T_0/\Tbar$ for
skin from 0.3 to 100\,GHz (dielectric data
from~\cite{Gabriel1996,ITISv5}).

\begin{table}[!htbp]
\centering
\caption{Normal-incidence transmission $T_0$, exact averaged
  transmission $\Tbar$, and their ratio $R=T_0/\Tbar$ for skin
  (IT'IS v5.0, Gabriel 4-Cole-Cole model).}\label{tab:Tbar}
\begin{tabular}{rcccc}
\toprule
$f$\,(GHz) & $|\ntilde|$ & $T_0$ & $\Tbar$ & $R$ \\
\midrule
  0.3 & 7.93 & 0.381 & 0.406 & 0.938 \\
  0.9 & 6.70 & 0.445 & 0.465 & 0.957 \\
  2.4 & 6.29 & 0.470 & 0.488 & 0.964 \\
  6.0 & 6.07 & 0.481 & 0.497 & 0.968 \\
 10   & 5.87 & 0.489 & 0.504 & 0.971 \\
 28   & 4.84 & 0.536 & 0.543 & 0.988 \\
 40   & 4.30 & 0.570 & 0.571 & 1.000 \\
 60   & 3.68 & 0.622 & 0.613 & 1.015 \\
100   & 3.01 & 0.701 & 0.677 & 1.035 \\
\bottomrule
\end{tabular}
\end{table}

Below 40\,GHz, $R<1$: the simplified formula with $T_0$
underestimates absorbed power and is therefore conservative.
Above 40\,GHz, $R>1$: the simplified formula overestimates by at
most 3.5\,\% (at 100\,GHz).  The RMS error of the $T_0$
approximation remains below 5\,\% across the entire 0.3--100\,GHz
range.

% -------------------------------------------------------
\subsection{Pseudo-Brewster compensation accuracy}\label{sec:Rf-landscape}

\Cref{fig:R-of-f} shows the ratio $R(f)$ across three decades of
frequency.

\begin{figure}[!htbp]
  \centering
  \includegraphics[width=0.7\textwidth]{R_of_f_landscape.pdf}
  \caption{Pseudo-Brewster accuracy across frequency (skin, IT'IS v5.0).
    Top: ratio $R(f) = T_0/\Tbar$ with conservative region
    ($R<1$, blue) and non-conservative region ($R>1$, red);
    the crossover at 40.4\,GHz marks the pseudo-Brewster sweet spot.
    Middle: $T_0(f)$ and $\Tbar(f)$.
    Bottom: absorption activation functions
    $\Tavg(\theta)\cos\theta$ at six frequencies; the dashed line
    is the $T_0\cos\theta$ approximation at 40\,GHz, where the
    curves nearly overlap.}
  \label{fig:R-of-f}
\end{figure}

The crossover $R=1$ occurs at 40.4\,GHz (interpolated), where the
pseudo-Brewster compensation is exact: $\Tavg(\theta)\cos\theta
\approx T_0\cos\theta$ to four significant figures.  This sweet spot
arises because the competing TE decrease and TM pseudo-Brewster
increase balance perfectly at $|\ntilde|\approx 4.3$.  Below this
frequency, $|\ntilde|$ is larger and the TM peak overshoots; above,
$|\ntilde|$ is smaller and the compensation undershoots.

The bottom panel of \cref{fig:R-of-f} displays the exact absorption
activation function $\Tavg(\theta)\cos\theta$ at six frequencies.
At 40\,GHz the exact and approximate curves are indistinguishable.
At lower frequencies the exact curve lies above $T_0\cos\theta$
at oblique angles because of the stronger TM enhancement; the
integrated excess produces $\Tbar > T_0$.

\begin{remark}[Multi-layer tissue below 6\,GHz]
At frequencies below 6\,GHz, the skin depth increases and the wave
can penetrate through the skin and fat into the muscle layer.
Destructive or constructive interference (Fabry--P\'erot effects)
between reflections from these interfaces can modify the transmission.
However, for the sphere-averaged transmission~$\Tbar$, these
modulations are largely suppressed when integrated over angle and
frequency.  The pseudo-Brewster compensation accuracy (\cref{fig:R-of-f})
remains the dominant qualitative feature, justifying the use of an
effective $\Tbar$ even in the presence of multi-layer interference.
A transfer-matrix analysis of a skin (2\,mm) / fat (10\,mm) / muscle (semi-infinite)
stack reveals Fabry--P\'erot resonances: at 0.9\,GHz the fat layer
acts as a quarter-wave impedance match and absorption increases by
40\,\%, while at 3.5\,GHz destructive interference reduces
absorption by 27\,\%.  Fat thickness varies from approximately 2\,mm
(limbs) to 30\,mm (abdomen), so different body regions hit different
resonance conditions.  Body-surface averaging suppresses the
oscillations, and the single-layer $\Tbar$ remains a reasonable
estimate for whole-body SAR.  Local surface maps require
tissue-layer modelling below 6\,GHz.
\end{remark}

% -------------------------------------------------------
\subsection{Validity and regime limits}\label{sec:traffic-light}

\Cref{tab:traffic-light} summarises the accuracy of each framework
capability across frequency.

\begin{table}[!htbp]
\centering
\caption{Framework validity across frequency.  Errors are worst-case
  over all body shapes and incidence directions; ``conservative''
  means the framework underestimates
  absorption.}\label{tab:traffic-light}
\begin{tabular}{p{0.42\textwidth}p{0.16\textwidth}p{0.28\textwidth}}
\toprule
Capability & Freq.\ range & Validity \\
\midrule
Geometric quantities ($D$, $\eta$, Cauchy, SH)
  & all & exact (purely geometric) \\[2pt]
Total $P_{\mathrm{abs}}$, $\SARwb$, compliance (with $\Tbar$)
  & $>$100\,MHz & exact \\[2pt]
Polarisation framework (Stokes, $\tilde{B}$): total power
  & all & exact \\[2pt]
Polarisation framework (Stokes, $\tilde{B}$): local map
  & $>$6\,GHz & $<$5\,\%, conservative \\[2pt]
Geometric absorption law
  $\Sab\approx\Sinc T_0\,\ReLU(\mu)$
  & $>$10\,GHz & $<$5\,\%, conservative \\[2pt]
Local $\Sab$ map: torso-scale objects
  & $>$6\,GHz & $<$10\,\%, conservative \\[2pt]
Local $\Sab$ map: finger-scale objects
  & $>$28\,GHz & $<$20\,\%, conservative \\[2pt]
Internal SAR depth distribution
  & $<$6\,GHz & Volumetric computation required \\[2pt]
Reactive near-field ($d<\lambda/(2\pi) \approx \text{cm--mm}$)
  & any & Full-wave simulation required \\
\bottomrule
\end{tabular}
\end{table}

The geometric results (projected area, directivity, Cauchy's formula,
spherical harmonics) are purely mathematical and frequency-independent.
The total-power and compliance results extend to any frequency where
the body is electromagnetically opaque ($f\gtrsim 100$\,MHz), provided
$T_0$ is replaced by~$\Tbar$.  The local surface map and the
curvature correction break down below approximately 6\,GHz, where
volumetric absorption, diffraction, and multi-layer resonances
dominate.  These frequency-dependent limitations are the irreducible
limitations of \cref{sec:limitations-v2}.





% ============================================================
\partheading{Part~III: Coherent MIMO Illumination}
{Sections~\ref{sec:coherent-setup}--\ref{sec:mu-mimo}: the
incoherent framework extended to beamformed sources with controlled
complex weights.}
% ============================================================

Parts~I and~II assume incoherent illumination: the $N$ incident
waves arrive from statistically independent sources with random
relative phases, so cross-terms average to zero and absorbed power
sums over paths.  When a massive-MIMO base station performs coherent
beamforming, $M$ antenna elements radiate with controlled complex
weights and their fields add as amplitudes.  For $N$ co-phased,
co-polarised paths the coherent peak exceeds the incoherent sum by a
factor of $N$; the framework of Parts~I--II cannot capture this
enhancement.  This Part derives the exact coherent absorption law
under two well-bounded approximations and collapses the result to a
finite-dimensional matrix problem.

% ============================================================
\section{Setup}\label{sec:coherent-setup}
% ============================================================

A base station has $M$ antenna elements.  It radiates with
precoding vector $\mathbf{x} = [x_1, \ldots, x_M]^T \in
\mathbb{C}^M$, $\|\mathbf{x}\|^2 = P$ (total radiated power in W, so
$[x_j] = \sqrt{\text{W}}$).  Between the base station and the user
body, the propagation environment creates $N$ paths.  Path
$n = 1,\ldots,N$ is characterised by its originating element
$j(n) \in \{1,\ldots,M\}$, its direction of arrival
$\hat{\mathbf{k}}_n$ at the body, and its complex
polarisation-amplitude vector $\boldsymbol{\psi}_n \in \mathbb{C}^3$
($\boldsymbol{\psi}_n \perp \hat{\mathbf{k}}_n$, units:
$\text{V}\,\text{m}^{-1}\text{W}^{-1/2}$), which encodes the path's
amplitude, polarisation, and all propagation interactions.

Three inherited assumptions apply throughout.
First, the body surface $\Sigma$ is locally flat on the scale of
$\lambda$.
Second, the skin depth $\delta$ is much smaller than any body dimension,
so all transmitted power is absorbed within the surface layer.
Third, every path's last scattering point is at distance $d > 3\lambda$
from the surface, so the field at each body-surface point is a
superposition of locally plane waves.  At 28\,GHz this requires
$d > 3.2$\,cm, satisfied for any device not in direct skin contact.

% ============================================================
\section{Field channel matrix}\label{sec:field-channel}
% ============================================================

The total electric field at a body-surface point determines the
absorbed power there.  We assemble that field from the propagation
paths before applying Fresnel physics.

% -------------------------------------------------------
\subsection{Propagation paths}\label{sec:antenna-paths}

Each antenna element's far field in free space is fully described by
its antenna pattern $\mathbf{C}(\khat) \in \mathbb{C}^2$,
whose two components give the $\hat{\boldsymbol{\theta}}$ and
$\hat{\boldsymbol{\varphi}}$ field amplitudes in direction~$\khat$,
normalised so that $\|\mathbf{C}(\khat)\|^2$ equals the directional
gain~$G(\khat)$.  Write $\mathbf{C}_{T,j}(\khat)$ for BS
element~$j$'s transmit pattern and $\mathbf{C}_R(\khat)$ for the UE
receive pattern.

Between the BS and the user's body, the propagation environment
creates $N$ paths via reflection, scattering, and diffraction.
Path~$n$ departs from element $j(n)$ in direction
$\hat{\mathbf{k}}_{T,n}$ and arrives at the body from direction
$\hat{\mathbf{k}}_n$.  Each path carries a $2\times 2$ transfer
matrix $\mathbf{T}_n \in \mathbb{C}^{2\times 2}$ that encodes all
scattering interactions, distance attenuation, and propagation phase
along the path.  In the standard multipath field
model~\cite{SionnaEM}, the arriving field phasor from path~$n$ (for
unit excitation of element $j(n)$) is
\begin{equation}\label{eq:e-n-phasor}
  \mathbf{e}_n = \sqrt{\frac{Z_0}{2\pi}}\;\frac{\lambda}{4\pi}\;
  \mathbf{T}_n\;\mathbf{C}_{T,j(n)}(\hat{\mathbf{k}}_{T,n}),
\end{equation}
expressed in the
$(\hat{\boldsymbol{\theta}},\hat{\boldsymbol{\varphi}})$ basis at
$\hat{\mathbf{k}}_n$. $Z_0$ is the vacuum impedance. Lifting to a 3D Cartesian vector via the
local frame
$[\hat{\boldsymbol{\theta}}(\hat{\mathbf{k}}_n),\,
\hat{\boldsymbol{\varphi}}(\hat{\mathbf{k}}_n)]$ gives the
polarisation-amplitude vector introduced in
\cref{sec:coherent-setup}:
\begin{equation}\label{eq:psi-def}
  \boldsymbol{\psi}_n = \sqrt{\frac{Z_0}{2\pi}}\;\frac{\lambda}{4\pi}\;
  \bigl[\hat{\boldsymbol{\theta}}(\hat{\mathbf{k}}_n),\;
  \hat{\boldsymbol{\varphi}}(\hat{\mathbf{k}}_n)\bigr]\;
  \mathbf{T}_n\;\mathbf{C}_{T,j(n)}(\hat{\mathbf{k}}_{T,n}).
\end{equation}
The per-path field contribution at surface point $\mathbf{r}$ is
$\boldsymbol{\psi}_n\,x_{j(n)}\,e^{-ik_0\hat{\mathbf{k}}_n\cdot\mathbf{r}}$:
the vector $\boldsymbol{\psi}_n$ is determined entirely by the
propagation environment and transmit hardware; the beamforming
weight $x_{j(n)}$ is the only degree of freedom.

% -------------------------------------------------------
\subsection{Superposition}\label{sec:G-def}

At any point $\mathbf{r}$, the fields from all paths add as
complex amplitudes.  Element $j$ contributes (with unit weight)
\begin{equation}
  \mathbf{g}_j(\mathbf{r})
  = \sum_{n:\,j(n)=j} \boldsymbol{\psi}_n\;
    e^{-ik_0\hat{\mathbf{k}}_n \cdot \mathbf{r}}.
\end{equation}
Only phase \emph{differences} between paths at the same point matter;
the absolute phase is convention.

\begin{definition}[Field channel matrix]\label{def:field-channel}
The field channel matrix is
\begin{equation}\label{eq:G-def}
  \mathbf{G}(\mathbf{r})
  = \bigl[\mathbf{g}_1(\mathbf{r}),\ldots,\mathbf{g}_M(\mathbf{r})\bigr]
  \;\in\; \mathbb{C}^{3\times M}.
\end{equation}
\end{definition}
The total electric field under precoding $\mathbf{x}$ is then
\begin{equation}\label{eq:E-Gx}
  \boxed{\mathbf{E}(\mathbf{r}) = \mathbf{G}(\mathbf{r})\,\mathbf{x},}
\end{equation}
exact for any precoder in the locally-plane-wave regime.

% -------------------------------------------------------
\subsection{Signal branch}\label{sec:ue-mrt}

A user equipment (UE) at position $\mathbf{r}_\mathrm{UE}$ has
receive antenna pattern $\mathbf{C}_R(\hat{\mathbf{k}}) \in
\mathbb{C}^2$ in each arrival direction.  The received signal is
$y = \mathbf{h}^T\mathbf{x}$, where the UE channel vector
$\mathbf{h} \in \mathbb{C}^M$ has entries
\begin{equation}\label{eq:h-def}
  h_j = \sum_{n:\,j(n)=j}
    \underbrace{\mathbf{C}_R(\hat{\mathbf{k}}_n)^H
    \boldsymbol{\psi}_n}_{=\,a_n}\;
    e^{-ik_0\hat{\mathbf{k}}_n\cdot\mathbf{r}_\mathrm{UE}}.
\end{equation}
The scalar $a_n \in \mathbb{C}$ projects each path's 3D
polarisation-amplitude vector onto the UE antenna response
(implicitly via
$[\hat{\boldsymbol{\theta}},\hat{\boldsymbol{\varphi}}]^T$, since
$\boldsymbol{\psi}_n \perp \hat{\mathbf{k}}_n$).  Entry $h_j$ sums
over all paths from element $j$; one weight $x_j$ controls them
simultaneously.  The UE antenna collapses each path's vector field
to a scalar, discarding the polarisation and direction structure
that the body surface retains (\cref{sec:Fn-def}).

Maximum-ratio transmission (MRT) sets
\begin{equation}\label{eq:MRT}
  \mathbf{x}_\mathrm{MRT} = \sqrt{P}\;\frac{\mathbf{h}^*}{\|\mathbf{h}\|},
\end{equation}
which maximises received signal power $|y|^2 = P\|\mathbf{h}\|^2$
subject to $\|\mathbf{x}\|^2 = P$ (Cauchy--Schwarz, equality when
$\mathbf{x} \propto \mathbf{h}^*$).  Because a single weight $x_j$
simultaneously modifies every co-originating path, MRT cannot steer
towards the UE without also modifying illumination of all bodies in
the scene.

The signal branch is complete: from the propagation paths we have
the UE channel~$\mathbf{h}$ and the MRT precoder.  The remainder of
Part~III derives the exposure branch: how the same paths contribute
to absorbed power on the body surface.

% -------------------------------------------------------
\subsection{Exact Fresnel amplitudes in the coherent regime}\label{sec:why-exact-fresnel}

The pseudo-Brewster approximation of \cref{sec:pseudo-brewster}
replaces the angle-dependent amplitude coefficients $t_s(\theta)$,
$t_p(\theta)$ by the real constant $\sqrt{T_0}$.  In the coherent
regime this replacement fails.  The absorbed power involves
cross-products $t_{s,n}^* t_{s,n'}$ at different incident angles,
which carry phase information; replacing $t_s$ by $\sqrt{T_0}$
collapses every cross-term to a real number, destroying the
interference structure entirely.  \Cref{fig:fresnel-curves} shows
the amplitude magnitudes and phases of $t_s$, $t_p$ for skin at
28\,GHz: the phases drift by $\sim\!20^\circ$ across the
front-facing hemisphere.  Exact Fresnel amplitude coefficients are
therefore retained throughout Part~III.

\begin{figure}[!htbp]
  \centering
  \includegraphics{fresnel_curves.pdf}
  \caption{Fresnel transmission for skin at 28\,GHz
    ($\tilde{n} = 4.49-1.79i$, $T_0 = 0.539$).
    Panel~(a): power transmittances $T_s$ (TE), $T_p$ (TM), their
    average $\Tavg$, and the normal-incidence reference $T_0$.
    The TE and TM curves diverge strongly; the $\Tavg\approx T_0$
    shortcut fails for coherent illumination.
    Panel~(b): amplitude magnitudes $|t_s|$, $|t_p|$ (left axis)
    and phases $\arg(t_s)$, $\arg(t_p)$ (right axis, degrees).
    Phase drift with angle is the information that coherent
    cross-terms exploit.}
  \label{fig:fresnel-curves}
\end{figure}

\Cref{fig:channel-diagram} shows the complete pipeline.  From
$\mathbf{G}$, two branches diverge.  The signal branch (left, red)
collapses each path to a scalar via the UE antenna projection~$\mathbf{C}_R^H$,
yielding~$\mathbf{h}$ and the MRT precoder.  The exposure branch
(right, blue) applies Fresnel transmission per path at each
body-surface point, yielding an exposure
channel~$\tilde{\mathbf{G}}$ and an exposure
operator~$\mathbf{Q}$; the exposure-constrained precoder (centre,
dashed) balances both objectives.  The key asymmetry is that the
signal branch discards the vector structure of the field, while the
exposure branch retains it.  The following sections derive each
block.

\begin{figure}[!htbp]
  \centering
  \includegraphics[width=\textwidth]{channel_block_diagram.pdf}
  \caption{Architecture of coherent geometric dosimetry.
    Propagation paths are assembled per transmit element into the
    field channel matrix~$\mathbf{G}(\mathbf{r})$.
    \textit{Left (signal branch):} each path is projected to a
    scalar via~$\mathbf{C}_R^H$, yielding the UE
    channel~$\mathbf{h}$ and the MRT precoder.
    \textit{Right (exposure branch):} the Fresnel
    operator~$\mathbf{F}_n(\mathbf{r})$ and depth-integral weight
    are applied per path, yielding~$\tilde{\mathbf{G}}(\mathbf{r})$
    and the exposure operator~$\mathbf{Q}$.
    \textit{Centre:} the exposure-constrained precoder balances
    signal quality and body absorption.}
  \label{fig:channel-diagram}
\end{figure}

% ============================================================
\section{Fresnel transmission}\label{sec:fresnel-operator}
% ============================================================

The absorbed power density at a surface point depends on the
transmitted field inside the tissue, not the incident field.
Each plane-wave component satisfies the boundary conditions
independently, so the total transmitted field inside tissue is
the sum of the individually Fresnel-transmitted fields.  This
section derives the per-path transmitted field under one
well-bounded approximation and defines the operator that maps
each path's polarisation-amplitude vector to its transmitted
contribution.

% -------------------------------------------------------
\subsection{Single-path transmitted field}\label{sec:single-path}

\Cref{fig:incidence-plane} illustrates the geometry.  At a surface
point $\mathbf{r}$, path $n$ has incidence cosine
$\mu_n = \hat{\mathbf{n}}(\mathbf{r})\cdot(-\hat{\mathbf{k}}_n)$.
For a front-facing path ($\mu_n > 0$), the TE and TM basis vectors
are
\begin{equation}\label{eq:TE-TM-basis}
  \hat{e}_{s,n} = \frac{\hat{\mathbf{k}}_n \times \hat{\mathbf{n}}}{|\hat{\mathbf{k}}_n \times \hat{\mathbf{n}}|},
  \qquad
  \hat{e}_{p,n} = \hat{e}_{s,n} \times \hat{\mathbf{k}}_n.
\end{equation}
The TE and TM scalar amplitudes are $\psi_{s,n} = \hat{e}_{s,n}
\cdot \boldsymbol{\psi}_n$ and $\psi_{p,n} = \hat{e}_{p,n}\cdot
\boldsymbol{\psi}_n$.  The Fresnel amplitude transmission
coefficients follow from \cref{eq:rs,eq:rp} via $t = 1 + r$:
\begin{equation}\label{eq:fresnel-t}
  t_{s,n} = \frac{2\mu_n}{\mu_n + \xi_n}, \qquad
  t_{p,n} = \frac{2\tilde{n}\,\mu_n}{\tilde{n}^2\mu_n + \xi_n},
\end{equation}
with $\xi_n = \sqrt{\tilde{n}^2 - 1 + \mu_n^2}$ and
$\operatorname{Re}(\xi_n) > 0$ (\cref{eq:xi-def}).
Writing $k_0\xi_n = \beta_n - i\alpha_n$ with $\alpha_n > 0$
(amplitude decay rate into tissue) and $\beta_n > 0$ (phase rate
into depth), the depth factor reads
$e^{-ik_0\xi_n z} = e^{-\alpha_n z}e^{-i\beta_n z}$.

\begin{figure}[!htbp]
  \centering
  \includegraphics{incidence_plane_fig.pdf}
  \caption{Incidence plane geometry for path $n$.
    The TE basis vector $\hat{e}_{s,n}$ (green) points out of the
    incidence plane; the incident TM direction $\hat{e}_{p,n}$
    (orange) lies in it with large normal component $\sin\theta_n$.
    The refracted ray $\hat{\mathbf{k}}'_n$ (dashed grey) enters at
    the much smaller angle $\theta_t$ due to the high refractive
    index; the refracted TM direction $\hat{e}'_{p,n}$ (blue) is
    nearly tangential, with normal component $\sin\theta_t =
    O(1/|\tilde{n}|)$.  Approximation~1 replaces $\hat{e}'_{p,n}$ by
    $\hat{e}_{p,n}$; the error in the normal component is
    $\sin\theta_n(1 - 1/|\tilde{n}|)$.}
  \label{fig:incidence-plane}
\end{figure}

\begin{proposition}[Approximation~1: TM polarisation direction]\label{prop:approx1}
The physically exact transmitted TM field points along the
refracted direction $\hat{e}'_{p,n}$ inside the tissue, not the
incident direction $\hat{e}_{p,n}$.  Let $\hat{\mathbf{t}}_n$ be the
unit tangent in the incidence plane (so
$\hat{\mathbf{k}}_n = \sin\theta_n\,\hat{\mathbf{t}}_n -
\cos\theta_n\,\hat{\mathbf{n}}$).  The two TM directions are
\begin{equation}\label{eq:ep-vs-epp}
  \hat{e}_{p,n}  = \cos\theta_n\,\hat{\mathbf{t}}_n
                  + \sin\theta_n\,\hat{\mathbf{n}},
  \qquad
  \hat{e}'_{p,n} = \cos\theta_{t,n}\,\hat{\mathbf{t}}_n
                  + \sin\theta_{t,n}\,\hat{\mathbf{n}},
\end{equation}
where $\sin\theta_{t,n} = \sin\theta_n/|\tilde{n}|$ by Snell's law.
These vectors differ by $\theta_n - \theta_{t,n}$ in angle (about
$32^\circ$ at $\theta_n = 40^\circ$, $|\tilde{n}| = 4.83$).

We replace $\hat{e}'_{p,n}$ by $\hat{e}_{p,n}$.  Self-terms are
unaffected ($|\hat{e}_{p,n}|^2 = |\hat{e}'_{p,n}|^2 = 1$).  The
error on TM--TM cross-terms is $O(1/|\tilde{n}|^2) \lesssim 4\,\%$
for skin at 28\,GHz; see \cref{tab:approx-accuracy}.
\hl{todo: this bound is unsound as stated. The substitution error on
the inner product is O(1), not O(1/n2): up to 90 percent per pair, up
to 35 percent on a depth-integrated TM-TM cross-term (JSAC2
agent-report audit of 2026-05-11, Issue 3 / T3). The honest statement
lives on the integrated operator Q, where the surface phase integral
decorrelates the bad pairs: error below 0.1 percent whole-body, about
6 percent median on a 4 cm2 patch (theory/scripts/validate\_approx1.py).
Either adopt that framing or drop the approximation and use the exact
refracted direction, as the TWC paper now does.}
\end{proposition}

Under Approximation~1 the transmitted field of path $n$ at depth $z$
inside the tissue is
\begin{equation}\label{eq:Etrans}
  \mathbf{E}_n^\mathrm{trans}(\mathbf{r}, z)
  = \bigl(t_{s,n}\psi_{s,n}\hat{e}_{s,n}
         + t_{p,n}\psi_{p,n}\hat{e}_{p,n}\bigr)\,x_{j(n)}\,
    e^{-ik_0\hat{\mathbf{k}}_n\cdot\mathbf{r}}\;
    e^{-ik_0\xi_n z}\;H(\mu_n).
\end{equation}
For a back-facing path ($\mu_n \leq 0$) the Heaviside factor $H(\mu_n)$ vanishes, so
$\mathbf{E}_n^\mathrm{trans} = \mathbf{0}$.

\begin{remark}[Approximation~1 error budget]\label{rem:approx1-why}
The vectors $\hat{e}_{p,n}$ and $\hat{e}'_{p,n}$ are not close in
angle, yet the error in absorbed power is small.  Absorbed power
depends on dot products between transmitted directions, not on the
directions themselves.  For $n \neq n'$, the exact cross-term
contains $\hat{e}'_{p,n}\cdot\hat{e}'_{p,n'} =
\cos\theta_{t,n}\cos\theta_{t,n'} +
\sin\theta_{t,n}\sin\theta_{t,n'}$.  Since $\sin\theta_{t,n} =
O(1/|\tilde{n}|)$, the tangential product dominates and the
normal--normal product is $O(1/|\tilde{n}|^2)$.  The physical origin
is that the normal electric field inside tissue is $E_\perp \propto
\sin\theta_t = O(1/|\tilde{n}|)$, so its contribution to
$|\mathbf{E}^\mathrm{trans}|^2$ is $O(1/|\tilde{n}|^2)$.  This is a
non-conservative error of $\lesssim 4\,\%$ on TM--TM cross-terms.
\end{remark}

% -------------------------------------------------------
\subsection{Fresnel transmission operator}\label{sec:Fn-def}

\begin{definition}[Fresnel transmission operator]\label{def:fresnel-op}
The \emph{Fresnel transmission operator} for path $n$ at $\mathbf{r}$
is
\begin{equation}\label{eq:Fn}
  \mathbf{F}_n(\mathbf{r}) =
  \Bigl(
    t_{s,n}\;\hat{e}_{s,n}\hat{e}_{s,n}^T
    + t_{p,n}\;\hat{e}_{p,n}\hat{e}_{p,n}^T
  \Bigr)H(\mu_n).
\end{equation}
\end{definition}

$\mathbf{F}_n$ is dimensionless and rank-$2$: it projects
$\boldsymbol{\psi}_n$ onto the TE and TM directions and scales by
the transmission coefficients.  Different paths at the same surface
point have different incidence planes and therefore different
$\mathbf{F}_n$.  The transmitted surface field is then
$\mathbf{E}_n^\mathrm{trans}(\mathbf{r},0) = \mathbf{F}_n(\mathbf{r})
\boldsymbol{\psi}_n\,x_{j(n)}\,e^{-ik_0\hat{\mathbf{k}}_n\cdot\mathbf{r}}$.

\begin{remark}[Relation to scalar calibration approaches]
Castellanos~\textit{et al.}~\cite{Castellanos2020} used per-element scalar
Fresnel transmission coefficients at 28\,GHz, replacing the full
operator $\mathbf{F}_n$ by a single calibration scalar $\beta\approx 0.70$
that lumps the Fresnel transmission with depth-decay effects.
The operator~$\mathbf{F}_n$ retains the TE/TM polarisation structure and the
angle-dependent phase; both are necessary in the coherent regime to correctly
compute cross-term interference.
\end{remark}

% ============================================================
\section{Coherent absorption law}\label{sec:coherent-absorption-law}
% ============================================================

The absorbed power density from the coherent superposition is the
ohmic loss integrated over the skin depth.  We show that the
resulting double sum over all path pairs reduces to a single squared
norm under one further approximation.

% -------------------------------------------------------
\subsection{Coherent depth integral}\label{sec:depth-integral}

The total transmitted field is
$\mathbf{E}^\mathrm{trans} = \sum_n \mathbf{E}_n^\mathrm{trans}$.
The absorbed power density is the ohmic loss integrated over depth:
\begin{equation}\label{eq:Sab-double-sum}
  \Sab(\rr) = \frac{\sigma}{2}\int_0^\infty
  \left|\sum_n \mathbf{E}_n^\mathrm{trans}(\rr,z)\right|^2 \diff z.
\end{equation}
Expanding the squared norm and using the depth dependence
$e^{-ik_0\xi_n z} = e^{-\alpha_n z}e^{-i\beta_n z}$ from
\cref{eq:Etrans}, every cross-pair $(n,n')$ contributes a
\emph{depth coupling factor}:
\begin{equation}\label{eq:Lambda}
  \Lambda_{nn'} = \int_0^\infty
  e^{-(\alpha_n+\alpha_{n'})z}\,e^{i(\beta_n-\beta_{n'})z}\,\diff z
  = \frac{1}{\alpha_n + \alpha_{n'} - i(\beta_{n'}-\beta_n)}.
\end{equation}
The matrix $\Lambda$ is Hermitian positive-semidefinite
($\Lambda_{n'n} = \Lambda_{nn'}^*$; diagonal entries
$\Lambda_{nn} = 1/(2\alpha_n) \in \mathbb{R}^+$).  The absorbed
power density is then
\begin{equation}\label{eq:Sab-Lambda}
  \Sab(\rr) = \frac{\sigma}{2}\sum_{n,n'}
  \bigl[\mathbf{E}_n^\mathrm{trans}(\rr,0)\bigr]^* \cdot
  \mathbf{E}_{n'}^\mathrm{trans}(\rr,0)\;\Lambda_{nn'}.
\end{equation}

The diagonal terms $n = n'$ recover the single-wave absorbed power
density.  A useful identity connects the depth-integral prefactor to
the power transmittance.

\begin{lemma}[Depth-integral identity]\label{lem:depth-id}
For TE polarisation,
\begin{equation}\label{eq:depth-id}
  \frac{\sigma}{4\alpha_n}\,|t_{s,n}|^2 = \frac{T_{s,n}\cos\theta_n}{Z_0}
\end{equation}
exactly, where $\alpha_n = k_0\operatorname{Im}(\xi_n)$.  The same identity holds
for TM polarisation to within $\sim\!2\,\%$; the correction is
$O(1/|\tilde{n}|^2)$, with the same physical origin as
Approximation~1 (\cref{rem:approx1-why}): the normal electric field
inside tissue is $O(1/|\tilde{n}|)$.  The proof is in
\cref{app:depth-identity}.
\end{lemma}

% -------------------------------------------------------
\subsection{Universal depth coupling}\label{sec:approx2}

\begin{proposition}[Approximation~2: universal depth coupling]\label{prop:approx2}
Define the normalised depth coupling:
\begin{equation}\label{eq:Gamma}
  \Gamma_{nn'} \equiv \frac{\Lambda_{nn'}}{\sqrt{\Lambda_{nn}\Lambda_{n'n'}}}
  = \frac{2\sqrt{\alpha_n\alpha_{n'}}}{\alpha_n + \alpha_{n'} - i(\beta_{n'}-\beta_n)}.
\end{equation}
By construction $\Gamma_{nn} = 1$.  We approximate
$\Gamma_{nn'} \approx 1$ for all pairs.

For skin at 28\,GHz, a sweep over all incidence-angle pairs
$(\theta_n, \theta_{n'}) \in [0^\circ, 85^\circ]^2$ gives
$|1 - \Gamma_{nn'}| \leq 0.44\,\%$ with mean $0.17\,\%$
(\cref{fig:approx2-validation}).  The high refractive index
$|\tilde{n}| \approx 4.83$ compresses all refraction angles to within
$\sim\!10^\circ$ of normal, making $\alpha_n$ and $\beta_n$ nearly
angle-independent.  This is a non-conservative error of $\leq 0.5\,\%$
on any cross-term.
\hl{todo: 0.44 / 0.17 percent came from an alpha-beta unpack swap in
validate\_approx2.py (the helper returns beta first, the caller reads
alpha first). Correct values at 28 GHz: max 2.5 percent, mean 1.0
percent. The error is a first-order phase term and can be cancelled by
referencing each path's phase at the mean absorption depth 1/(2 alpha),
which brings it below 0.05 percent at no cost. See erratum 2 in
errata\_monograph.txt.}
\end{proposition}

\begin{figure}[!htbp]
  \centering
  \includegraphics{approx2_gamma_validation.pdf}
  \caption{Numerical validation of Approximation~2 for skin at
    28\,GHz ($\tilde{n} = 4.49 - 1.79i$).
    Panel~(a): $|1 - \Gamma_{nn'}|$\,(\%) over all
    $(\theta_n, \theta_{n'}) \in [0^\circ,85^\circ]^2$; the error
    stays below $0.44\,\%$ everywhere.
    Panel~(b): two representative 1-D slices confirming the
    approximation is uniform.  Skin's high refractive index
    compresses all refraction angles below $\sim\!10^\circ$, making
    $\alpha_n$ and $\beta_n$ nearly angle-independent.
    \hl{todo: figure shows the swapped computation. Regenerate with
    the fixed validate\_approx2.py, colorbar tops near 2.6 percent.}}
  \label{fig:approx2-validation}
\end{figure}

% -------------------------------------------------------
\subsection{Absorbed power as a squared norm}\label{sec:factorisation}

Setting $\Gamma_{nn'} = 1$ in \cref{eq:Sab-Lambda}, so that
$\Lambda_{nn'} = \sqrt{\Lambda_{nn}\Lambda_{n'n'}}$:
\begin{equation}
  \Sab(\rr) \approx \frac{\sigma}{2}\sum_{n,n'}
  \sqrt{\Lambda_{nn}}\;[\mathbf{E}_n^\mathrm{trans}(\rr,0)]^*
  \cdot \sqrt{\Lambda_{n'n'}}\;\mathbf{E}_{n'}^\mathrm{trans}(\rr,0).
\end{equation}
The double sum factors into the expansion of a squared norm
(since $\sqrt{\Lambda_{nn}} \in \mathbb{R}^+$):
\begin{equation}
  \Sab({\rr}) \approx \left\|\sum_n \sqrt{\frac{\sigma}{4\alpha_n}}\;
      \mathbf{F}_n({\rr})\,\boldsymbol{\psi}_n\,x_{j(n)}\,
      e^{-ik_0\hat{\mathbf{k}}_n\cdot{\rr}}\right\|^2.
\end{equation}
For fixed element $j$, the weight $x_{j(n)} = x_j$ is the same for
every co-originating path, which factors out of the inner sum.
Grouping by originating element $j$ therefore defines the
\emph{exposure channel vector}:
\begin{equation}\label{eq:gtilde-def}
  \tilde{\mathbf{g}}_j(\rr) =
  \sum_{n:\,j(n)=j} \sqrt{\frac{\sigma}{4\alpha_n}}\;
  \mathbf{F}_n(\rr)\,\boldsymbol{\psi}_n\,
  e^{-ik_0\hat{\mathbf{k}}_n\cdot\rr}
  \;\in\; \mathbb{C}^3,
\end{equation}
and the \emph{exposure channel matrix}
\begin{equation}\label{eq:Gtilde-def}
  \tilde{\mathbf{G}}(\rr) =
  \bigl[\tilde{\mathbf{g}}_1(\rr),\ldots,\tilde{\mathbf{g}}_M(\rr)\bigr]
  \;\in\; \mathbb{C}^{3\times M}.
\end{equation}

The three channel objects derived from the same propagation paths
are compared in the following table.

\begin{center}
\begin{tabular}{lccc}
\toprule
 & $\mathbf{G}(\mathbf{r})$ & $\mathbf{h}$ & $\tilde{\mathbf{G}}(\mathbf{r})$ \\
\midrule
Domain & $\mathbb{C}^{3\times M}$, varies over $\mathbb{R}^3$ &
  $\mathbb{C}^M$, at $\mathbf{r}_\mathrm{UE}$ &
  $\mathbb{C}^{3\times M}$, varies over $\Sigma$ \\
Per-path operation & identity & $\mathbf{C}_R^H$ (scalar projection) &
  $\mathbf{F}_n$ (Fresnel filter) \\
Output & 3D field $\mathbf{E}(\mathbf{r})$ & scalar signal $y$ &
  absorbed-power field \\
Depends on & paths & + UE antenna & + body normal, tissue \\
\bottomrule
\end{tabular}
\end{center}

The exposure channel retains the full 3D vector structure
of the field; the UE channel discards it.  This structural difference
is the geometric origin of the decoupling between signal quality and
body absorption.

% -------------------------------------------------------
\subsection{Main theorem}\label{sec:coherent-law-statement}

\begin{theorem}[Coherent absorption law]\label{thm:coherent-law}
Under Approximations~1 and~2,
\begin{equation}\label{eq:coherent-law}
  \boxed{\Sab(\rr) = \|\tilde{\mathbf{G}}(\rr)\,\mathbf{x}\|^2
  = \mathbf{x}^H\tilde{\mathbf{G}}(\rr)^H\tilde{\mathbf{G}}(\rr)\,\mathbf{x}.}
\end{equation}
This holds for any precoding vector $\mathbf{x}$.
\end{theorem}

\begin{corollary}[Single-wave limit]\label{cor:single-wave}
For $N = M = 1$ with $|x_1|^2 = P$, the sum collapses to one
term.  Using $\Lambda_{11} = 1/(2\alpha_1)$, \cref{lem:depth-id},
and $\ReLU(\mu_1) = H(\mu_1)\mu_1$:
\begin{equation}\label{eq:Sab-n-alone}
  \Sab^{(1)} = \frac{\sigma}{4\alpha_1}
  \|\mathbf{F}_1\boldsymbol{\psi}_1\,x_1\|^2
  = \frac{P}{Z_0}\left(T_{s,1}\,|\psi_{s,1}|^2
    + T_{p,1}\,|\psi_{p,1}|^2\right)\ReLU(\mu_1).
\end{equation}
This is exact for TE; the TM contribution has $\sim\!2\,\%$ correction
from \cref{lem:depth-id}. Neither approximation is involved since
self-terms are unaffected. By factoring out the incident power density
$S_{\mathrm{inc}}^{(1)} \equiv \frac{P}{Z_0}|\boldsymbol{\psi}_1|^2$,
this can be rewritten as
\begin{equation}\label{eq:Sab-n-PartII}
  \Sab^{(1)} = S_{\mathrm{inc}}^{(1)}\,\Teff^{(1)}\,\ReLU(\mu_1),
  \qquad \text{where} \quad
  \Teff^{(1)} \equiv T_{s,1}\frac{|\psi_{s,1}|^2}{|\boldsymbol{\psi}_1|^2}
    + T_{p,1}\frac{|\psi_{p,1}|^2}{|\boldsymbol{\psi}_1|^2},
\end{equation}
recovering exactly the polarisation-aware absorption law of Part~II
(\cref{eq:Sab-exact}). Under the additional pseudo-Brewster approximation
$T_{s,1}\approx T_{p,1}\approx T_0$, the effective transmission
collapses to $\Teff^{(1)} \approx T_0$, yielding
\begin{equation}\label{eq:Sab-n-PB}
  \Sab^{(1)} = S_{\mathrm{inc}}^{(1)}\,T_0\,\ReLU(\mu_1),
\end{equation}
which is the fully unpolarised geometric absorption law of Part~I.
\end{corollary}

\begin{corollary}[Incoherent limit]\label{cor:incoherent}
When the $N$ paths have mutually random phases, cross-terms
$n \neq n'$ average to zero.  Only the self-terms survive:
\begin{equation}
  \langle\Sab(\rr)\rangle = \sum_n
  \frac{|x_{j(n)}|^2}{Z_0}
  \left(T_{s,n}\,|\psi_{s,n}|^2 + T_{p,n}\,|\psi_{p,n}|^2\right)
  \ReLU(\mu_n).
\end{equation}
Under the pseudo-Brewster approximation this reduces to the
multi-source formula of Part~I\@.  The coherent law strictly
contains the incoherent result as a special case.
\end{corollary}

Both approximations affect only cross-terms: Approximation~1 because
$|\hat{e}_{p,n}|^2 = |\hat{e}'_{p,n}|^2 = 1$; Approximation~2
because $\Gamma_{nn} = 1$ by definition.  The bounded error
accumulates only in cross-terms.

% ============================================================
\section{Exposure operator}\label{sec:exposure-operator}
% ============================================================

Integrating the local coherent absorption law over the body yields
a global quadratic form in the precoder that encodes the complete
body-absorption response.

Integrating \cref{eq:coherent-law} over the body surface:
\begin{equation}
  P_\mathrm{abs} = \int_\Sigma \Sab(\rr)\,\diff A
  = \mathbf{x}^H \underbrace{\int_\Sigma
    \tilde{\mathbf{G}}(\rr)^H\tilde{\mathbf{G}}(\rr)\,\diff A}_{=\,\mathbf{Q}}\mathbf{x}.
\end{equation}

\begin{definition}[Exposure operator]\label{def:Q}
The \emph{exposure operator} is the Gram integral
\begin{equation}\label{eq:Q-def}
  \boxed{\mathbf{Q} = \int_\Sigma
    \tilde{\mathbf{G}}(\rr)^H\tilde{\mathbf{G}}(\rr)\,\diff A
    \;\in\; \mathbb{C}^{M\times M},
  \qquad P_\mathrm{abs} = \mathbf{x}^H\mathbf{Q}\,\mathbf{x}.}
\end{equation}
\end{definition}

The quadratic form $P_\mathrm{abs} = \mathbf{x}^H\mathbf{Q}\mathbf{x}$
shares its algebraic structure with the SAR matrix introduced by
Hochwald~\textit{et al.}~\cite{Hochwald2014}, who write
$\mathrm{SAR}(\mathbf{x}) = \mathbf{x}^H\mathbf{S}_V\,\mathbf{x}$,
with $\mathbf{S}_V$ assembled from volume integrals of the per-antenna
electric fields over a fixed tissue volume~$V$, calibrated via FDTD.
The object $\mathbf{S}_V$ is therefore a \emph{peak} (spatially local)
SAR quantity: it captures absorption in a specific $V$ rather than over
the whole body.
Ying~\textit{et al.}~\cite{Ying2015} call the same object the ``SAR matrix''
and proved that the capacity-optimal precoder subject to
$\mathbf{x}^H\mathbf{S}_V\mathbf{x}\le\mathrm{SAR}_0$ has the closed form
$\mathbf{x}^{\star}\propto(\lambda\mathbf{S}_V+\nu\mathbf{I})^{-1}\mathbf{h}^*$.
The present $\mathbf{Q}$ differs in two respects: it integrates over the
entire body surface rather than a local volume, yielding whole-body
absorbed power directly; and its entries are derived in closed form
from the propagation paths and Fresnel coefficients, without calibration.
Ebadi-Shahrivar~\textit{et al.}~\cite{Ebadi2019} showed that
worst-case SAR equals $\lambda_{\max}(\mathbf{S}_V)$; the same
identity holds for $\mathbf{Q}$ with respect to worst-case whole-body
absorbed power per unit transmit power.

$\mathbf{Q}$ is Hermitian positive-semidefinite by construction
($\mathbf{Q}^H = \mathbf{Q}$, $\mathbf{x}^H\mathbf{Q}\mathbf{x} \geq 0$)
and has rank $\leq \min(M, 3M_\mathrm{tri})$ where $M_\mathrm{tri}$
is the number of mesh triangles.  Its eigendecomposition
$\mathbf{Q} = \sum_k \lambda_k \mathbf{v}_k\mathbf{v}_k^H$
($\lambda_1 \geq \cdots \geq \lambda_M \geq 0$) reveals the
\emph{exposure modes}: $\mathbf{v}_1$ is the precoder direction that
maximises total body absorption;
$\lambda_1 = \max_{\|\mathbf{x}\|=1} P_\mathrm{abs}$ is the
maximum absorbed power per unit transmit power; and
$\operatorname{tr}(\mathbf{Q}) = \sum_k \lambda_k$ equals the
expected absorbed power under isotropic transmission, scaled by~$P$.

This eigendecomposition plays the same role in Part~III as the
absorption directivity $D(\khat)$ in Part~I: it reduces the
infinite-dimensional space of precoders to a finite spectral
description of the body's response.

% ============================================================
\section{MRT and body absorption}\label{sec:MRT-body-section}
% ============================================================

% -------------------------------------------------------
\subsection{Exposure-signal alignment}\label{sec:rho}

\begin{definition}[Exposure-signal alignment]\label{def:rho}
For MRT with $\mathbf{x} = \sqrt{P}\,\mathbf{h}^*/\|\mathbf{h}\|$,
the total absorbed power is
$P_\mathrm{abs}^\mathrm{MRT} = (P/\|\mathbf{h}\|^2)
\mathbf{h}^T\mathbf{Q}\mathbf{h}^*$.
Define
\begin{equation}\label{eq:rho-def}
  \boxed{\rho = \frac{\mathbf{h}^T\mathbf{Q}\,\mathbf{h}^*}
  {\|\mathbf{h}\|^2 \cdot \lambda_\mathrm{max}(\mathbf{Q})}
  \;\in\; [0,1].}
\end{equation}
\end{definition}

The bound $\rho \in [0,1]$ follows from the Rayleigh--Ritz theorem
applied to the Hermitian PSD matrix $\mathbf{Q}$: setting
$\mathbf{v} = \mathbf{h}^*$ gives $\mathbf{h}^T\mathbf{Q}\mathbf{h}^*
= \mathbf{v}^H\mathbf{Q}\mathbf{v} \leq \lambda_\mathrm{max}(\mathbf{Q})
\|\mathbf{h}\|^2$.  The absorbed power under MRT reads
\begin{equation}\label{eq:Pabs-MRT}
  P_\mathrm{abs}^\mathrm{MRT} = P\cdot\rho\cdot\lambda_\mathrm{max}(\mathbf{Q}).
\end{equation}

The parameter $\rho$ is the single number that quantifies the
exposure cost of MRT.  A large $\rho$ means the dominant signal
paths also illuminate the body at moderate incidence; a small $\rho$
means those paths arrive at grazing incidence or from directions
where the body is self-shadowed.

% -------------------------------------------------------
\subsection{Absorption under MRT precoding}\label{sec:MRT-body}

MRT sets $\mathbf{x} = \sqrt{P}\,\mathbf{h}^*/\|\mathbf{h}\|$, so
the absorbed power density at surface point~$\rr$ is
\begin{equation}\label{eq:Sab-MRT}
  \Sab(\rr) = \frac{P}{\|\mathbf{h}\|^2}\;
  \bigl\|\tilde{\mathbf{G}}(\rr)\,\mathbf{h}^*\bigr\|^2
  = \frac{P}{\|\mathbf{h}\|^2}\;
  \left\|\sum_{j=1}^M h_j^*\;\tilde{\mathbf{g}}_j(\rr)\right\|^2.
\end{equation}
Compare to the received signal, which sums \emph{scalars}:
$y \propto \sum_j h_j^* h_j = \sum_j|h_j|^2 \in \mathbb{R}^+$.
The signal sums $h_j^* h_j$; the body field sums \emph{vectors}
$h_j^*\tilde{\mathbf{g}}_j(\rr)$.  Three structural differences
prevent the body hotspot from being a copy of the signal focus.

First, the signal is a scalar inner product, whereas the body field is a
3D vector.  Even when all magnitudes $|h_j^*\tilde{g}_j|$ are
equal, the vectors may partially cancel when transmitted
polarisations differ at $\mathbf{r}$.
Second, the UE channel projects each path via
$\mathbf{C}_R(\hat{\mathbf{k}}_n)^H\boldsymbol{\psi}_n$, while the
exposure channel applies
$\mathbf{F}_n(\mathbf{r})\boldsymbol{\psi}_n$ instead.  A path
with strong signal contribution may contribute weakly to absorption
if it arrives at grazing incidence.
Finally, the UE sits at $\mathbf{r}_\mathrm{UE}$, but the body surface
$\Sigma$ lies elsewhere.  Moving to a surface point, each path
acquires an additional phase
$e^{-ik_0\hat{\mathbf{k}}_n\cdot(\mathbf{r}-\mathbf{r}_\mathrm{UE})}$,
degrading coherence. 

Hence MRT is generally neither worst-case nor best-case for body
absorption.  The alignment~$\rho$ is large when the dominant signal
paths also illuminate the body surface at moderate incidence angles
(good Fresnel transmission); it is small when those paths arrive at
grazing incidence on the body or from directions where the body is
self-shadowed.  In the latter case the body is geometrically
protected: strong signal quality coexists with low absorption.

% ============================================================
\section{Coherent hotspot formation}\label{sec:coherent-hotspot}
% ============================================================

Coherent superposition creates spatial maxima of $\Sab$ that can
exceed the incoherent average by a factor of~$N$.  Their size and
intensity depend jointly on phase alignment and the Fresnel filter.

% -------------------------------------------------------
\subsection{Definition and formation conditions}\label{sec:hotspot-def}

A \emph{coherent hotspot} is a local maximum of $\Sab(\rr)$ over
the body surface arising from constructive interference.  Recall the structure
of the hotspot from \cref{thm:coherent-law}:
\begin{equation}\label{eq:Sab-hotspot}
  \Sab(\rr) = \left\|\sum_{n}
  \sqrt{\frac{\sigma}{4\alpha_n}}\;x_{j(n)}\;
  \mathbf{F}_n(\rr)\,\boldsymbol{\psi}_n\;
  e^{-ik_0\hat{\mathbf{k}}_n\cdot\rr}\right\|^2.
\end{equation}
On a locally flat patch, all quantities except the phase factors
$e^{-ik_0\hat{\mathbf{k}}_n\cdot\rr}$ are position-independent;
the spatial variation of $\Sab$ is therefore a pure interference
pattern.  At a point $\rr_0$, the sum has large
norm when the contributing vectors satisfy two conditions simultaneously.
First, they must be co-phased, meaning the complex scalars
$x_{j(n)}\,e^{-ik_0\hat{\mathbf{k}}_n\cdot\rr_0}$ have similar
arguments for the dominant paths.
Second, they must be co-polarised, such that the Fresnel-transmitted vectors
$\mathbf{F}_n(\rr_0)\boldsymbol{\psi}_n$ point in similar
directions in $\mathbb{C}^3$.
Both conditions must hold simultaneously.  Phase alignment alone is
insufficient when transmitted polarisations are orthogonal; polarisation
alignment alone is insufficient when phases cancel.

Under MRT, the phases at the UE are aligned by design.  On the body
at $\rr \neq \rr_\mathrm{UE}$, each path acquires an additional phase
$e^{-ik_0\hat{\mathbf{k}}_n\cdot(\rr - \rr_\mathrm{UE})}$. For $N$ co-phased, co-polarised paths of equal amplitude, the
coherent sum gives $\Sab^\mathrm{peak} \propto N^2$, while the
incoherent average (random relative phases) gives only
$\langle\Sab\rangle \propto N$.  The coherent peak exceeds the
incoherent average by a factor of~$N$.  Doubling the number of paths
doubles the incoherent floor but quadruples the coherent peak;
\cref{tab:coherent-scaling} quantifies this.

\begin{table}[!htbp]
\centering
\caption{Coherent vs.\ incoherent peak scaling with $N$ co-phased,
  co-polarised paths.  The coherent peak grows as $N^2$; the
  incoherent average as $N$.  The enhancement factor equals $N$.
  At 32 paths, the coherent hotspot is 32$\times$ more intense than
  the incoherent average.}\label{tab:coherent-scaling}
\begin{tabular}{rrrr}
\toprule
Paths $N$ & Incoherent ($\propto N$) & Coherent peak ($\propto N^2$)
  & Enhancement \\
\midrule
$2$  & $2$  & $4$     & $2\times$  \\
$4$  & $4$  & $16$    & $4\times$  \\
$8$  & $8$  & $64$    & $8\times$  \\
$16$ & $16$ & $256$   & $16\times$ \\
$32$ & $32$ & $1024$  & $32\times$ \\
$64$ & $64$ & $4096$  & $64\times$ \\
\bottomrule
\end{tabular}
\end{table}

% -------------------------------------------------------
\subsection{Hotspot size and angular spread}\label{sec:hotspot-size}

\begin{remark}[Hotspot size as a Fourier uncertainty]
The rate at which constructive interference degrades away from a
hotspot peak is set by the spread of the transverse wave-vector
components.  On a locally flat surface patch, the transverse
wave-vector of path $n$ is
$\mathbf{q}_n = k_0\,\mathrm{proj}_{\hat{\mathbf{n}}^\perp}(\hat{\mathbf{k}}_n)$.
A displacement $\boldsymbol{\delta}$ along the surface shifts the
phase of path $n$ by $e^{-i\mathbf{q}_n\cdot\boldsymbol{\delta}}$.
The effective transverse wave-vector spread, weighted by path
contribution, is
\begin{equation}
  \sigma_q^2 = \frac{\sum_{n,n'} w_n^* w_{n'}\;
  |\mathbf{q}_n - \mathbf{q}_{n'}|^2}
  {\bigl|\sum_n w_n\bigr|^2},
\end{equation}
where $w_n = x_{j(n)}\,\hat{\mathbf{u}}\cdot
\mathbf{F}_n(\rr_0)\boldsymbol{\psi}_n\,
e^{-ik_0\hat{\mathbf{k}}_n\cdot\rr_0}$ is the scalar contribution
of path~$n$ projected onto the dominant polarisation
direction~$\hat{\mathbf{u}}$ at the hotspot.
The coherence length, over which $\Sab$ decays to half its peak,
scales as
\begin{equation}
  \delta_\mathrm{coh} \sim \frac{1}{\sigma_q} \sim \frac{\lambda}{2\pi\sigma_\Omega},
\end{equation}
where $\sigma_\Omega$ is the
corresponding angular spread projected onto the surface tangent
plane.  The hotspot diameter is $D_\mathrm{spot} \sim 2\delta_\mathrm{coh}$.
Wider angular spread gives a tighter hotspot; for near-LOS channels
($\sigma_\Omega \ll 1$\,rad) no compact hotspot forms.  This is a
Fourier uncertainty relation between direction and position.
\end{remark}

% -------------------------------------------------------
\subsection{Fresnel filtering effects}\label{sec:fresnel-filter}

The Fresnel operator $\mathbf{F}_n(\mathbf{r})$ and the
depth-integral weight $\sqrt{\sigma/(4\alpha_n)}$ modify each path's
contribution before it enters the coherent sum, reshaping the hotspot
relative to the free-space intensity pattern through three distinct effects.

First, grazing suppression occurs because Fresnel transmission coefficients
$|t_s|$, $|t_p|$ decrease for paths arriving at large incidence
angle.  Grazing paths therefore contribute less to $\Sab$ than to
$|\mathbf{E}|^2$.  The absorbed hotspot is shaped by the paths the
body actually couples to, not those that produce the strongest
free-space field.

Second, polarisation mixing arises because paths arriving from different
directions decompose differently into TE and TM. Their
Fresnel-transmitted polarisations $\mathbf{F}_n\boldsymbol{\psi}_n$
may be less aligned than the incident polarisations.  This reduces
vector coherence at the hotspot and tends to broaden it.

Finally, for visibility gating, back-facing paths ($\mu_n \leq 0$)
contribute nothing.  On a curved body, different surface points see
different path subsets.  The hotspot can only form from paths that
are front-facing at that location.

\begin{remark}[Free-space peak vs.\ absorbed peak]
The incident power density at the surface is
$S_\mathrm{inc}(\mathbf{r}) = \|\mathbf{G}(\mathbf{r})\mathbf{x}\|^2/Z_0$
(units: $\mathrm{W\,m^{-2}}$); the absorbed power density
is $\Sab(\mathbf{r}) = \|\tilde{\mathbf{G}}(\mathbf{r})\mathbf{x}\|^2$
(units: $\mathrm{W\,m^{-2}}$).
The relation between these two is mediated by the per-path Fresnel
filter, which acts on each path component separately.  The absorbed
peak is generally weaker (Fresnel reflection removes power),
displaced (different paths dominate after Fresnel weighting), and
possibly broader or narrower (depending on which paths survive the
filter).  No general inequality holds between the locations of the
two peaks.
\end{remark}

% ============================================================
\section{Exposure-constrained beamforming}\label{sec:ecbf}
% ============================================================

The operator $\mathbf{Q}$ reduces the absorbed-power constraint to a
quadratic form in the precoder.  The resulting optimisation has a
closed-form solution.

% -------------------------------------------------------
\subsection{Problem formulation}\label{sec:ecbf-problem}

The standard communications objective is to maximise received signal
power at the UE subject to a bound on total body absorption:
\begin{equation}\label{eq:QCQP}
  \max_{\mathbf{x}}\; |\mathbf{h}^T\mathbf{x}|^2
  \quad\text{subject to}\quad
  \mathbf{x}^H\mathbf{Q}\,\mathbf{x} \leq P_\mathrm{abs}^\mathrm{max},
  \quad \|\mathbf{x}\|^2 \leq P.
\end{equation}
Both constraints are convex (the matrices $\mathbf{Q}$ and
$\mathbf{I}$ are PSD).  The objective is a quadratic being maximised,
making \cref{eq:QCQP} a non-convex QCQP.  A global maximum exists
because the feasible set is compact and non-empty for
$P, P_\mathrm{abs}^\mathrm{max} > 0$.  For QCQPs with at most two
quadratic constraints, the S-procedure guarantees that the KKT
point is a global maximum.

% -------------------------------------------------------
\subsection{Optimal precoder}\label{sec:ecbf-solution}

The KKT conditions give:
\begin{equation}\label{eq:optimal-x}
  \boxed{\mathbf{x}^\star = \sqrt{P}\;\frac{(\lambda\mathbf{Q} + \nu\mathbf{I})^{-1}\mathbf{h}^*}
  {\|(\lambda\mathbf{Q} + \nu\mathbf{I})^{-1}\mathbf{h}^*\|},}
\end{equation}
with Lagrange multipliers $\lambda \geq 0$ (exposure) and $\nu \geq
0$ (power) determined by complementary slackness.  The matrix
$\lambda\mathbf{Q} + \nu\mathbf{I}$ is positive definite whenever
$\nu > 0$ (satisfied when the power constraint is active, which is
generically the case).
Ying~\textit{et al.}~\cite{Ying2015} derived the same closed-form structure
for their SAR-matrix formulation, with $\mathbf{Q}$ replaced by the
empirically calibrated SAR matrix~$\mathbf{S}$ and with strong duality
proved by induction.  The present derivation applies the same result with
$\mathbf{Q}$ derived analytically from the propagation geometry and Fresnel
coefficients.

The matrix $(\lambda\mathbf{Q}+\nu\mathbf{I})^{-1}$ regularises MRT
by reshaping the precoder in the $\mathbf{Q}$-eigenspace: along
eigenvectors with large $\lambda_k$ (high-absorption directions), the
precoding is suppressed by $\sim 1/(\lambda\lambda_k)$.  Along null-space
directions (zero absorption), the precoding is unpenalised.

% -------------------------------------------------------
\subsection{Limiting regimes}\label{sec:ecbf-limits}

\paragraph{Unconstrained ($\lambda = 0$).} $\mathbf{x}^\star \propto \mathbf{h}^*$.
Standard MRT.  The exposure constraint is slack: MRT already satisfies
the absorption limit.

\paragraph{Tightly constrained ($\lambda \gg \nu$).}
$\mathbf{x}^\star \approx \mathbf{Q}^{-1}\mathbf{h}^*$ (up to normalisation;
pseudoinverse if $\mathbf{Q}$ is rank-deficient).  The precoder steers
away from the body's absorption modes.  Signal power is reduced, but
absorption is controlled.

\paragraph{Orthogonal regime ($\rho \ll 1$).} The MRT direction
$\mathbf{h}^*$ already has low overlap with the top eigenvectors of
$\mathbf{Q}$.  The exposure constraint is satisfied with minimal
signal loss.  Body geometry provides natural protection.


% ============================================================
\section{Multi-user MIMO}\label{sec:mu-mimo}
% ============================================================

The single-user framework extends to $K$ simultaneous users without
new approximations.

In a multi-user MIMO (MU-MIMO) downlink, the BS transmits to
$K$~users simultaneously.  User~$k$ has data symbol~$s_k$ and
precoding vector $\mathbf{w}_k \in \mathbb{C}^M$; the total
transmitted signal is
$\mathbf{x}_\mathrm{total} = \sum_{k=1}^K \mathbf{w}_k s_k$
with total power constraint $\sum_k\|\mathbf{w}_k\|^2 \leq P$.
The field channel $\mathbf{G}(\mathbf{r})$ and the body-surface
channel $\tilde{\mathbf{G}}^{(u)}(\mathbf{r})$ for each
person~$u$'s body are exactly as defined in
\cref{sec:field-channel,sec:fresnel-operator,sec:coherent-absorption-law}.

% -------------------------------------------------------
\subsection{Per-person absorbed power}\label{sec:mu-per-person}

The absorbed power density on person~$u$'s body at
$\mathbf{r} \in \Sigma^{(u)}$ is
\begin{equation}
  S_\mathrm{ab}^{(u)}(\mathbf{r})
  = \bigl\|\tilde{\mathbf{G}}^{(u)}(\mathbf{r})\,
    \mathbf{x}_\mathrm{total}\bigr\|^2
  = \left\|\sum_{k=1}^K
    \tilde{\mathbf{G}}^{(u)}(\mathbf{r})\,\mathbf{w}_k\,s_k\right\|^2.
\end{equation}
Expanding the squared norm and taking the expectation over the data
symbols ($\mathbb{E}[s_k] = 0$, $\mathbb{E}[|s_k|^2] = 1$,
$\mathbb{E}[s_k s_{k'}^*] = 0$ for $k \neq k'$):

\begin{proposition}[Incoherent stream superposition]\label{prop:mu-superposition}
\begin{equation}\label{eq:Sab-MU-expect}
  \mathbb{E}\bigl[S_\mathrm{ab}^{(u)}(\mathbf{r})\bigr]
  = \sum_{k=1}^K
  \bigl\|\tilde{\mathbf{G}}^{(u)}(\mathbf{r})\,\mathbf{w}_k\bigr\|^2.
\end{equation}
Cross-terms between streams vanish because different users' data
symbols are uncorrelated.  The time-averaged absorbed power density
from $K$ simultaneous streams is the incoherent sum of each stream's
individual contribution.  Each term
$\|\tilde{\mathbf{G}}^{(u)}\mathbf{w}_k\|^2$ is itself a coherent
dosimetry problem in the sense of \cref{thm:coherent-law}.
\end{proposition}

% -------------------------------------------------------
\subsection{Per-person exposure operator}\label{sec:mu-Q}

Integrating over person~$u$'s body surface defines a per-person
exposure operator:
\begin{equation}\label{eq:Q-u}
  \mathbf{Q}^{(u)} = \int_{\Sigma^{(u)}}
  \tilde{\mathbf{G}}^{(u)}(\mathbf{r})^H
  \tilde{\mathbf{G}}^{(u)}(\mathbf{r})\,\diff A
  \;\in\; \mathbb{C}^{M\times M}.
\end{equation}
The expected total absorbed power in person~$u$'s body is
\begin{equation}\label{eq:Pabs-MU}
  \boxed{P_\mathrm{abs}^{(u)}
  = \sum_{k=1}^K \mathbf{w}_k^H\,\mathbf{Q}^{(u)}\,\mathbf{w}_k.}
\end{equation}
Each person has a separate Hermitian PSD matrix $\mathbf{Q}^{(u)}$
built from the paths arriving at their body.  A single stream's
precoder~$\mathbf{w}_k$ contributes to the absorption of every body
it illuminates, not just the intended recipient.

% -------------------------------------------------------
\subsection{Exposure-constrained MU-MIMO design}\label{sec:mu-ecbf}

The exposure-constrained MU-MIMO design problem maximises sum-rate
subject to per-person absorption limits:
\begin{equation}\label{eq:MU-QCQP}
  \max_{\{\mathbf{w}_k\}} \;\sum_{k=1}^K
  \log_2\!\left(1 + \frac{|\mathbf{h}_k^T\mathbf{w}_k|^2}
  {\sum_{k'\neq k}|\mathbf{h}_k^T\mathbf{w}_{k'}|^2 + \sigma_n^2}\right)
\end{equation}
subject to
\begin{equation}
  \sum_{k=1}^K \mathbf{w}_k^H\,\mathbf{Q}^{(u)}\,\mathbf{w}_k
  \leq P_\mathrm{abs}^\mathrm{max},\quad u = 1,\ldots,K,
  \qquad
  \sum_{k=1}^K \|\mathbf{w}_k\|^2 \leq P.
\end{equation}
The per-person exposure constraints are convex in
$\{\mathbf{w}_k\}$ (each a sum of quadratic forms with PSD
matrices), and they integrate directly into WMMSE or
alternating-optimisation frameworks.

\begin{remark}[Self-exposure vs.\ cross-exposure]
Person~$u$'s body is illuminated both by the stream intended for
them ($k = u$) and by streams intended for other users
($k \neq u$).  Zero-forcing precoders, which null inter-user
interference at each UE, simultaneously reduce cross-exposure of
bystander bodies that are near the UE: nulling the signal at a UE
also reduces the field incident on the nearby body.  When bystander
bodies are spatially separated from all UEs, this alignment between
signal nulls and body illumination is not guaranteed.
\end{remark}

\begin{remark}[Channel dynamics]
When users move, both $\mathbf{h}_k$ and $\mathbf{Q}^{(u)}$ evolve
with position and body orientation.  Over a coherence interval~$T_c$
the time-averaged constraint is
\begin{equation}
  \frac{1}{T_c}\int_0^{T_c}
  P_\mathrm{abs}^{(u)}(t)\,\diff t \leq P_\mathrm{abs}^\mathrm{max};
\end{equation}
$\tilde{\mathbf{G}}^{(u)}$ is re-evaluated at each channel
estimation step from the current body pose, requiring no additional
assumptions beyond the inherited locally-flat-surface condition.
\end{remark}

\begin{remark}[Per-user exposure tracking]
Computing $\tilde{\mathbf{G}}^{(u)}(\mathbf{r})$ requires the
per-path polarisation-amplitude vectors $\boldsymbol{\psi}_n$, not
just the scalar channel $\mathbf{h}$.  Site-specific ray tracing
provides these vectors directly.  An alternative is to measure
$\mathbf{Q}^{(u)}$ via a calibrated phantom; the leading eigenvalue
$\lambda_1^{(u)}$ then bounds absorbed power for any precoder:
$P_\mathrm{abs}^{(u)} \leq \lambda_1^{(u)} P$.
\end{remark}

The operators $\{\mathbf{Q}^{(u)}\}_{u=1}^K$ give a finite-dimensional
description of how any precoder distributes absorbed power across all
bodies in the scene.  No assumption beyond those of
\cref{sec:coherent-setup} is required.  The per-person exposure
constraints in \cref{eq:MU-QCQP} are quadratic with
positive-semidefinite matrices; they fit directly into standard WMMSE
or successive-convex-approximation solvers.

\bigskip
\noindent\textbf{Symbol table for Part~III.}
The following table collects all new symbols introduced in Part~III.

\begin{center}
\begin{tabular}{lp{6.5cm}l}
\toprule
Symbol & Meaning & Space / units \\
\midrule
$M$ & number of BS antenna elements & scalar \\
$N$ & number of propagation paths & scalar \\
$\mathbf{x}$ & precoding vector & $\mathbb{C}^M$, $[\sqrt{\mathrm{W}}]$ \\
$\mathbf{C}_{T,j}(\khat)$ & BS transmit antenna pattern, element $j$
  (\cref{sec:antenna-paths}) & $\mathbb{C}^2$ \\
$\mathbf{C}_R(\khat)$ & UE receive antenna pattern
  (\cref{sec:antenna-paths}) & $\mathbb{C}^2$ \\
$\mathbf{T}_n$ & transfer matrix, path $n$
  (\cref{sec:antenna-paths}) & $\mathbb{C}^{2\times 2}$ \\
$\boldsymbol{\psi}_n$ & complex polarisation-amplitude vector, path $n$ &
  $\mathbb{C}^3$, $[\mathrm{V\,m^{-1}\,W^{-1/2}}]$ \\
$\hat{\mathbf{k}}_n$ & direction of arrival, path $n$ & unit vector \\
$\mu_n(\mathbf{r})$ & incidence cosine,
  $\hat{\mathbf{n}}(\mathbf{r})\cdot(-\hat{\mathbf{k}}_n)$ & real \\
$\xi_n$ & normal wave-vector in tissue (\cref{eq:xi-def}) & complex \\
$\alpha_n$, $\beta_n$ & depth-decay and phase rates
  ($k_0\xi_n = \beta_n - i\alpha_n$) & $[\mathrm{m}^{-1}]$ \\
$\mathbf{G}(\mathbf{r})$ & field channel matrix
  (\cref{eq:G-def}) & $\mathbb{C}^{3\times M}$ \\
$\mathbf{h}$ & UE channel vector (\cref{eq:h-def}) & $\mathbb{C}^M$ \\
$\mathbf{F}_n(\mathbf{r})$ & Fresnel transmission operator, path $n$
  (\cref{def:fresnel-op}) & $\mathbb{C}^{3\times 3}$ \\
$\tilde{\mathbf{G}}(\mathbf{r})$ & absorption-normalised body-surface
  channel (\cref{eq:Gtilde-def}) & $\mathbb{C}^{3\times M}$ \\
$\mathbf{Q}$ & exposure operator (\cref{def:Q}) & $\mathbb{C}^{M\times M}$ \\
$\rho$ & exposure-signal alignment (\cref{eq:rho-def}) & $[0,1]$ \\
\bottomrule
\end{tabular}
\end{center}


\appendix
\crefalias{section}{appendix}
\crefalias{subsection}{appendix}
% ============================================================

\section{Complex refractive index}\label{app:refractive-index}

For a lossy dielectric with relative permittivity~$\varepsilon_r$ and
conductivity~$\sigma$, the complex relative permittivity is
\begin{equation}
  \epstilde = \varepsilon_r - i\frac{\sigma}{\omega\varepsilon_0}.
\end{equation}
The complex refractive index is $\ntilde=\sqrt{\epstilde}$, with the
branch chosen so that $\RE(\ntilde)>0$ and $\IM(\ntilde)<0$ (the wave
decays in the $+z$ direction).  Writing $\ntilde=n-i\kappa$ with
$n,\kappa>0$:
\begin{equation}
  n = \frac{1}{\sqrt{2}}\sqrt{|\epstilde|+\varepsilon_r}, \qquad
  \kappa = \frac{1}{\sqrt{2}}\sqrt{|\epstilde|-\varepsilon_r}.
\end{equation}

\Cref{tab:skin-data} lists the dielectric properties of human skin
(Gabriel 1996 four-pole Cole--Cole model~\cite{Gabriel1996}, IT'IS
v5.0 database~\cite{ITISv5}) at representative frequencies across the
mmWave band.  All numerical examples in this work use these
values.  The normal-incidence
power-absorption coefficient~$T_0$ is computed from
\cref{eq:T0}.  As frequency increases, permittivity decreases
and the loss tangent approaches unity, so that
$|\ntilde|$ decreases and $T_0$ increases.

\begin{table}[!htbp]
\centering
\caption{Dielectric properties of human skin (Gabriel 1996 /
  IT'IS v5.0) and derived optical quantities at selected mmWave
  frequencies.}\label{tab:skin-data}
\begin{tabular}{rrrrrrrr}
\toprule
$f$ (GHz) & $\varepsilon_r$ & $\sigma$ (S/m) &
  $n$ & $\kappa$ & $|\ntilde|$ & $T_0$ \\
\midrule
6   & 34.95 & \phantom{0}3.9 & 5.991 & 0.973 & 6.07 & 0.481 \\
10  & 31.29 & \phantom{0}8.0 & 5.733 & 1.256 & 5.87 & 0.489 \\
20  & 21.96 &           19.2 & 4.995 & 1.729 & 5.29 & 0.513 \\
28  & 16.55 &           25.8 & 4.471 & 1.854 & 4.84 & 0.536 \\
40  & 11.69 &           31.8 & 3.882 & 1.840 & 4.30 & 0.570 \\
60  & \phantom{0}7.98 & 36.4 & 3.278 & 1.663 & 3.68 & 0.622 \\
80  & \phantom{0}6.40 & 38.4 & 2.927 & 1.474 & 3.28 & 0.665 \\
100 & \phantom{0}5.60 & 39.4 & 2.705 & 1.310 & 3.01 & 0.701 \\
\bottomrule
\end{tabular}
\end{table}

\section{Derivation of Fresnel coefficients from boundary conditions}
\label{app:fresnel-derivation}

This appendix derives the Fresnel reflection coefficients stated in
\cref{sec:fresnel} from the electromagnetic boundary conditions at a
planar interface.

% -------------------------------------------------------
\subsection{Setup}
\label{app:fresnel-setup}

Consider a planar interface at $z=0$ between free space ($z>0$) and a
lossy medium with complex permittivity
$\tilde{\varepsilon}=\varepsilon_r-i\sigma/\omega\varepsilon_0$ and
refractive index $\ntilde=\sqrt{\tilde{\varepsilon}}$ ($z<0$).  The
standard boundary conditions require continuity of the tangential
components of $\EE$ and $\HH$ across the interface.

A plane wave in a medium with refractive index $\ntilde_j$ satisfies
the dispersion relation $k=k_0\ntilde_j$ (with
$k_0=\omega/c$), and its magnetic field is
\begin{equation}\label{eq:app-H-from-E}
  \HH = \frac{\ntilde_j}{Z_0}\,\khat\times\EE,
\end{equation}
where $Z_0=\mu_0 c\approx 377\,\Omega$.

% -------------------------------------------------------
\subsection{Snell's law from phase matching}
\label{app:snell}

An incident wave propagating in the $xz$-plane makes angle~$\theta_i$
with the surface normal~$\hat{z}$, so
\begin{equation}
  \mathbf{k}_i = k_0(\sin\theta_i\,\hat{x} - \cos\theta_i\,\hat{z}).
\end{equation}
The reflected wave has $\mathbf{k}_r = k_0(\sin\theta_r\,\hat{x}
+\cos\theta_r\,\hat{z})$ and the transmitted wave has
$\mathbf{k}_t = k_0\ntilde(\sin\theta_t\,\hat{x}
-\cos\theta_t\,\hat{z})$.

The boundary conditions must hold at every point on the interface
($z=0$, all~$x$).  This requires
all three waves to have the same tangential wave vector, i.e.
\begin{equation}
  k_0\sin\theta_i = k_0\sin\theta_r = k_0\ntilde\sin\theta_t.
\end{equation}
The first equality gives $\theta_r=\theta_i$ (law of reflection).
The second gives Snell's law, generalised to complex media:
\begin{equation}\label{eq:app-snell}
  \boxed{\sin\theta_t = \frac{\sin\theta_i}{\ntilde}.}
\end{equation}
For a lossy medium $\ntilde$ is complex and so is $\theta_t$.  The
tangential phase velocity is conserved across the interface;
$\theta_t$ does not correspond to a geometric angle.

The normal component of the transmitted wave vector is
\begin{equation}\label{eq:app-ktz}
  k_{t,z} = k_0\ntilde\cos\theta_t
  = k_0\sqrt{\ntilde^2 - \sin^2\!\theta_i}
  \;\equiv\; k_0\,\xi,
\end{equation}
defining $\xi=\sqrt{\ntilde^2-\sin^2\!\theta_i}$ with
$\RE(\xi)>0$ (the wave decays into the lossy medium, not away from
it).  Using $\cos\theta_i=\mu$:
\begin{equation}
  \xi = \sqrt{\ntilde^2 - 1 + \mu^2},
\end{equation}
which is \cref{eq:xi-def} of the main text.

% -------------------------------------------------------
\subsection{TE (s) polarisation}
\label{app:fresnel-TE}

In TE polarisation the electric field is perpendicular to the plane
of incidence: $\EE=E\,\hat{y}$.  The three waves at the interface
($z=0$) are
\begin{alignat}{2}
  \EE_i &= E_i\,\hat{y}\,\exp[-ik_0(x\sin\theta_i - z\cos\theta_i)],
  &\qquad
  \HH_i &= \frac{E_i}{Z_0}(-\cos\theta_i\,\hat{x} - \sin\theta_i\,\hat{z})\,e^{-i\mathbf{k}_i\cdot\rr},
  \notag\\
  \EE_r &= E_r\,\hat{y}\,\exp[-ik_0(x\sin\theta_i + z\cos\theta_i)],
  &\qquad
  \HH_r &= \frac{E_r}{Z_0}(\cos\theta_i\,\hat{x} - \sin\theta_i\,\hat{z})\,e^{-i\mathbf{k}_r\cdot\rr},
  \notag\\
  \EE_t &= E_t\,\hat{y}\,\exp[-ik_0(x\sin\theta_i - z\xi)],
  &\qquad
  \HH_t &= \frac{E_t}{Z_0}(-\xi\,\hat{x} - \sin\theta_i\,\hat{z})\,e^{-i\mathbf{k}_t\cdot\rr},
\end{alignat}
where the magnetic fields follow from \cref{eq:app-H-from-E} and
Snell's law has been used to eliminate $\theta_t$ in favour of~$\xi$.

\paragraph{Tangential-$E$ condition} ($y$-component at $z=0$):
\begin{equation}\label{eq:app-TE-Ey}
  E_i + E_r = E_t.
\end{equation}

\paragraph{Tangential-$H$ condition} ($x$-component at $z=0$):
\begin{equation}\label{eq:app-TE-Hx}
  -\frac{E_i}{Z_0}\cos\theta_i + \frac{E_r}{Z_0}\cos\theta_i
  = -\frac{E_t}{Z_0}\,\xi.
\end{equation}
Multiplying through by $Z_0$ and writing $\mu=\cos\theta_i$:
\begin{equation}
  -E_i\,\mu + E_r\,\mu = -E_t\,\xi.
\end{equation}

Eliminating $E_t$ using \cref{eq:app-TE-Ey}:
\begin{equation}
  -E_i\,\mu + E_r\,\mu = -(E_i+E_r)\,\xi
  \quad\Longrightarrow\quad
  E_r(\mu+\xi) = E_i(\mu-\xi),
\end{equation}
and therefore
\begin{equation}\label{eq:app-rs-result}
  \boxed{r_s \;\equiv\; \frac{E_r}{E_i}
    \;=\; \frac{\mu - \xi}{\mu + \xi},}
\end{equation}
which is \cref{eq:rs} of the main text.

% -------------------------------------------------------
\subsection{TM (p) polarisation}
\label{app:fresnel-TM}

In TM polarisation the magnetic field is perpendicular to the plane
of incidence: $\HH=H\,\hat{y}$.  The electric fields lie in the
$xz$-plane.  From \cref{eq:app-H-from-E}, $\EE=(Z_0/\ntilde_j)\,
\HH\times\khat_j$ in medium~$j$.  Explicitly:
\begin{alignat}{2}
  \HH_i &= H_i\,\hat{y}\,\exp(-i\mathbf{k}_i\cdot\rr),
  &\qquad
  \EE_i &= \frac{H_i Z_0}{1}(\cos\theta_i\,\hat{x} + \sin\theta_i\,\hat{z})\,\exp(-i\mathbf{k}_i\cdot\rr),
  \notag\\
  \HH_r &= H_r\,\hat{y}\,\exp(-i\mathbf{k}_r\cdot\rr),
  &\qquad
  \EE_r &= \frac{H_r Z_0}{1}(-\cos\theta_i\,\hat{x} + \sin\theta_i\,\hat{z})\,\exp(-i\mathbf{k}_r\cdot\rr),
  \notag\\
  \HH_t &= H_t\,\hat{y}\,\exp(-i\mathbf{k}_t\cdot\rr),
  &\qquad
  \EE_t &= \frac{H_t Z_0}{\ntilde^2}(\xi\,\hat{x} + \sin\theta_i\,\hat{z})\,\exp(-i\mathbf{k}_t\cdot\rr),
\end{alignat}
where the factor $1/\ntilde^2$ arises because the wave impedance in
the lossy medium is $Z_0/\ntilde$ and
$\EE=(1/\tilde{\varepsilon}\varepsilon_0)\,\mathbf{D}$.

\paragraph{Tangential-$H$ condition} ($y$-component at $z=0$):
\begin{equation}\label{eq:app-TM-Hy}
  H_i + H_r = H_t.
\end{equation}

\paragraph{Tangential-$E$ condition} ($x$-component at $z=0$):
\begin{equation}\label{eq:app-TM-Ex}
  H_i Z_0\,\mu - H_r Z_0\,\mu
  = H_t\,\frac{Z_0\,\xi}{\ntilde^2}.
\end{equation}
Cancelling $Z_0$ and eliminating $H_t$ using \cref{eq:app-TM-Hy}:
\begin{equation}
  (H_i - H_r)\,\mu = (H_i + H_r)\,\frac{\xi}{\ntilde^2}
  \quad\Longrightarrow\quad
  H_r\!\left(\mu + \frac{\xi}{\ntilde^2}\right)
  = H_i\!\left(\mu - \frac{\xi}{\ntilde^2}\right).
\end{equation}
Multiplying numerator and denominator by $\ntilde^2$:
\begin{equation}\label{eq:app-rp-result}
  \boxed{r_p \;\equiv\; \frac{H_r}{H_i}
    \;=\; \frac{\ntilde^2\,\mu - \xi}{\ntilde^2\,\mu + \xi},}
\end{equation}
which is \cref{eq:rp} of the main text.

\begin{remark}[Sign convention]
The sign of $r_p$ depends on the choice of reference directions for
$\EE_i$ and $\EE_r$.  Here and in the main text, $r_p$ is defined
through the $H$-field ratio, giving $r_p>0$ at normal incidence.  Some references define
$r_p$ through the $E$-field ratio with opposite sign.  The power
reflectance $|r_p|^2$ is independent of convention.
\end{remark}

% -------------------------------------------------------
\subsection{Power-absorption coefficient}
\label{app:power-T}

The time-averaged Poynting vector of a plane wave is
$\langle\mathbf{S}\rangle=\frac{1}{2}\RE(\EE\times\HH^*)$.  For the
incident wave, the $z$-component is
$S_i^z = \frac{|E_i|^2}{2Z_0}\cos\theta_i$.  For the reflected wave,
$S_r^z = \frac{|E_r|^2}{2Z_0}\cos\theta_i$ (outward).

The power reflectance is
\begin{equation}
  \mathcal{R} = \frac{S_r^z}{S_i^z}
  = \frac{|E_r|^2}{|E_i|^2} = |r|^2.
\end{equation}

For a lossy half-space in which the skin depth is much smaller than
any body dimension, all transmitted power is absorbed within a thin
surface layer.  Energy conservation at the interface then gives
\begin{equation}
  \text{(absorbed power)} = \text{(incident power)}
  - \text{(reflected power)},
\end{equation}
and the power-absorption coefficient is
\begin{equation}\label{eq:app-T-result}
  \boxed{T(\theta_i) = 1 - |r(\theta_i)|^2,}
\end{equation}
for each polarisation separately ($T_s=1-|r_s|^2$,
$T_p=1-|r_p|^2$).  This is \cref{eq:T-def} of the main text.

At normal incidence ($\mu=1$, $\xi=\ntilde$), both polarisations
give
\begin{equation}
  r_s(0)=r_p(0) = \frac{1-\ntilde}{1+\ntilde},
  \qquad
  T_0 = 1 - \left|\frac{1-\ntilde}{1+\ntilde}\right|^2
  = \frac{4\,\RE(\ntilde)}{|1+\ntilde|^2},
\end{equation}
recovering \cref{eq:T0}.  Expanding:
\begin{align}
  1-\left|\frac{1-\ntilde}{1+\ntilde}\right|^2
  &= \frac{(1+n)^2+\kappa^2 - (1-n)^2 - \kappa^2}{|1+\ntilde|^2} \notag\\
  &= \frac{4n}{|1+\ntilde|^2},
\end{align}
where $\ntilde=n-i\kappa$.

\section{Pseudo-Brewster angle}\label{app:pseudo-brewster}

For a lossless dielectric, the Brewster angle satisfies
$\tan\theta_B=n_2/n_1$ and $r_p=0$~\cite{BornWolf1999}.  For a lossy
dielectric $|r_p|$ never vanishes but has a
minimum~\cite{Potter1970,Ohman1977}.  The pseudo-Brewster angle where
$T_p$ is maximised is well approximated by
\begin{equation}
  \theta_{pB} \approx \arctan|\ntilde|.
\end{equation}
At this angle: $\cos\theta_{pB}\approx 1/|\ntilde|$,
$T_s(\theta_{pB})\approx 0$, $T_p(\theta_{pB})\approx 1$, and
$\Tavg(\theta_{pB})\approx 0.5$.  The near-coincidence
$\Tavg(\theta_{pB})\approx T_0\approx 0.5$ to $0.6$ for
$|\ntilde|\approx 4$ to $6$ is what makes the constant-$T_0$
approximation work for biological tissue at mmWave frequencies.

\section{Proofs}\label{app:proofs}

\subsection{Proof of \texorpdfstring{\cref{eq:q-exact}}{Equation \ref*{eq:q-exact}}}\label{app:q-proof}

Let the incident wave have Jones vector
$\mathbf{E}_0=E_0(\cos\psi_0\cos\chi-i\sin\psi_0\sin\chi,\;
\sin\psi_0\cos\chi+i\cos\psi_0\sin\chi)^\top$ in the reference
basis~$(\hat{e}_1,\hat{e}_2)$.  The local TE/TM basis is rotated by
angle~$\alpha(\rr)$:
$\hat{e}_s=\cos\alpha\,\hat{e}_1+\sin\alpha\,\hat{e}_2$,
$\hat{e}_p=-\sin\alpha\,\hat{e}_1+\cos\alpha\,\hat{e}_2$.  The
projections are $a_s=E_1\cos\alpha+E_2\sin\alpha$ and
$a_p=-E_1\sin\alpha+E_2\cos\alpha$.  Expanding $|a_p|^2-|a_s|^2$
and substituting the Jones components:
\begin{align}
  |a_p|^2-|a_s|^2
  &= (|E_1|^2-|E_2|^2)(\sin^2\!\alpha-\cos^2\!\alpha)
    - 4\RE(E_1E_2^*)\cos\alpha\sin\alpha \notag\\
  &= -(|E_1|^2-|E_2|^2)\cos 2\alpha - 2\RE(E_1E_2^*)\sin 2\alpha.
\end{align}
From the Jones vector: $|E_1|^2-|E_2|^2=E_0^2\cos 2\chi\cos 2\psi_0$
and $\RE(E_1E_2^*)=\frac{E_0^2}{2}\cos 2\chi\sin 2\psi_0$.  Dividing
by $|E_0|^2=E_0^2$:
\begin{align}
  q &= -\cos 2\chi\cos 2\psi_0\cos 2\alpha
     - \cos 2\chi\sin 2\psi_0\sin 2\alpha \notag\\
    &= -\cos(2\chi)\cos(2\alpha-2\psi_0). \qquad\blacksquare
\end{align}

\subsection{Polarisation independence on a sphere}
\label{app:sphere-proof}

\begin{proposition}
For a sphere illuminated by a single plane wave of arbitrary
polarisation, the total absorbed power is polarisation-independent.
\end{proposition}

\begin{proof}
The sphere has rotational symmetry around the propagation axis $\khat$.
Under a rotation by angle~$\delta$ around~$\khat$, each surface point
maps to another point with the same incidence angle~$\theta$ (hence
the same $\Tavg$, $\Delta T$, and $\cos\theta$), but with
$\alpha\to\alpha+\delta$.  The surface measure~$\diff A$ is invariant.
Therefore
\[
  \int_\Sigma\cos(2\alpha)\,\Delta T(\theta)\cos\theta\,\diff A
  = \int_\Sigma\cos(2\alpha+2\delta)\,\Delta T(\theta)\cos\theta\,\diff A
\]
for all~$\delta$.  Since $\cos(2\alpha+2\delta)$ oscillates
with~$\delta$ and the equation must hold for all~$\delta$, the
integral vanishes.  The same argument applies to the $\sin(2\alpha)$
integral.  Hence $\tilde{B}=0$ and the polarisation correction to
total power vanishes.
\end{proof}

\subsection{Peak \texorpdfstring{$\Sab$}{S\_ab} is polarisation-universal}
\label{app:peak-proof}

The local absorbed power density is
$\Sab=\Sinc\,\Teff(\theta)\cos\theta$.  At $\theta=0$,
$\Delta T=0$, so $\Teff(0)=T_0$ for any polarisation and
$\Sab(0)=\Sinc\,T_0$.

For TE polarisation, $T_s(\theta)$ is monotonically decreasing (since
$|r_s|^2$ is monotonically increasing), so the product
$T_s(\theta)\cos\theta$ is the product of two non-negative decreasing
functions of~$\theta$ and is therefore monotonically decreasing.
Hence $T_s(\theta)\cos\theta \le T_0$ for all~$\theta$.

For TM polarisation, $T_p(\theta)$ can increase towards the
pseudo-Brewster peak, but the flux factor $\cos\theta$ decreases
faster.  Numerical verification across all lossy dielectrics with
$\RE(n)\in[1,100]$ and $\kappa\in[0,50]$ (including all biological
tissues from 0.9 to 100\,GHz) confirms that
$T_p(\theta)\cos\theta$ is monotonically decreasing in~$\theta$.

Combining both cases:
$\max_{\theta,\mathrm{pol}}\Sab = \Sinc\,T_0$, with equality at
normal incidence.  The bound is exact, not approximate (see also
\cref{sec:exact-bounds}).

\subsection{TE one-sided bound}\label{app:TE-bound}

For lossy dielectrics with $|\ntilde|>1$, $T_s(\theta)$ is
monotonically decreasing (because $|r_s|^2$ is monotonically
increasing).  Therefore $T_s(\theta)\le T_0$ for all~$\theta$, and
for a wave that is locally pure TE:
\[
  P_{\mathrm{abs,TE}}
  = \Sinc\int T_s(\theta)\cos\theta\,\diff A
  \le \Sinc\,T_0\int\cos\theta\,\diff A
  = \Sinc\,T_0\,\Aperp(\khat).
\]
This is an exact upper bound.

% ============================================================
\section{Depth-integral identity}\label{app:depth-identity}
% ============================================================

This appendix proves \cref{lem:depth-id}: the identity
$\frac{\sigma}{4\alpha_n}|t_{s,n}|^2 = \frac{T_{s,n}\cos\theta_n}{Z_0}$
for TE polarisation.

For a TE-polarised transmitted wave in the tissue, energy flux in the
$z$-direction (normal into the body) at the surface ($z = 0$) is
\begin{equation}
  S_z^{\mathrm{trans}} =
  \frac{\operatorname{Re}(\xi_n)}{Z_0\,|\tilde{n}|^2}\,|t_{s,n}|^2
  |\psi_{s,n}|^2.
\end{equation}
Meanwhile, the power absorbed per unit surface area from this wave is
\begin{equation}
  \Pi_n = \frac{\sigma}{2}\int_0^\infty |t_{s,n}|^2|\psi_{s,n}|^2
  e^{-2\alpha_n z}\,\diff z
  = \frac{\sigma}{4\alpha_n}|t_{s,n}|^2|\psi_{s,n}|^2.
\end{equation}
By conservation of energy (Poynting's theorem, no backward wave in
an absorbing medium), $\Pi_n = S_z^{\mathrm{trans}}$.  Hence
\begin{equation}
  \frac{\sigma}{4\alpha_n}|t_{s,n}|^2 = \frac{\operatorname{Re}(\xi_n)}{Z_0|\tilde{n}|^2}.
\end{equation}
The power transmittance is $T_{s,n} = 1 - |r_{s,n}|^2$; by the
same energy argument applied at the interface,
$T_{s,n}\cos\theta_n =
\operatorname{Re}(\xi_n)/|\tilde{n}|^2$
(see \cref{app:fresnel-derivation}).  This gives
\begin{equation}
  \frac{\sigma}{4\alpha_n}|t_{s,n}|^2
  = \frac{T_{s,n}\cos\theta_n}{Z_0},
\end{equation}
completing the proof.

\subsection*{TM polarisation}

For TM, the transmitted field has both a tangential and a normal
component.  From \cref{app:fresnel-TM}, at depth~$z$:
\begin{equation}\label{eq:app-Et-TM}
  \mathbf{E}_t(z)
  = \frac{H_t Z_0}{\tilde{n}^2}
    \bigl(\xi\,\hat{x} + \sin\theta_i\,\hat{z}\bigr)\,
    e^{-ik_0\xi z},
\end{equation}
where $\hat{x}$ is the tangential direction in the incidence plane
and $\hat{z}$ the inward normal.  The absorbed power per unit area is
\begin{equation}
  \Pi_n^{(\mathrm{TM})}
  = \frac{\sigma}{2}\int_0^\infty |\mathbf{E}_t(z)|^2\,\diff z
  = \frac{\sigma}{4\alpha_n}\,
    \frac{|H_t|^2 Z_0^2}{|\tilde{n}|^4}
    \bigl(|\xi|^2 + \sin^2\!\theta_i\bigr).
\end{equation}
The tangential component of $\mathbf{E}_t$ at $z=0$ has magnitude
$|t_p||\psi_{p,n}|$, so the full surface field satisfies
\begin{equation}
  |\mathbf{E}_t(0)|^2
  = |t_p|^2|\psi_{p,n}|^2
    \cdot\frac{|\xi|^2 + \sin^2\!\theta_i}{|\xi|^2}.
\end{equation}
The ratio $(|\xi|^2+\sin^2\!\theta_i)/|\xi|^2$ is the contribution
of the normal component $E_\perp \propto \sin\theta_i/\tilde{n}^2$,
absent in TE.  The absorbed power is therefore
\begin{equation}\label{eq:app-depth-TM-exact}
  \Pi_n^{(\mathrm{TM})}
  = \frac{\sigma}{4\alpha_n}\,|t_{p,n}|^2\,|\psi_{p,n}|^2
    \cdot\underbrace{
      \frac{|\xi_n|^2 + \sin^2\!\theta_n}{|\xi_n|^2}
    }_{\displaystyle\equiv\;\mathcal{C}_p(\theta_n)}.
\end{equation}
The TE-style identity holds when $\mathcal{C}_p = 1$.  Since
$\mathcal{C}_p = 1 + \sin^2\!\theta_n/|\xi_n|^2$ and
$|\xi_n|^2 \geq |\tilde{n}|^2 - 1$:
\begin{equation}
  \mathcal{C}_p - 1 = \frac{\sin^2\!\theta_n}{|\xi_n|^2}
  \leq \frac{1}{|\tilde{n}|^2 - 1}.
\end{equation}
For skin at 28\,GHz ($|\tilde{n}|\approx 4.83$):
$\mathcal{C}_p - 1 \leq 4.5\,\%$ at $\theta_n = 90^\circ$ and
$\leq 2\,\%$ for $\theta_n \leq 60^\circ$.  The correction has the
same origin as Approximation~1 (\cref{prop:approx1}): the normal
electric field inside tissue is
$E_\perp \propto \sin\theta_t = O(1/|\tilde{n}|)$, so its squared
contribution to $|\mathbf{E}^\mathrm{trans}|^2$ is
$O(1/|\tilde{n}|^2)$.

% ============================================================
\section{Approximation accuracy}\label{app:depth-coupling}
% ============================================================

\Cref{tab:approx-accuracy} quantifies the errors of
Approximations~1 and~2 for skin at representative millimetre-wave
frequencies.  All values are derived from numerical evaluation of
the exact Fresnel coefficients using the IT'IS v5.0
database~\cite{ITISv5}.

\begin{table}[!htbp]
\centering
\caption{Error bounds for Approximations~1 and~2 for biological
  skin.  Approximation~1 error (TM polarisation direction
  replacement) is $\max_{\theta\in[0^\circ,85^\circ]}
  |\hat{e}'_{p,n} - \hat{e}_{p,n}|$ on cross-terms, as a fraction.
  Approximation~2 error is $\max_{\theta,\theta'}
  |1 - \Gamma_{nn'}|$ over all angle pairs.
  $|\tilde{n}|$ from \\cref{tab:skin-data}.}%
\label{tab:approx-accuracy}
\begin{tabular}{rrrrrr}
\toprule
$f$ (GHz) & $|\tilde{n}|$ &
  App.~1 error (\%) & App.~2 error (\%) \\
\midrule
10  & 5.87 & $\lesssim 3$   & $\leq 0.30$ \\
28  & 4.83 & $\lesssim 4$   & $\leq 0.44$ \\
40  & 4.30 & $\lesssim 5$   & $\leq 0.55$ \\
60  & 3.68 & $\lesssim 7$   & $\leq 0.75$ \\
100 & 3.01 & $\lesssim 11$  & $\leq 1.1$  \\
\bottomrule
\end{tabular}
\end{table}

Approximation~1 error increases with frequency because $|\tilde{n}|$
decreases, enlarging the refracted angle $\theta_t = \arcsin(\sin\theta/|\tilde{n}|)$.
Approximation~2 error follows the same trend.  Both remain within
the $\sim\!20\,\%$ tissue dielectric uncertainty across the full
frequency range.
\hl{todo: the App 2 column of the table above and this trend sentence
are wrong (same alpha-beta swap). Correct App 2 max errors: 3.3
percent at 10 GHz, 2.5 at 28, 2.8 at 40, 3.6 at 60, 5.5 at 100. The
correct error is U-shaped with a minimum near 20 to 28 GHz, not
monotone increasing: toward low frequency the loss tangent drops and
the 1/(2 alpha) normalisation amplifies the phase spread. The App 1
column carries the unsound bound, see the todo at prop approx1.}

% ============================================================
\section{Computational cost of the coherent pipeline}\label{app:coherent-cost}
% ============================================================

\Cref{tab:coherent-cost} lists the cost of each step of the
coherent dosimetry pipeline (\cref{sec:field-channel}--\cref{sec:ecbf})
for a single-user scenario with $M_\mathrm{tri}$ mesh triangles,
$N$~paths, and $M$~antenna elements.

\begin{table}[!htbp]
\centering
\caption{Computational cost of the coherent dosimetry pipeline
  (single user).  The dominant cost is the exposure-operator assembly
  at $O(M_\mathrm{tri}\!\cdot\! M^2)$.  For $M_\mathrm{tri}=10^4$,
  $N=10^2$, $M=64$: approximately $4\times 10^7$ operations,
  milliseconds on a GPU.  FDTD at 28\,GHz on an anatomical phantom
  requires $\sim\!10^{12}$ cells, giving $\sim\!10^6\times$
  speed-up.}\label{tab:coherent-cost}
\begin{tabular}{lll}
\toprule
Step & Operation & Cost \\
\midrule
Fresnel parameters & $\mu_n$, $\xi_n$, $t_{s,n}$, $t_{p,n}$,
  $\alpha_n$, TE/TM basis & $O(M_\mathrm{tri}\!\cdot\! N)$ \\
Exposure channel $\tilde{\mathbf{G}}(\mathbf{r}_m)$ &
  sum $N$ Fresnel-weighted paths per triangle &
  $O(M_\mathrm{tri}\!\cdot\! N)$ \\
Local $S_\mathrm{ab}(\mathbf{r}_m)$ &
  $\|\tilde{\mathbf{G}}_m\mathbf{x}\|^2$ per triangle &
  $O(M_\mathrm{tri}\!\cdot\! M)$ \\
Exposure operator $\mathbf{Q}$ &
  sum $M_\mathrm{tri}$ rank-$\leq$3 outer products &
  $O(M_\mathrm{tri}\!\cdot\! M^2)$ \\
Optimal precoder &
  $(\lambda\mathbf{Q}+\nu\mathbf{I})^{-1}\mathbf{h}^*$ &
  $O(M^3)$ \\
\bottomrule
\end{tabular}
\end{table}

The speed-up over FDTD is the same order as that of the incoherent
framework (\cref{sec:computation}), but the coherent pipeline
captures the hotspot structure that FDTD resolves and the incoherent
formula cannot.

% ============================================================
% References
% ============================================================
\begingroup
\raggedright
\begin{thebibliography}{99}

% [1]
\bibitem{ICNIRP2020}
  {ICNIRP}, ``Guidelines for limiting exposure to electromagnetic fields
  (100\,kHz to 300\,GHz),'' \textit{Health Phys.}, vol.~118, no.~5,
  pp.~483--524, 2020.
  \doi{10.1097/HP.0000000000001210}

% [2]
\bibitem{Dimbylow2002}
  P.~J.~Dimbylow, ``Fine resolution calculations of SAR in the human body
  for frequencies up to 3\,GHz,'' \textit{Phys.\ Med.\ Biol.}, vol.~47,
  no.~16, pp.~2835--2846, 2002.
  \doi{10.1088/0031-9155/47/16/301}

% [3]
\bibitem{Hirata2021}
  A.~Hirata, S.~Kodera, K.~Sasaki, J.~Gomez-Tames, I.~Laakso, A.~Wood,
  S.~Watanabe, and K.~R.~Foster, ``Human exposure to radiofrequency energy
  above 6\,GHz: Review of computational dosimetry studies,''
  \textit{Phys.\ Med.\ Biol.}, vol.~66, no.~8, art.~no.~08TR01, 2021.
  \doi{10.1088/1361-6560/abf1b7}

% [4]
\bibitem{Cauchy1841}
  A.-L.~Cauchy, ``M\'{e}moire sur la rectification des courbes et la
  quadrature des surfaces courbes,'' \textit{Comptes Rendus Acad.\ Sci.},
  vol.~13, pp.~1129--1146, 1841.

% [5]
\bibitem{Hadwiger1957}
  H.~Hadwiger, \textit{Vorlesungen \"{u}ber Inhalt, Oberfl\"{a}che und
  Isoperimetrie}, ser.\ Grundlehren der mathematischen Wissenschaften,
  vol.~93. Berlin: Springer-Verlag, 1957.

% [6]
\bibitem{Zhukov1998}
  S.~Zhukov, A.~Iones, and G.~Kronin, ``An ambient light illumination
  model,'' in \textit{Rendering Techniques~'98} (Proc.\ Eurographics
  Workshop on Rendering), G.~Drettakis and N.~Max, Eds. Springer, 1998,
  pp.~45--55.

% [7]
\bibitem{NairHinton2010}
  V.~Nair and G.~E.~Hinton, ``Rectified linear units improve restricted
  Boltzmann machines,'' in \textit{Proc.\ 27th Int.\ Conf.\ Machine
  Learning (ICML)}, 2010, pp.~807--814.

% [8]
\bibitem{SionnaRT}
  J.~Hoydis, F.~A\"{i}t~Aoudia, S.~Cammerer, M.~Nimier-David, N.~Binder,
  G.~Marcus, and A.~Keller, ``{Sionna RT}: Differentiable ray tracing for
  radio propagation modeling,'' \textit{arXiv:2303.11103}, 2023.
  \doi{10.48550/arXiv.2303.11103}

% [9]
\bibitem{Christ2006}
  A.~Christ, T.~Samaras, A.~Klingenb\"{o}ck, and N.~Kuster,
  ``Characterization of the electromagnetic near-field absorption in layered
  biological tissue in the frequency range from 30\,MHz to 6000\,MHz,''
  \textit{Phys.\ Med.\ Biol.}, vol.~51, no.~19, pp.~4951--4965, 2006.
  \doi{10.1088/0031-9155/51/19/014}

% [10]
\bibitem{Hirata2007corr}
  A.~Hirata, Y.~Nagaya, O.~Fujiwara, T.~Nagaoka, and S.~Watanabe,
  ``Correlation between absorption cross section and body surface area of
  human for far-field exposure at GHz bands,'' in \textit{Proc.\ IEEE Int.\
  Symp.\ Electromagn.\ Compat.\ (ISEMC)}, 2007, pp.~1--4.
  \doi{10.1109/ISEMC.2007.143}

% [11]
\bibitem{Zhadobov2011}
  M.~Zhadobov, N.~Chahat, R.~Sauleau, C.~Le~Qu\'{e}ment, and Y.~Le~Dr\'{e}an,
  ``Millimeter-wave interactions with the human body: State of knowledge and
  recent advances,'' \textit{Int.\ J.\ Microw.\ Wireless Technol.}, vol.~3,
  no.~2, pp.~237--247, 2011.
  \doi{10.1017/S1759078711000122}

% [12]
\bibitem{Colombi2018}
  D.~Colombi, B.~Thors, C.~Torn\'{e}vik, and Q.~Balzano, ``RF energy
  absorption by biological tissues in close proximity to millimeter-wave 5G
  wireless equipment,'' \textit{IEEE Access}, vol.~6, pp.~4974--4981, 2018.
  \doi{10.1109/ACCESS.2018.2790038}

% [13]
\bibitem{Gosselin2011}
  M.-C.~Gosselin, G.~Vermeeren, S.~Kuhn, V.~Kellerman, S.~Benkler,
  T.~M.~I.~Uusitupa, W.~Joseph, A.~Gati, J.~Wiart, F.~J.~C.~Meyer,
  \textit{et al.}, ``Estimation formulas for the specific absorption rate in
  humans exposed to base-station antennas,'' \textit{IEEE Trans.\ Electromagn.\
  Compat.}, vol.~53, no.~4, pp.~909--922, 2011.
  \doi{10.1109/TEMC.2011.2139216}

% [14]
\bibitem{Kodera2024}
  S.~Kodera, K.~Taguchi, Y.~Diao, T.~Kashiwa, and A.~Hirata, ``Computation
  of whole-body average SAR in realistic human models from 1 to 100\,GHz,''
  \textit{IEEE Trans.\ Microw.\ Theory Tech.}, vol.~72, no.~1, pp.~91--100,
  2024.
  \doi{10.1109/TMTT.2023.3289562}

% [15]
\bibitem{Diao2024}
  Y.~Diao, S.~Kodera, K.~Li, and A.~Hirata, ``Assessment of whole-body-average
  SAR for exposure to electromagnetic fields up to 30\,GHz using a body model
  with scaled dielectric parameters,'' \textit{IEEE Trans.\ Electromagn.\
  Compat.}, vol.~66, no.~5, pp.~1351--1360, 2024.
  \doi{10.1109/TEMC.2024.3421521}

% [16]
\bibitem{Funahashi2018}
  D.~Funahashi, A.~Hirata, S.~Kodera, and K.~R.~Foster, ``Area-averaged
  transmitted power density at skin surface as metric to estimate surface
  temperature elevation,'' \textit{IEEE Access}, vol.~6, pp.~77665--77674,
  2018.
  \doi{10.1109/ACCESS.2018.2883733}

% [17]
\bibitem{Li2019}
  K.~Li, K.~Sasaki, S.~Watanabe, and H.~Shirai, ``Relationship between power
  density and surface temperature elevation for human skin exposure to
  electromagnetic waves with oblique incidence angle from 6\,GHz to 1\,THz,''
  \textit{Phys.\ Med.\ Biol.}, vol.~64, art.~no.~065016, 2019.
  \doi{10.1088/1361-6560/ab057a}

% [18]
\bibitem{Samaras2019}
  T.~Samaras and N.~Kuster, ``Theoretical evaluation of the power transmitted
  to the body as a function of angle of incidence and polarization at
  frequencies ${>}$6\,GHz and its relevance for standardization,''
  \textit{Bioelectromagnetics}, vol.~40, no.~2, pp.~136--139, 2019.
  \doi{10.1002/bem.22170}

% [19]
\bibitem{Diao2021}
  Y.~Diao, K.~Li, K.~Sasaki, S.~Kodera, I.~Laakso, W.~El~Hajj, and
  A.~Hirata, ``Effect of incidence angle on the spatial-average of incident
  power density definition to correlate skin temperature rise for millimeter
  wave exposures,'' \textit{IEEE Trans.\ Electromagn.\ Compat.}, vol.~63,
  no.~5, pp.~1709--1716, 2021.
  \doi{10.1109/TEMC.2021.3098594}

% [20]
\bibitem{Sasaki2017}
  K.~Sasaki, M.~Mizuno, K.~Wake, and S.~Watanabe, ``Monte Carlo simulations
  of skin exposure to electromagnetic field from 10\,GHz to 1\,THz,''
  \textit{Phys.\ Med.\ Biol.}, vol.~62, no.~17, pp.~6993--7010, 2017.
  \doi{10.1088/1361-6560/aa81fc}

% [21]
\bibitem{Christ2020a}
  A.~Christ, T.~Samaras, E.~Neufeld, and N.~Kuster, ``RF-induced temperature
  increase in a stratified model of the skin for plane-wave exposure at
  6--100\,GHz,'' \textit{Radiat.\ Prot.\ Dosim.}, vol.~188, no.~3,
  pp.~350--360, 2020.
  \doi{10.1093/rpd/ncz293}

% [22]
\bibitem{Potter1970}
  R.~F.~Potter, ``Pseudo-{Brewster} angle and the optical constants,''
  \textit{Appl.\ Opt.}, vol.~9, no.~7, p.~1717, 1970.
  \doi{10.1364/AO.9.001717}

% [23]
\bibitem{Ohman1977}
  G.~P.~\"{O}hman, ``The pseudo-{Brewster} angle,'' \textit{IEEE Trans.\
  Antennas Propag.}, vol.~25, no.~6, pp.~903--904, 1977.
  \doi{10.1109/TAP.1977.1141718}

% [24]
\bibitem{Azzam2015}
  R.~M.~A.~Azzam, ``High-index dielectric substrates with nearly constant
  reflectance for incident unpolarized or circularly polarized light over a
  wide range of incidence angles,'' \textit{J.\ Mod.\ Opt.}, vol.~62, no.~10,
  pp.~811--815, 2015.
  \doi{10.1080/09500340.2015.1007103}

% [25]
\bibitem{BornWolf1999}
  M.~Born and E.~Wolf, \textit{Principles of Optics}, 7th~(expanded)~ed.
  Cambridge: Cambridge Univ.\ Press, 1999.

% [26]
\bibitem{Andersen2007}
  J.~B.~Andersen, J.~{\O}.~Nielsen, G.~F.~Pedersen, G.~Bauch, and
  J.~M.~Herdin, ``Room electromagnetics,'' \textit{IEEE Antennas Propag.\
  Mag.}, vol.~49, no.~2, pp.~27--33, 2007.
  \doi{10.1109/MAP.2007.376642}

% [27]
\bibitem{Bamba2012}
  A.~Bamba, W.~Joseph, J.~B.~Andersen, E.~Tanghe, G.~Vermeeren, D.~Plets,
  J.~{\O}.~Nielsen, and L.~Martens, ``Experimental assessment of specific
  absorption rate using room electromagnetics,'' \textit{IEEE Trans.\
  Electromagn.\ Compat.}, vol.~54, no.~4, pp.~747--757, 2012.
  \doi{10.1109/TEMC.2012.2189572}

% [28]
\bibitem{Bamba2014}
  A.~Bamba, W.~Joseph, G.~Vermeeren, E.~Tanghe, D.~P.~Gaillot,
  J.~B.~Andersen, J.~{\O}.~Nielsen, M.~Li\'{e}nard, and L.~Martens,
  ``Validation of experimental whole-body SAR assessment method in a complex
  indoor environment,'' \textit{Bioelectromagnetics}, vol.~34, no.~2,
  pp.~122--132, 2013.
  \doi{10.1002/bem.21749}

% [29]
\bibitem{Flintoft2014}
  I.~D.~Flintoft, M.~P.~Robinson, G.~C.~R.~Melia, A.~C.~Marvin, and
  J.~F.~Dawson, ``Average absorption cross-section of the human body measured
  at 1--12\,GHz in a reverberant chamber: Results of a human volunteer
  study,'' \textit{Phys.\ Med.\ Biol.}, vol.~59, no.~13, pp.~3297--3317,
  2014.
  \doi{10.1088/0031-9155/59/13/3297}

% [30]
\bibitem{Vouk1948}
  V.~Vouk, ``Projected area of convex bodies,'' \textit{Nature}, vol.~162,
  no.~4113, pp.~330--331, 1948.
  \doi{10.1038/162330a0}

% [31]
\bibitem{Zhang2017thesis}
  X.~Zhang, ``Accurate wideband measurement of human body absorption cross
  section in a reverberation chamber: A morphological parameters study from
  1\,GHz to 18\,GHz,'' Ph.D.\ thesis, Univ.\ York, 2017.

% [32]
\bibitem{ZhangRobinson2020}
  X.~Zhang, M.~P.~Robinson, I.~D.~Flintoft, J.~F.~Dawson, and S.~Parker,
  ``Morphological study on human body absorption cross section in a
  reverberation chamber from 1\,GHz to 16\,GHz,'' \textit{IEEE Trans.\
  Electromagn.\ Compat.}, vol.~62, no.~2, pp.~330--337, 2020.
  \doi{10.1109/TEMC.2019.2904642}

% [33]
\bibitem{Hallbjorner2005}
  P.~Hallbj\"{o}rner, U.~Carlberg, K.~Madsen, and J.~Andersson, ``Extracting
  electrical material parameters of electrically large dielectric objects from
  reverberation chamber measurements of absorption cross section,''
  \textit{IEEE Trans.\ Electromagn.\ Compat.}, vol.~47, no.~2, pp.~291--303,
  2005.
  \doi{10.1109/TEMC.2005.847391}

% [34]
\bibitem{SchneiderWeil2008}
  R.~Schneider and W.~Weil, \textit{Stochastic and Integral Geometry},
  ser.\ Probability and Its Applications. Berlin: Springer-Verlag, 2008.
  \doi{10.1007/978-3-540-78859-1}

% [35]
\bibitem{Landis2002}
  H.~Landis, ``Production-ready global illumination,'' \textit{ACM SIGGRAPH
  2002 Course Notes}, vol.~16, 2002.

% [36]
\bibitem{Pharr2016}
  M.~Pharr, W.~Jakob, and G.~Humphreys, \textit{Physically Based Rendering:
  From Theory to Implementation}, 3rd~ed. San Francisco, CA: Morgan Kaufmann,
  2016.

% [37]
\bibitem{AkenineMoller2018}
  T.~Akenine-M\"{o}ller, E.~Haines, and N.~Hoffman, \textit{Real-Time
  Rendering}, 4th~ed. Boca Raton, FL: A~K Peters/CRC Press, 2018.
  \doi{10.1201/9781315365459}

% [38]
\bibitem{Hochwald2014}
  B.~M.~Hochwald, D.~J.~Love, S.~Yan, P.~Fay, and J.-M.~Jin, ``Incorporating
  specific absorption rate constraints into wireless signal design,''
  \textit{IEEE Commun.\ Mag.}, vol.~52, no.~9, pp.~126--133, Sep.\ 2014.
  \doi{10.1109/MCOM.2014.6894463}

% [39]
\bibitem{Ying2015}
  D.~Ying, D.~J.~Love, and B.~M.~Hochwald, ``Closed-loop precoding and
  capacity analysis for multiple-antenna wireless systems with user radiation
  exposure constraints,'' \textit{IEEE Trans.\ Wireless Commun.}, vol.~14,
  no.~10, pp.~5859--5870, 2015.
  \doi{10.1109/TWC.2015.2443777}

% [40]
\bibitem{Ying2017}
  D.~Ying, D.~J.~Love, and B.~M.~Hochwald, ``Sum-rate analysis for multi-user
  MIMO systems with user exposure constraints,'' \textit{IEEE Trans.\ Wireless
  Commun.}, vol.~16, no.~11, pp.~7376--7388, 2017.
  \doi{10.1109/TWC.2017.2748103}

% [41]
\bibitem{Ebadi2019}
  A.~Ebadi-Shahrivar, P.~Fay, D.~J.~Love, and B.~M.~Hochwald, ``Determining
  electromagnetic exposure compliance of multi-antenna devices in linear
  time,'' \textit{IEEE Trans.\ Antennas Propag.}, vol.~67, no.~12,
  pp.~7585--7596, 2019.
  \doi{10.1109/TAP.2019.2925135}

% [42]
\bibitem{Castellanos2020}
  M.~R.~Castellanos, Y.~Liu, D.~J.~Love, B.~Peleato, J.-M.~Jin, and
  B.~M.~Hochwald, ``Signal-level models of pointwise electromagnetic exposure
  for millimeter wave communication,'' \textit{IEEE Trans.\ Antennas Propag.},
  vol.~68, no.~5, pp.~3963--3977, 2020.
  \doi{10.1109/TAP.2019.2946742}

% [43] *** FLAGGED: Zhou2026 was cited in text but had no .bib entry.
%          Entry reconstructed from in-text information. Verify before submission.
\bibitem{Zhou2026}
  C.~Zhou, Y.~Chen, S.~Chen, and B.~Wang, ``Exposure-aware beamforming for
  mmWave systems: From EM theory to thermal compliance,''
  \textit{arXiv:2601.19587}, 2026.

% [44]
\bibitem{Shikhantsov2020}
  S.~Shikhantsov, A.~Thielens, G.~Vermeeren, P.~Demeester, L.~Martens,
  G.~Torfs, and W.~Joseph, ``Massive MIMO propagation modeling with
  user-induced coupling effects using ray-tracing and FDTD,'' \textit{IEEE J.\
  Sel.\ Areas Commun.}, vol.~38, no.~9, pp.~1955--1963, 2020.
  \doi{10.1109/JSAC.2020.3000874}

% [45]
\bibitem{Wydaeghe2026npj}
  R.~Wydaeghe, S.~Shikhantsov, G.~Vermeeren, L.~Martens, E.~Tanghe, and
  W.~Joseph, ``Hybrid ray-tracing-{QuaDRiGa}/{FDTD} method for realistic
  28\,GHz exposure with 6G CF-MaMIMO in 3D outdoor environments,''
  \textit{npj Wireless Technol.}, 2026.

% [46]
\bibitem{WydaegheEtAl2026PMB}
  R.~Wydaeghe \textit{et~al.}, ``Systematic FDTD campaign for whole-body
  dosimetry at 450\,MHz to 26\,GHz,'' \emph{in preparation}, 2026.

% [47]
\bibitem{Kushiyama2025}
  Y.~Kushiyama and T.~Nagaoka, ``Whole-body exposure assessment from
  millimeter-wave base stations using spherical-wave lookup tables,''
  \textit{IEEE Trans.\ Electromagn.\ Compat.}, 2025.

% [48]
\bibitem{DuBois1916}
  D.~Du~Bois and E.~F.~Du~Bois, ``Clinical calorimetry: Tenth paper.
  A formula to estimate the approximate surface area if height and weight be
  known,'' \textit{Arch.\ Intern.\ Med.}, vol.~17, pp.~863--871, 1916.
  \doi{10.1001/archinte.1916.00080130010002}

% [49]
\bibitem{Gabriel1996}
  S.~Gabriel, R.~W.~Lau, and C.~Gabriel, ``The dielectric properties of
  biological tissues: III.\ Parametric models for the dielectric spectrum of
  tissues,'' \textit{Phys.\ Med.\ Biol.}, vol.~41, no.~11, pp.~2271--2293,
  1996.
  \doi{10.1088/0031-9155/41/11/003}

% [50]
\bibitem{ITISv5}
  {IT'IS Foundation}, ``Tissue properties database v5.0,'' 2024. [Online].
  Available: \url{https://itis.swiss/virtual-population/tissue-properties/database/}

% [51]
\bibitem{AlekseevZiskin2007}
  S.~I.~Alekseev and M.~C.~Ziskin, ``Human skin permittivity determined by
  millimeter wave reflection measurements,'' \textit{Bioelectromagnetics},
  vol.~28, no.~5, pp.~331--339, 2007.
  \doi{10.1002/bem.20308}

% [52]
\bibitem{BohrenHuffman1983}
  C.~F.~Bohren and D.~R.~Huffman, \textit{Absorption and Scattering of Light
  by Small Particles}. New York: Wiley, 1983.

% [53]
\bibitem{Goodfellow2016}
  I.~Goodfellow, Y.~Bengio, and A.~Courville, \textit{Deep Learning}.
  Cambridge, MA: MIT Press, 2016.

% [54]
\bibitem{Vermeeren2008}
  G.~Vermeeren, W.~Joseph, C.~Olivier, and L.~Martens, ``Statistical multipath
  exposure of a human in a realistic electromagnetic environment,''
  \textit{Health Phys.}, vol.~94, no.~4, pp.~345--354, 2008.
  \doi{10.1097/01.HP.0000298816.66888.05}

% [55] SionnaEM is the same paper as SionnaRT [8]. Kept as a separate key
%      because it is cited independently in Part III (Sec. field-channel).
\bibitem{SionnaEM}
  J.~Hoydis, F.~A\"{i}t~Aoudia, S.~Cammerer, M.~Nimier-David, N.~Binder,
  G.~Marcus, and A.~Keller, ``{Sionna RT}: Differentiable ray tracing for
  radio propagation modeling,'' \textit{arXiv:2303.11103}, 2023.
  \doi{10.48550/arXiv.2303.11103}

\end{thebibliography}
\endgroup

\end{document}
